% Généralités sur les fonctions — Mathématiques 1ère C — Cours signé OnBuch+
% Programme officiel MINESEC. Compiler avec : tectonic N04-generalites-fonctions.tex
\newcommand{\DOCMATIERE}{Mathématiques}
\newcommand{\DOCNIVEAU}{1ère C}
\newcommand{\DOCMODULE}{Relations et opérations fondamentales dans l'ensemble des nombres réels}
\newcommand{\DOCLECON}{N04}
\newcommand{\DOCTITRE}{Généralités sur les fonctions}
% OnBuch+ — préambule « cours signés » ENRICHI (compilé avec tectonic / XeLaTeX)
% Système visuel « Pop » d'OnBuch+ : fond crème (jamais blanc), bordures encre
% épaisses, ombres « dures » décalées, titres Archivo Black en capitales.
% Version MATHÉMATIQUES : graphes (pgfplots/tikz), boîtes pédagogiques
% (définition, propriété, méthode, attention, le savais-tu, exercice résolu,
% exercice, TP, bilan), page de garde et signature OnBuch+ ancrée en bas.
%
% Macros attendues dans main.tex AVANT \input{preamble} :
%   \DOCMATIERE  \DOCNIVEAU  \DOCMODULE  \DOCLECON (numéro)  \DOCTITRE
\documentclass[11pt]{article}

\usepackage{fontspec}
\usepackage[french]{babel}
\usepackage[a4paper,margin=2cm,top=3.4cm,bottom=2.8cm,headheight=1.4cm,footskip=1.3cm]{geometry}
\usepackage{amsmath,amssymb,amsfonts}
\usepackage{mathtools} % \xmapsto, \coloneqq... souvent utilisés par les modèles
\usepackage{cancel} % simplifications barrées
% \abs{z} (module, valeur absolue) et \norm{u} (norme), utilisés par les modèles
\providecommand{\abs}{}\let\abs\relax\DeclarePairedDelimiter{\abs}{\lvert}{\rvert}
\providecommand{\norm}{}\let\norm\relax\DeclarePairedDelimiter{\norm}{\lVert}{\rVert}
\usepackage[table]{xcolor} % [table] : fournit aussi \rowcolors (alternance de lignes)
\usepackage{pagecolor}
\usepackage{tikz}
\usetikzlibrary{arrows.meta,positioning,calc,shapes.geometric,shapes.misc,shapes.arrows,decorations.pathmorphing,decorations.pathreplacing,decorations.markings,patterns,shadows,mindmap,trees,fit,backgrounds,matrix,intersections,angles,quotes}
\usepackage{pgfplots}
\usepackage{pgf-pie} % diagrammes circulaires \pie (statistiques)
\pgfplotsset{compat=1.18}
\usepgfplotslibrary{fillbetween,groupplots} % aires entre courbes (calcul intégral)
\usepackage{enumitem}
\usepackage{fancyhdr}
% [most] charge toutes les bibliothèques tcolorbox (listings, minted, raster,
% documentation...) et gonfle inutilement la mémoire TeX sur un document long
% (~300 boîtes) : on ne charge que ce qui est réellement utilisé.
\usepackage[skins,breakable]{tcolorbox}
\usepackage{graphicx}
\usepackage{colortbl}
\usepackage{booktabs}
\usepackage{tabularx}
\usepackage{diagbox} % case d'en-tête diagonale des tableaux à double entrée
% Colonnes extensibles centrée / gauche / droite (C, L, R), souvent utilisées par les modèles
\newcolumntype{C}{>{\centering\arraybackslash}X}
\newcolumntype{L}{>{\raggedright\arraybackslash}X}
\newcolumntype{R}{>{\raggedleft\arraybackslash}X}
\usepackage{array}
\usepackage{multicol}
\usepackage{siunitx}
% parse-numbers=false : accepte aussi \SI{2,5 \times 10^{-3}}{...}
\sisetup{locale=FR,output-decimal-marker={,},group-minimum-digits=4,parse-numbers=false}
\DeclareSIUnit\ohm{\ensuremath{\symup{\Omega}}} % U+2126 absent de la police
\usepackage{qrcode}
% \degree : commande fréquemment utilisée par les modèles pour un angle ou une
% température en dehors de \SI{}{} (non fournie par les paquets chargés).
\providecommand{\degree}{^{\circ}}
% \qed (fin de démonstration), utilisé par les modèles sans amsthm
\providecommand{\qed}{\hfill$\square$}
% \barre : double barre des valeurs interdites dans un tableau de variations (tkz-tab)
\providecommand{\barre}{\ensuremath{\|}}
% environnement proof (amsthm non chargé) utilisé par les modèles
\ifdefined\proof\else\newenvironment{proof}[1][Démonstration]{\par\noindent\textit{#1.}\ }{\qed\par}\fi
\usepackage{lastpage}
\usepackage{hyperref}
\hypersetup{hidelinks}

% ============================================================
% Palette « Pop » OnBuch+ — mêmes hex que la classe P (plus_ui.dart)
% ============================================================
\definecolor{poporange}{HTML}{F59321}
\definecolor{popdark}{HTML}{E07A0C}
\definecolor{popcream}{HTML}{FFF6E8}
\definecolor{popink}{HTML}{1C1714}
\definecolor{poppurple}{HTML}{7A5AE0}
\definecolor{popgreen}{HTML}{2E9E6A}
\definecolor{poppink}{HTML}{E0457B}
\definecolor{popblue}{HTML}{2F7DE1}
\definecolor{popgold}{HTML}{C98A12}
\definecolor{popmuted}{HTML}{8A7F76}
\definecolor{popline}{HTML}{EADFD2}
\definecolor{popgray}{HTML}{8A7F76}
\colorlet{popgrey}{popgray}
% Alias tolérants (le modèle nomme parfois les couleurs différemment)
\colorlet{popwhite}{white}
\colorlet{popblack}{popink}
\colorlet{popred}{poppink}
\colorlet{popyellow}{popgold}
% Teintes claires pour fonds de tableaux / zones
\colorlet{popblueL}{popblue!10!white}
\colorlet{poporangeL}{poporange!14!white}
\colorlet{popgreenL}{popgreen!12!white}
\colorlet{poppurpleL}{poppurple!12!white}
\colorlet{poppinkL}{poppink!12!white}
\colorlet{popgoldL}{popgold!14!white}

\pagecolor{popcream}

% ============================================================
% Typographie
% ============================================================
\defaultfontfeatures{Ligatures=TeX}
\setmainfont{PlusJakartaSans-Medium.ttf}[Path=../../../fonts/,BoldFont=PlusJakartaSans-Bold.ttf]
% Maths en sans-serif (Fira Math) avec chiffres et lettres droites de Plus
% Jakarta, pour que les formules se fondent dans le texte.
\usepackage[math-style=ISO,bold-style=ISO]{unicode-math}
\setmathfont{FiraMath-Regular.otf}[Path=../../../fonts/]
\setmathfont{PlusJakartaSans-Medium.ttf}[Path=../../../fonts/,range={up/num,up/{Latin,latin},\mathup}]
\usepackage{icomma} % virgule décimale sans espace en mode math
% Caractères mathématiques Unicode absents des polices de texte (Plus Jakarta,
% Archivo Black) : rendus par leur équivalent mathématique, en texte comme en
% formule — sinon ils s'affichent en carré vide (ex. « Arithmétique dans ℤ »).
\usepackage{newunicodechar}
\newunicodechar{ℝ}{\ensuremath{\mathbb{R}}}
\newunicodechar{ℤ}{\ensuremath{\mathbb{Z}}}
\newunicodechar{ℂ}{\ensuremath{\mathbb{C}}}
\newunicodechar{ℕ}{\ensuremath{\mathbb{N}}}
\newunicodechar{ℚ}{\ensuremath{\mathbb{Q}}}
\newunicodechar{𝔻}{\ensuremath{\mathbb{D}}}
\newunicodechar{α}{\ensuremath{\alpha}}
\newunicodechar{β}{\ensuremath{\beta}}
\newunicodechar{γ}{\ensuremath{\gamma}}
\newunicodechar{δ}{\ensuremath{\delta}}
\newunicodechar{ε}{\ensuremath{\varepsilon}}
\newunicodechar{θ}{\ensuremath{\theta}}
\newunicodechar{λ}{\ensuremath{\lambda}}
\newunicodechar{μ}{\ensuremath{\mu}}
\newunicodechar{σ}{\ensuremath{\sigma}}
\newunicodechar{τ}{\ensuremath{\tau}}
\newunicodechar{φ}{\ensuremath{\varphi}}
\newunicodechar{ψ}{\ensuremath{\psi}}
\newunicodechar{ω}{\ensuremath{\omega}}
\newunicodechar{Σ}{\ensuremath{\Sigma}}
% majuscules produites par \MakeUppercase dans les titres (α-aminés -> Α)
\newunicodechar{Α}{\ensuremath{\alpha}}
\newunicodechar{Β}{\ensuremath{\beta}}
\newunicodechar{Γ}{\ensuremath{\Gamma}}
\newunicodechar{Δ}{\ensuremath{\Delta}}
\newunicodechar{Ε}{\ensuremath{\varepsilon}}
\newunicodechar{Θ}{\ensuremath{\Theta}}
\newunicodechar{Λ}{\ensuremath{\Lambda}}
\newunicodechar{Μ}{\ensuremath{\mu}}
\newunicodechar{Τ}{\ensuremath{\tau}}
\newunicodechar{Φ}{\ensuremath{\Phi}}
\newunicodechar{Ψ}{\ensuremath{\Psi}}
\newunicodechar{Ω}{\ensuremath{\Omega}}
\newunicodechar{↦}{\ensuremath{\mapsto}}
\newunicodechar{∈}{\ensuremath{\in}}
\newunicodechar{∉}{\ensuremath{\notin}}
\newunicodechar{∧}{\ensuremath{\wedge}}
\newunicodechar{⇔}{\ensuremath{\Leftrightarrow}}
\newunicodechar{⟺}{\ensuremath{\Longleftrightarrow}}
\newunicodechar{⇒}{\ensuremath{\Rightarrow}}
\newunicodechar{⟹}{\ensuremath{\Longrightarrow}}
\newunicodechar{∘}{\ensuremath{\circ}}
\newunicodechar{∪}{\ensuremath{\cup}}
\newunicodechar{∩}{\ensuremath{\cap}}
\newunicodechar{≡}{\ensuremath{\equiv}}
\newunicodechar{⊂}{\ensuremath{\subset}}
\newunicodechar{⊥}{\ensuremath{\perp}}
\newunicodechar{∃}{\ensuremath{\exists}}
\newunicodechar{∀}{\ensuremath{\forall}}
\newunicodechar{∛}{\ensuremath{\sqrt[3]{\,}}}
\newunicodechar{∜}{\ensuremath{\sqrt[4]{\,}}}
\newunicodechar{⃗}{\ensuremath{{}^{\to}}}
\newunicodechar{ⁿ}{\ifmmode^{n}\else\textsuperscript{\ensuremath{n}}\fi}
\newunicodechar{ˣ}{\ifmmode^{x}\else\textsuperscript{\ensuremath{x}}\fi}
\newunicodechar{ᵇ}{\ifmmode^{b}\else\textsuperscript{\ensuremath{b}}\fi}
\newunicodechar{ʳ}{\ifmmode^{r}\else\textsuperscript{\ensuremath{r}}\fi}
\newunicodechar{ᵘ}{\ifmmode^{u}\else\textsuperscript{\ensuremath{u}}\fi}
\newunicodechar{ʸ}{\ifmmode^{y}\else\textsuperscript{\ensuremath{y}}\fi}
\newunicodechar{ᵃ}{\ifmmode^{a}\else\textsuperscript{\ensuremath{a}}\fi}
\newunicodechar{ᵉ}{\ifmmode^{e}\else\textsuperscript{\ensuremath{e}}\fi}
\newunicodechar{ᵏ}{\ifmmode^{k}\else\textsuperscript{\ensuremath{k}}\fi}
\newunicodechar{ᵖ}{\ifmmode^{p}\else\textsuperscript{\ensuremath{p}}\fi}
\newunicodechar{ᴺ}{\ifmmode^{N}\else\textsuperscript{\ensuremath{N}}\fi}
\newunicodechar{ᶜ}{\ifmmode^{c}\else\textsuperscript{\ensuremath{c}}\fi}
\newunicodechar{ʰ}{\ifmmode^{h}\else\textsuperscript{\ensuremath{h}}\fi}
\newunicodechar{ᵠ}{\ifmmode^{q}\else\textsuperscript{\ensuremath{q}}\fi}
\newunicodechar{ₙ}{\ifmmode_{n}\else\textsubscript{\ensuremath{n}}\fi}
\newunicodechar{ₐ}{\ifmmode_{a}\else\textsubscript{\ensuremath{a}}\fi}
\newunicodechar{ᵢ}{\ifmmode_{i}\else\textsubscript{\ensuremath{i}}\fi}
\newunicodechar{ᵣ}{\ifmmode_{r}\else\textsubscript{\ensuremath{r}}\fi}
\newunicodechar{ₖ}{\ifmmode_{k}\else\textsubscript{\ensuremath{k}}\fi}
\newunicodechar{ᵦ}{\ifmmode_{\beta}\else\textsubscript{\ensuremath{\beta}}\fi}
\newunicodechar{ₓ}{\ifmmode_{x}\else\textsubscript{\ensuremath{x}}\fi}
\newfontfamily\popdisplay{ArchivoBlack-Regular.ttf}[Path=../../../fonts/]
\newfontfamily\poplogo{SpaceGrotesk-Bold.ttf}[Path=../../../fonts/]
\newfontfamily\popsemi{PlusJakartaSans-SemiBold.ttf}[Path=../../../fonts/]
\newfontfamily\popxbold{PlusJakartaSans-ExtraBold.ttf}[Path=../../../fonts/]
\newfontfamily\popxboldls{PlusJakartaSans-ExtraBold.ttf}[Path=../../../fonts/,LetterSpace=3.0]

\setlist[enumerate]{leftmargin=1.7em,itemsep=5pt,topsep=4pt}
\setlist[itemize]{leftmargin=1.7em,itemsep=4pt,topsep=4pt,label={\color{poporange}\textbullet}}
\setlength{\parindent}{0pt}
\setlength{\parskip}{5pt}
\renewcommand{\arraystretch}{1.25}

% pgfplots : style Pop par défaut
\pgfplotsset{
  every axis/.append style={axis line style={popink,line width=1pt},tick style={popink},
    grid style={popline,line width=0.5pt},label style={font=\small},tick label style={font=\footnotesize}},
  popcurve/.style={poporange,line width=1.8pt,no markers,smooth},
}
% Même style disponible dans un \draw TikZ simple (hors pgfplots)
\tikzset{popcurve/.style={poporange,line width=1.8pt}}

% ============================================================
% Pill / squiggle
% ============================================================
\newcommand{\poppill}[3]{%
  \tikz[baseline=-0.55ex]\node[draw=popink,line width=1.2pt,rounded corners=2.6pt,%
    fill=#2,inner xsep=6.5pt,inner ysep=3pt]{%
    \popxboldls\fontsize{7}{7}\selectfont\color{#3}\MakeUppercase{#1}};%
}
\newcommand{\popsquiggle}[1]{%
  \tikz[baseline=-0.4ex]{
    \draw[#1,line width=2.6pt,line cap=round]
      plot [smooth] coordinates {(0,0.09) (0.34,0.01) (0.68,0.21) (1.02,0.01) (1.36,0.10)};
  }%
}
% Mot-clé surligné (marqueur orange sous le texte)
\newcommand{\cle}[1]{{\popxbold\color{popdark}#1}}

% ============================================================
% En-tête / pied
% ============================================================
\pagestyle{fancy}
\fancyhf{}
\renewcommand{\headrulewidth}{0pt}
\renewcommand{\footrulewidth}{0pt}
\lhead{\raisebox{-0.15ex}{\poplogo\fontsize{13}{13}\selectfont\color{popink}OnBuch\color{poporange}+}}
\rhead{\poppill{\DOCMATIERE\ \textbullet\ \DOCNIVEAU\ \textbullet\ Leçon \DOCLECON}{white}{popink}}
\lfoot{\popxbold\fontsize{7.6}{7.6}\selectfont\color{popmuted}\MakeUppercase{Cours signé OnBuch+}\ \textbullet\ {\popsemi\DOCTITRE}}
\rfoot{\tikz[baseline=-0.6ex]\node[rounded corners=8.5pt,fill=popink,minimum height=17pt,minimum width=17pt,inner xsep=5pt,inner sep=0]{\popxbold\fontsize{8}{8}\selectfont\color{popcream}\thepage\,/\,\pageref*{LastPage}};}

% ============================================================
% Titres
% ============================================================
\newcommand{\chaptitre}[2]{%
  \begin{tcolorbox}[enhanced,colback=poporange,colframe=popink,boxrule=2.6pt,arc=11pt,
      drop shadow=popink,left=15pt,right=15pt,top=12pt,bottom=11pt,before skip=3pt,after skip=18pt]
    \poppill{#2}{popink}{popcream}\par\vspace{9pt}
    {\raggedright\hyphenpenalty=10000\exhyphenpenalty=10000\popdisplay\fontsize{22}{24}\selectfont\color{popcream}\MakeUppercase{#1}\par}
  \end{tcolorbox}
}
% Section numérotée : pastille orange avec le numéro + titre Archivo
\newcounter{popsec}
\newcommand{\coursec}[1]{%
  \stepcounter{popsec}\setcounter{popsub}{0}%
  \par\vspace{16pt}\noindent
  \begin{minipage}{\linewidth}
  \tikz[baseline=-0.7ex]\node[draw=popink,line width=1.6pt,fill=poporange,rounded corners=4pt,minimum size=22pt,inner sep=0]{\popdisplay\fontsize{12}{12}\selectfont\color{popcream}\thepopsec};%
  \hspace{8pt}{\popdisplay\fontsize{15}{17}\selectfont\color{popink}\MakeUppercase{#1}}\par
  \vspace{4pt}\noindent{\color{poporange}\rule{36pt}{3.6pt}}
  \end{minipage}\par\nopagebreak\vspace{8pt}
}
\newcounter{popsub}
\newcommand{\courssub}[1]{%
  \stepcounter{popsub}%
  \par\vspace{9pt}\noindent{\popxbold\fontsize{11.5}{13}\selectfont\color{popink}{\color{poporange}\thepopsec.\thepopsub}\hspace{6pt}#1}\par\nopagebreak\vspace{4pt}
}

% ============================================================
% Boîtes pédagogiques — même recette : fond blanc, bordure épaisse colorée,
% ombre dure encre, badge plein. Titre optionnel [#1].
% ============================================================
\tcbset{popbase/.style={breakable,enhanced,colback=white,boxrule=2.3pt,arc=9pt,
  shadow={3pt}{-3pt}{0pt}{popink},left=14pt,right=14pt,top=11pt,bottom=11pt,before skip=11pt,after skip=15pt,
  fontupper=\popsemi\fontsize{10.6}{14.5}\selectfont\color{popink},
  fontlower=\popsemi\fontsize{10.6}{14.5}\selectfont\color{popink}}}
% Coche dessinée (le glyphe ✓ n'existe pas dans les polices du document)
\newcommand{\popcheck}[1]{\tikz[baseline=-0.5ex]\draw[#1,line width=1.6pt,line cap=round,line join=round](-0.13,0.0)--(-0.04,-0.09)--(0.13,0.1);}
\newcommand{\popboxhead}[3]{\poppill{#1}{#2}{white}\ifx\relax#3\relax\else\hspace{6pt}{\popxbold\fontsize{10.6}{12}\selectfont\color{popink}#3}\fi\par\vspace{7pt}}

\newtcolorbox{aretenir}[1][]{popbase,colframe=popgreen,before upper={\popboxhead{À retenir}{popgreen}{#1}}}
\newtcolorbox{exemplebox}[1][]{popbase,colframe=poporange,before upper={\popboxhead{Exemple}{poporange}{#1}}}
\newtcolorbox{definition}[1][]{popbase,colframe=popblue,before upper={\popboxhead{Définition}{popblue}{#1}}}
\newtcolorbox{propriete}[1][]{popbase,colframe=poppurple,before upper={\popboxhead{Propriété}{poppurple}{#1}}}
\newtcolorbox{methode}[1][]{popbase,colframe=popgold,before upper={\popboxhead{Méthode}{popgold}{#1}}}
\newtcolorbox{attention}[1][]{popbase,colframe=poppink,before upper={\popboxhead{Attention}{poppink}{#1}}}
\newtcolorbox{experience}[1][]{popbase,colframe=popblue,colback=popblueL,before upper={\popboxhead{Expérience}{popblue}{#1}}}
% « Le savais-tu ? » : carte violette pleine (contexte, culture, Cameroun)
\newtcolorbox{savaistu}[1][]{popbase,colback=poppurple,colframe=popink,
  fontupper=\popsemi\fontsize{10.4}{14.2}\selectfont\color{white},
  before upper={\poppill{Le savais-tu\,?}{white}{poppurple}\ifx\relax#1\relax\else\hspace{6pt}{\popxbold\color{white}#1}\fi\par\vspace{7pt}}}
% Exercice résolu : énoncé en haut, \tcblower puis la solution sur fond crème
\newcounter{popexo}\newcounter{popexr}
\newtcolorbox{exoresolu}[1][]{popbase,colframe=poporange,colbacklower=popcream,
  segmentation style={popink,line width=1.2pt,dashed},
  before upper={\stepcounter{popexr}\popboxhead{Exercice résolu \thepopexr}{poporange}{#1}},
  before lower={\poppill{Solution}{popink}{popcream}\par\vspace{6pt}}}
% Exercice (non résolu en place) — la correction est dans la partie Corrigés
\newtcolorbox{exercice}[1][]{popbase,colframe=popink,
  before upper={\stepcounter{popexo}\popboxhead{Exercice \thepopexo}{popink}{#1}}}
% Corrigé
\newtcolorbox{corrige}[1][]{popbase,colframe=popgreen,colback=popgreenL,
  before upper={\popboxhead{Corrigé}{popgreen}{#1}}}
% Objectifs / prérequis / bilan
\newtcolorbox{objectifs}{popbase,colframe=popink,colback=poporangeL,
  before upper={\popboxhead{Objectifs de la leçon}{popink}{}}}
\newtcolorbox{prerequis}{popbase,colframe=popink,colback=popblueL,
  before upper={\popboxhead{Prérequis}{popblue}{}}}
\newtcolorbox{bilan}{popbase,colframe=popink,colback=white,boxrule=2.8pt,
  before upper={\popboxhead{Fiche bilan}{poporange}{L'essentiel à connaître pour l'examen}}}
% Figure encadrée avec légende
\newtcolorbox{popfigure}[1][]{enhanced,breakable=false,colback=white,colframe=popline,boxrule=1.4pt,arc=8pt,
  left=10pt,right=10pt,top=10pt,bottom=8pt,before skip=10pt,after skip=12pt,halign=center,
  lower separated=false,halign lower=center,
  fontlower=\popsemi\fontsize{9}{11.5}\selectfont\color{popmuted},
  #1}
\newcommand{\legende}[1]{\par\vspace{6pt}{\popsemi\fontsize{9}{11.5}\selectfont\color{popmuted}#1}}

% ============================================================
% Signature OnBuch+
% ============================================================
\newcommand{\poplogoinline}[1]{{\poplogo\fontsize{#1}{#1}\selectfont\color{popink}OnBuch\color{poporange}+}}
\newcommand{\signatureonbuch}{%
  \begin{tcolorbox}[enhanced,colback=popink,colframe=popink,boxrule=2.2pt,arc=10pt,
      drop shadow=poporange,left=14pt,right=14pt,top=10pt,bottom=10pt,before skip=0pt,after skip=0pt]
    \tikz[baseline=-0.7ex]\node[circle,fill=popgreen,draw=popcream,line width=1.3pt,minimum size=22pt,inner sep=0]{\popcheck{white}};%
    \hspace{9pt}\begin{minipage}[c]{0.62\linewidth}
      {\popxbold\fontsize{12}{13}\selectfont\color{popcream}Signé\ \textbullet\ {\poplogo OnBuch\color{poporange}+}}\par\vspace{2pt}
      {\raggedright\popsemi\fontsize{8}{10.5}\selectfont\color{popline}\MakeUppercase{Équipe pédagogique OnBuch+}\\\MakeUppercase{Conforme au programme officiel MINESEC}\par}
    \end{minipage}\hfill
    {\poplogo\fontsize{22}{22}\selectfont\color{popcream}OnBuch\color{poporange}+}
  \end{tcolorbox}%
}

% ============================================================
% Page de garde
% #1 titre · #2 matière · #3 niveau · #4 module · #5 n° leçon · #6 durée ·
% #7 description / objectifs (texte ou itemize)
% ============================================================
\newcommand{\pagegarde}[7]{%
  \thispagestyle{empty}
  \newgeometry{a4paper,margin=1.7cm,top=1.6cm,bottom=1.4cm}%
  \begin{tikzpicture}[remember picture,overlay]
    \fill[poporange!18] ([xshift=-3.2cm,yshift=-3.6cm]current page.north east) circle (4.6cm);
    \fill[poppurple!14] ([xshift=2.4cm,yshift=7.6cm]current page.south west) circle (3.8cm);
  \end{tikzpicture}%
  {\poplogo\fontsize{30}{30}\selectfont\color{popink}OnBuch\color{poporange}+}\hfill
  \raisebox{0.6ex}{\poppill{Cours signé \textbullet\ Premium}{popink}{popcream}}\par
  \vspace{1pt}\popsquiggle{poppurple}\par
  \vspace{8pt}
  \begin{tcolorbox}[enhanced,colback=poporange,colframe=popink,boxrule=2.8pt,arc=13pt,
      drop shadow=popink,left=18pt,right=18pt,top=12pt,bottom=13pt,after skip=10pt]
    \poppill{#2 \textbullet\ #3}{popink}{popcream}\hspace{5pt}\poppill{#4}{white}{popink}\par\vspace{8pt}
    {\popdisplay\fontsize{14}{14}\selectfont\color{popink}LEÇON #5}\par\vspace{4pt}
    {\raggedright\hyphenpenalty=10000\exhyphenpenalty=10000\popdisplay\fontsize{25}{27}\selectfont\color{popcream}\MakeUppercase{#1}\par}
    \vspace{8pt}
    {\popxbold\fontsize{9.5}{9.5}\selectfont\color{popink}\MakeUppercase{Durée conseillée : #6}}
  \end{tcolorbox}
  \begin{tcolorbox}[enhanced,colback=white,colframe=popink,boxrule=2.2pt,arc=10pt,
      drop shadow=popink,left=16pt,right=16pt,top=10pt,bottom=10pt,after skip=8pt]
    \poppill{Dans cette leçon}{poppurple}{white}\par\vspace{5pt}
    \popsemi\fontsize{9.3}{12.6}\selectfont\color{popink}#7
  \end{tcolorbox}
  \noindent\begin{minipage}[t]{0.487\linewidth}
    \begin{tcolorbox}[enhanced,colback=popgreen,colframe=popink,boxrule=2.2pt,arc=10pt,
        drop shadow=popink,left=11pt,right=11pt,top=10pt,bottom=10pt,height=3.1cm,valign=center]
      \tcbox[colback=white,colframe=popink,boxrule=1.3pt,arc=5pt,boxsep=0pt,nobeforeafter,
        left=3pt,right=3pt,top=3pt,bottom=3pt,valign=center]{\qrcode[height=1.9cm]{https://play.google.com/store/apps/details?id=com.luvvix.onbuch&hl=fr}}%
      \hspace{8pt}\begin{minipage}[c]{0.5\linewidth}
        {\popdisplay\fontsize{11}{12.5}\selectfont\color{white}\MakeUppercase{Télécharge OnBuch}}\par\vspace{4pt}
        {\raggedright\popsemi\fontsize{8.4}{11}\selectfont\color{white}Gratuit \textbullet\ Google Play\\Tuteur IA, annales, concours, résultats\par}
      \end{minipage}
    \end{tcolorbox}
  \end{minipage}\hfill
  \begin{minipage}[t]{0.487\linewidth}
    \begin{tcolorbox}[enhanced,colback=poppurple,colframe=popink,boxrule=2.2pt,arc=10pt,
        drop shadow=popink,left=11pt,right=11pt,top=10pt,bottom=10pt,height=3.1cm,valign=center]
      \tikz[baseline=-0.7ex]\node[circle,fill=white,draw=popink,line width=1.3pt,minimum size=1.6cm,inner sep=0]{\popdisplay\fontsize{24}{24}\selectfont\color{poppurple}+};%
      \hspace{8pt}\begin{minipage}[c]{0.6\linewidth}
        {\popdisplay\fontsize{11}{12.5}\selectfont\color{white}\MakeUppercase{Passe à OnBuch+}}\par\vspace{4pt}
        {\raggedright\popsemi\fontsize{8.4}{11}\selectfont\color{white}Tous les cours signés, Tuteur IA illimité, fiches PDF — onglet OnBuch+ de l'app\par}
      \end{minipage}
    \end{tcolorbox}
  \end{minipage}
  \vspace{8pt}
  \popcontenu
  \vfill
  \signatureonbuch
  \restoregeometry
  \newpage
}

% Rangée « Ce cours contient » (page de garde)
\newcommand{\poptile}[3]{% #1 couleur #2 gros texte #3 légende
  \begin{tcolorbox}[enhanced,colback=#1,colframe=popink,boxrule=2pt,arc=9pt,shadow={2.5pt}{-2.5pt}{0pt}{popink},
      left=6pt,right=6pt,top=9pt,bottom=9pt,halign=center,valign=center,height=1.75cm,width=\linewidth,nobeforeafter]
    {\popdisplay\fontsize{12.5}{13}\selectfont\color{white}\MakeUppercase{#2}}\par\vspace{3pt}
    {\centering\popsemi\fontsize{7.8}{9.5}\selectfont\color{white}#3\par}
  \end{tcolorbox}}
\newcommand{\popcontenu}{%
  \par\noindent{\popxboldls\fontsize{7.5}{7.5}\selectfont\color{popmuted}CE COURS SIGNÉ CONTIENT}\par\vspace{7pt}
  \noindent\begin{minipage}[t]{0.235\linewidth}\poptile{popblue}{Cours}{complet, illustré, pas à pas}\end{minipage}\hfill
  \begin{minipage}[t]{0.235\linewidth}\poptile{poporange}{Méthodes}{et exercices résolus}\end{minipage}\hfill
  \begin{minipage}[t]{0.235\linewidth}\poptile{poppink}{Exercices}{type examen + corrigés}\end{minipage}\hfill
  \begin{minipage}[t]{0.235\linewidth}\poptile{popgreen}{Fiche bilan}{l'essentiel à retenir}\end{minipage}\par}

% Sommaire
\newtcolorbox{sommaire}{enhanced,colback=white,colframe=popink,boxrule=2.2pt,arc=10pt,
  drop shadow=popink,left=18pt,right=18pt,top=13pt,bottom=13pt,before skip=6pt,after skip=20pt,
  before upper={\poppill{Sommaire}{poppurple}{white}\par\vspace{9pt}\popsemi\fontsize{10.6}{14.5}\selectfont\color{popink}}}

% Signature de fin de cours — poussée tout en bas de la dernière page
\newcommand{\findecours}{%
  \par\vfill\nopagebreak
  \begin{center}\popsquiggle{poporange}\end{center}\vspace{4pt}
  \signatureonbuch
}
\AtBeginDocument{\def\llbracket{\mathopen{[\mkern-3.5mu[}}\def\rrbracket{\mathclose{]\mkern-3.5mu]}}} % intervalles d'entiers (glyphe ⟦ absent de la police)
% Glyphes absents de Fira Math : redéfinis à partir de glyphes disponibles
\AtBeginDocument{%
  \def\setminus{\mathbin{\backslash}}\let\smallsetminus\setminus
  \def\vdots{\mathord{\vbox{\baselineskip3.2pt\lineskiplimit0pt\kern4pt\hbox{.}\hbox{.}\hbox{.}}}}%
  \def\ddots{\mathinner{\mkern1mu\raise6.5pt\hbox{.}\mkern2mu\raise3.6pt\hbox{.}\mkern2mu\raise0.7pt\hbox{.}\mkern1mu}}%
  \def\triangle{\mathord{\scalebox{0.9}{$\symup{\Delta}$}}}%
  \def\nabla{\mathord{\raisebox{\depth}{\scalebox{1}[-1]{$\symup{\Delta}$}}}}%
  \def\checkmark{\popcheck{popdark}}%
  \def\star{\mathbin{\ast}}\def\bigstar{\mathord{\ast}}%
  \def\diamond{\mathbin{\scalebox{0.6}{$\Diamond$}}}%
  \def\bigcup{\mathop{\mathchoice{\raisebox{-0.3em}{\scalebox{1.7}{$\cup$}}}{\scalebox{1.25}{$\cup$}}{\cup}{\cup}}\displaylimits}%
  \def\bigcap{\mathop{\mathchoice{\raisebox{-0.3em}{\scalebox{1.7}{$\cap$}}}{\scalebox{1.25}{$\cap$}}{\cap}{\cap}}\displaylimits}%
  \def\lesseqqgtr{\mathrel{\lessgtr}}\def\bigoplus{\mathop{\oplus}\displaylimits}%
}
\providecommand{\ssi}{si et seulement si} % abréviation fréquente
\AtBeginDocument{\providecommand{\wideparen}{\overparen}} % arc de cercle

%% FIGFIX-BEGIN v3 — correctifs globaux des figures (docs/onbuch-generation/tools/figfix)
\usepackage{adjustbox}\usepackage{etoolbox}
% toute figure TikZ/pgfplots plus large que la ligne est réduite à la largeur disponible
\BeforeBeginEnvironment{tikzpicture}{\begin{adjustbox}{max width=\linewidth}}
\AfterEndEnvironment{tikzpicture}{\end{adjustbox}}
% légendes pgfplots lisibles par-dessus les courbes
\pgfplotsset{every axis/.append style={legend style={fill=white,fill opacity=0.92,text opacity=1,draw=black!15,inner sep=3pt}}}
% étiquettes d'arêtes (syntaxe edge["..."]) avec fond blanc
\tikzset{every edge quotes/.append style={fill=white,inner sep=1.2pt}}
% bulles de cartes mentales : texte à la ligne dans la bulle, bulle un peu plus grande
\tikzset{every concept/.append style={text width=2.0cm,align=center,minimum size=2.3cm,inner sep=2pt}}
%% FIGFIX-END
\begin{document}

\pagegarde{Généralités sur les fonctions}{Mathématiques}{1ère C}{Relations et opérations fondamentales dans l'ensemble des nombres réels}{N04}{8 h}{Cette leçon pose les fondations de l'étude des fonctions numériques : ensemble de définition, opérations sur les fonctions et les polynômes, parité, périodicité, symétries et composition. Elle introduit aussi les notions d'applications injectives, surjectives et bijectives, outils indispensables pour toute l'année de Première et la Terminale.
\vspace{4pt}
\begin{multicols}{2}\small
\begin{itemize}
\item À la fin de cette leçon, je sais déterminer par le calcul l'ensemble de définition d'une fonction numérique
\item À la fin de cette leçon, je sais déterminer la restriction d'une fonction à un intervalle
\item À la fin de cette leçon, je sais justifier qu'une application est injective, surjective ou bijective et expliciter sa bijection réciproque
\item À la fin de cette leçon, je sais déterminer la composée de deux applications et son ensemble de définition
\item À la fin de cette leçon, je sais calculer la somme, le produit de deux polynômes et donner la condition d'existence d'un quotient de polynômes
\item À la fin de cette leçon, je sais déterminer l'ensemble de définition d'une somme, d'un produit, d'un quotient de deux fonctions
\end{itemize}\end{multicols}}

\begin{sommaire}
\begin{multicols}{2}
\begin{enumerate}
\item Fonctions numériques et ensemble de définition
\item Applications injectives, surjectives, bijectives
\item Composition des applications
\item Opérations sur les polynômes et les fonctions
\item Parité, périodicité et éléments de symétrie
\item Courbe d'une fonction et fonctions associées
\item Activité d'intégration
\item Méthodes et exercices résolus
\item Exercices
\item Corrigés des exercices
\item Fiche bilan
\end{enumerate}
\end{multicols}
\end{sommaire}

\begin{objectifs}
\begin{itemize}
\item À la fin de cette leçon, je sais déterminer par le calcul l'ensemble de définition d'une fonction numérique
\item À la fin de cette leçon, je sais déterminer la restriction d'une fonction à un intervalle
\item À la fin de cette leçon, je sais justifier qu'une application est injective, surjective ou bijective et expliciter sa bijection réciproque
\item À la fin de cette leçon, je sais déterminer la composée de deux applications et son ensemble de définition
\item À la fin de cette leçon, je sais calculer la somme, le produit de deux polynômes et donner la condition d'existence d'un quotient de polynômes
\item À la fin de cette leçon, je sais déterminer l'ensemble de définition d'une somme, d'un produit, d'un quotient de deux fonctions
\item À la fin de cette leçon, je sais montrer qu'une fonction est paire, impaire ou périodique
\item À la fin de cette leçon, je sais justifier qu'un point est centre de symétrie ou qu'une droite x = a est axe de symétrie d'une courbe
\item À la fin de cette leçon, je sais conjecturer par lecture graphique un ensemble de définition, des variations, des asymptotes et des éléments de symétrie
\end{itemize}
\end{objectifs}

\begin{prerequis}
\begin{itemize}
\item Ensembles de nombres (ℕ, ℤ, 𝔻, ℚ, ℝ) et intervalles de ℝ
\item Résolution d'équations et d'inéquations du premier et du second degré
\item Notion de fonction et calcul d'images et d'antécédents (Seconde)
\item Calcul littéral : développement, factorisation, identités remarquables
\item Repérage dans le plan et lecture de coordonnées de points
\end{itemize}
\end{prerequis}

\newpage

\coursec{Fonctions numériques et ensemble de définition}

\textbf{Situation d'accroche.} À la gare routière de Mokolo à Yaoundé, le prix d'une course en moto-taxi dépend de la distance parcourue : à chaque distance correspond \emph{un seul} prix. Un client qui parcourt \SI{2}{km} paiera par exemple 300 FCFA, et jamais deux montants différents pour la même distance. En revanche, deux clients partant de quartiers différents peuvent payer le même prix si leurs trajets ont la même longueur ! Comment modéliser mathématiquement cette correspondance ? Quand peut-on, à partir du prix payé, retrouver la distance parcourue sans ambiguïté ? Ces questions, au cœur de la vie quotidienne camerounaise, trouvent leurs réponses dans la notion de \cle{fonction}, puis d'\cle{injection}, de \cle{surjection} et de \cle{bijection}, que nous étudierons dans cette leçon. Commençons par le point de départ de toute l'analyse : la fonction numérique et son ensemble de définition.

\courssub{Fonction numérique d'une variable réelle : rappels}

\textbf{Intuition.} Une fonction est une \emph{machine} : on introduit un nombre $x$ en entrée, la machine effectue un calcul, et elle sort \emph{un et un seul} résultat, noté $f(x)$. Le prix de la course de moto-taxi est une fonction de la distance : à une distance donnée correspond un prix unique. En revanche, « le client qui a payé $p$ francs » n'est pas une fonction du prix, car plusieurs clients (plusieurs distances) peuvent correspondre au même prix.

\begin{definition}[Fonction numérique, image, antécédent]
Une \cle{fonction numérique d'une variable réelle} est un procédé qui, à certains nombres réels $x$, associe \textbf{un unique} nombre réel noté $f(x)$. On écrit :
\[ f : x \longmapsto f(x). \]
\begin{itemize}
\item Le nombre $f(x)$ est appelé l'\cle{image} de $x$ par $f$.
\item Si $y = f(x)$, on dit que $x$ est \textbf{un} \cle{antécédent} de $y$ par $f$.
\end{itemize}
\end{definition}

\begin{attention}[Image unique, antécédents multiples]
Tout réel $x$ de l'\textbf{ensemble de définition} $D_f$ possède \textbf{exactement une} image : c'est l'exigence fondamentale. En revanche, un nombre $y$ peut avoir \textbf{plusieurs antécédents} (comme deux distances donnant le même prix) ou \textbf{aucun}. Ne confonds jamais les deux notions : l'unicité porte sur l'\emph{image}, pas sur l'antécédent.
\end{attention}

\begin{popfigure}
\begin{tikzpicture}[scale=0.9, every node/.style={font=\small}]
% Diagramme 1 : pas une fonction
\draw[thick, popblue] (0,0.4) ellipse (1.1 and 1.7);
\draw[thick, poporange] (4.5,0.4) ellipse (1.1 and 1.7);
\node[popblue] at (0,2.4) {$E$};
\node[poporange] at (4.5,2.4) {$F$};
\fill[popdark] (-0.3,1.3) circle (1.5pt) node[left]{$a$};
\fill[popdark] (-0.3,0.4) circle (1.5pt) node[left]{$b$};
\fill[popdark] (-0.3,-0.5) circle (1.5pt) node[left]{$c$};
\fill[popdark] (4.8,1.3) circle (1.5pt) node[right]{$1$};
\fill[popdark] (4.8,0.4) circle (1.5pt) node[right]{$2$};
\fill[popdark] (4.8,-0.5) circle (1.5pt) node[right]{$3$};
\draw[-{Stealth}, thick, poppink] (0,1.3) -- (4.4,1.35);
\draw[-{Stealth}, thick, poppink] (0,1.25) -- (4.4,0.45);
\draw[-{Stealth}, thick, popgreen] (0,0.4) -- (4.4,0.35);
\node[align=center] at (2.25,-2) {Correspondance qui \textbf{n'est pas} une fonction :\\ $a$ a deux images et $c$ n'a aucune image};
% Diagramme 2 : fonction
\begin{scope}[xshift=8.5cm]
\draw[thick, popblue] (0,0.4) ellipse (1.1 and 1.7);
\draw[thick, poporange] (4.5,0.4) ellipse (1.1 and 1.7);
\node[popblue] at (0,2.4) {$E$};
\node[poporange] at (4.5,2.4) {$F$};
\fill[popdark] (-0.3,1.3) circle (1.5pt) node[left]{$a$};
\fill[popdark] (-0.3,0.4) circle (1.5pt) node[left]{$b$};
\fill[popdark] (-0.3,-0.5) circle (1.5pt) node[left]{$c$};
\fill[popdark] (4.8,1.3) circle (1.5pt) node[right]{$1$};
\fill[popdark] (4.8,0.4) circle (1.5pt) node[right]{$2$};
\fill[popdark] (4.8,-0.5) circle (1.5pt) node[right]{$3$};
\draw[-{Stealth}, thick, popgreen] (0,1.3) -- (4.4,1.35);
\draw[-{Stealth}, thick, popgreen] (0,0.4) -- (4.4,1.25);
\draw[-{Stealth}, thick, popgreen] (0,-0.5) -- (4.4,0.35);
\node[align=center] at (2.25,-2) {Une fonction : chaque élément de $E$\\ a \textbf{exactement une} image dans $F$};
\end{scope}
\end{tikzpicture}
\legende{Schémas sagittaux : à gauche, une correspondance qui n'est pas une fonction ; à droite, une fonction. Remarque que deux éléments ($a$ et $b$) peuvent avoir la même image.}
\end{popfigure}

\begin{exemplebox}[Lecture d'images et d'antécédents]
Soit $f$ la fonction définie par $f(x) = x^2 - 3$.
\begin{itemize}
\item L'image de $2$ est $f(2) = 2^2 - 3 = 1$.
\item L'image de $-1$ est $f(-1) = (-1)^2 - 3 = -2$.
\item Cherchons les antécédents de $6$ : on résout $x^2 - 3 = 6$, soit $x^2 = 9$, d'où $x = 3$ ou $x = -3$. Le nombre $6$ a \textbf{deux} antécédents.
\item Cherchons les antécédents de $-5$ : $x^2 - 3 = -5$ donne $x^2 = -2$, impossible dans $\mathbb{R}$ : $-5$ n'a \textbf{aucun} antécédent.
\end{itemize}
\end{exemplebox}

\courssub{Ensemble de définition d'une fonction numérique}

\textbf{Intuition.} La machine $f$ n'accepte pas forcément tous les nombres en entrée. Si tu introduis $x = 5$ dans la machine $f(x) = \dfrac{1}{x-5}$, la machine « casse » : la division par zéro est impossible. L'ensemble de définition est précisément la liste de tous les nombres que la machine accepte.

\begin{definition}[Ensemble de définition]
L'\cle{ensemble de définition} d'une fonction numérique $f$, noté $D_f$ (ou $\mathcal{D}_f$), est l'ensemble de tous les nombres réels $x$ pour lesquels $f(x)$ \textbf{existe}, c'est-à-dire pour lesquels le calcul de $f(x)$ est possible et donne un nombre réel.
\[ D_f = \{\, x \in \mathbb{R} \;|\; f(x) \text{ existe} \,\}. \]
\end{definition}

\begin{aretenir}[Les trois conditions d'existence fondamentales]
\begin{enumerate}
\item \textbf{Quotient :} un dénominateur ne doit \textbf{jamais être nul}. L'expression $\dfrac{A(x)}{B(x)}$ existe si et seulement si $B(x) \neq 0$.
\item \textbf{Racine carrée :} la quantité sous une racine carrée (le \emph{radicande}) doit être \textbf{positive ou nulle}. L'expression $\sqrt{A(x)}$ existe si et seulement si $A(x) \geq 0$.
\item \textbf{Polynôme :} un polynôme (somme de termes en $x^n$) est défini pour \textbf{tout} réel : $D_f = \mathbb{R}$.
\end{enumerate}
Quand plusieurs conditions apparaissent dans la même expression, elles doivent être vérifiées \textbf{simultanément}.
\end{aretenir}

\begin{attention}[Erreur fréquente : équation ou inéquation ?]
La condition sur un \textbf{dénominateur} est une \textbf{équation} : on résout $B(x) = 0$ et l'on \emph{exclut} les solutions (ce sont les \cle{valeurs interdites}). La condition sur un \textbf{radicande} est une \textbf{inéquation} : on résout $A(x) \geq 0$ et l'on \emph{garde} les solutions. Écrire « $x - 5 \geq 0$ » pour un dénominateur, ou « $\sqrt{2x-4}$ existe si $2x - 4 \neq 0$ », sont deux erreurs classiques qui coûtent cher au Probatoire !
\end{attention}

\begin{methode}[Déterminer l'ensemble de définition d'une fonction]
\begin{enumerate}
\item \textbf{Repérer} dans l'expression de $f(x)$ tous les dénominateurs et toutes les racines carrées.
\item \textbf{Traduire} chaque contrainte : dénominateur $\neq 0$ (équation à résoudre, valeurs à exclure) ; radicande $\geq 0$ (inéquation à résoudre, ensemble à garder).
\item \textbf{Résoudre} chaque équation ou inéquation dans $\mathbb{R}$.
\item \textbf{Combiner} les conditions : on prend l'intersection des ensembles de validité, en excluant les valeurs interdites.
\item \textbf{Conclure} en écrivant $D_f$ comme un intervalle ou une \textbf{réunion d'intervalles}.
\end{enumerate}
\end{methode}

\begin{popfigure}
\begin{tikzpicture}[
  grow=right, level distance=4.2cm,
  level 1/.style={sibling distance=2.6cm},
  level 2/.style={sibling distance=1.6cm},
  every node/.style={font=\small},
  edge from parent/.style={draw=popmuted, thick, -{Stealth}}
]
\node[draw, thick, fill=popblueL, rounded corners, align=center] {Forme de $f(x)$ ?}
  child { node[draw, fill=popgreenL, rounded corners, align=center] {Racine carrée\\ $\sqrt{A(x)}$}
    child { node[align=center] {$A(x) \geq 0$\\ (inéquation)} } }
  child { node[draw, fill=poporangeL, rounded corners, align=center] {Quotient\\ $\dfrac{A(x)}{B(x)}$}
    child { node[align=center] {$B(x) \neq 0$\\ (équation, exclure)} } }
  child { node[draw, fill=poppurpleL, rounded corners, align=center] {Polynôme}
    child { node[align=center] {$D_f = \mathbb{R}$} } };
\end{tikzpicture}
\legende{Arbre de décision pour déterminer $D_f$ selon la forme de l'expression de $f$. En cas de combinaison (quotient sous une racine, racine au dénominateur...), on applique toutes les règles concernées.}
\end{popfigure}

\begin{exoresolu}[Détermination complète d'un ensemble de définition]
Déterminer l'ensemble de définition de la fonction $f$ définie par
\[ f(x) = \dfrac{\sqrt{2x - 4}}{x - 5}. \]
\tcblower
\textbf{Étape 1 — Repérage.} L'expression contient une racine carrée ($\sqrt{2x-4}$) et un dénominateur ($x - 5$). Il y a donc deux conditions.

\textbf{Étape 2 — Traduction.}
\begin{itemize}
\item Condition sur le radicande (inéquation) : $2x - 4 \geq 0$.
\item Condition sur le dénominateur (équation, valeur à exclure) : $x - 5 \neq 0$.
\end{itemize}

\textbf{Étape 3 — Résolution.}
\begin{align*}
2x - 4 \geq 0 &\iff 2x \geq 4 \iff x \geq 2, \quad \text{soit } x \in [2\,;\,+\infty[.\\
x - 5 \neq 0 &\iff x \neq 5.
\end{align*}

\textbf{Étape 4 — Combinaison.} On garde l'intervalle $[2\,;\,+\infty[$ et l'on en exclut la valeur interdite $5$, qui appartient à cet intervalle.

\textbf{Étape 5 — Conclusion.}
\[ \boxed{D_f = [2\,;\,5[ \;\cup\; ]5\,;\,+\infty[} \]
Vérification : $f(2) = \dfrac{\sqrt{0}}{-3} = 0$ existe bien (donc $2 \in D_f$) ; $f(5)$ n'existe pas (dénominateur nul) ; $f(1)$ n'existe pas (radicande $-2 < 0$).
\end{exoresolu}

\begin{exemplebox}[Autres situations types]
\begin{itemize}
\item $g(x) = 3x^3 - 2x + 7$ est un polynôme : $D_g = \mathbb{R} = ]-\infty\,;\,+\infty[$.
\item $h(x) = \dfrac{4x+1}{x^2-9}$ : on résout $x^2 - 9 = 0 \iff x = 3$ ou $x = -3$. Donc $D_h = \mathbb{R} \setminus \{-3\,;\,3\} = ]-\infty\,;\,-3[ \cup ]-3\,;\,3[ \cup ]3\,;\,+\infty[$.
\item $k(x) = \sqrt{5 - x}$ : $5 - x \geq 0 \iff x \leq 5$, donc $D_k = ]-\infty\,;\,5]$.
\end{itemize}
\end{exemplebox}

\begin{attention}[Rédaction de la conclusion]
La conclusion doit être un ensemble de nombres réels, écrit proprement avec des intervalles. Écris $D_f = [2\,;\,5[ \cup ]5\,;\,+\infty[$ et non « $x \geq 2$ et $x \neq 5$ » seul : la phrase de condition doit être traduite en ensemble. Attention aussi aux crochets : la valeur $5$ est exclue, donc les crochets sont ouverts en $5$ ; la valeur $2$ est admise, donc le crochet est fermé en $2$.
\end{attention}

\courssub{Restriction d'une fonction}

\textbf{Intuition.} Le tarif du moto-taxi de Mokolo n'est valable que pour des courses de $0$ à \SI{15}{km} : au-delà, un autre barème s'applique. On peut donc regarder la fonction « prix » \emph{uniquement} sur l'intervalle $[0\,;\,15]$, même si la formule de calcul aurait un sens ailleurs. C'est exactement l'idée de restriction.

\begin{definition}[Restriction d'une fonction]
Soit $f$ une fonction d'ensemble de définition $D_f$ et $I$ une partie de $D_f$ (par exemple un intervalle $I \subset D_f$). La \cle{restriction} de $f$ à $I$, notée $f_{|I}$, est la fonction définie sur $I$ par :
\[ f_{|I} : x \longmapsto f(x), \quad \text{pour tout } x \in I. \]
Autrement dit : \textbf{même formule, ensemble de départ réduit}.
\end{definition}

\begin{exemplebox}[Restriction à un intervalle]
Soit $f$ la fonction définie sur $\mathbb{R}$ par $f(x) = x^2$. Sa restriction à $I = [0\,;\,+\infty[$ est la fonction
\[ f_{|I} : [0\,;\,+\infty[ \longrightarrow \mathbb{R}, \quad x \longmapsto x^2. \]
Elle ne calcule les carrés que des nombres positifs : $f_{|I}(3) = 9$, mais $f_{|I}(-2)$ n'a \textbf{aucun sens}, car $-2 \notin I$. Nous verrons dans la section suivante que cette restriction est injective, alors que $f$ ne l'est pas : restreindre une fonction est un outil précieux pour lui donner de meilleures propriétés.
\end{exemplebox}

\begin{methode}[Déterminer la restriction d'une fonction sur un intervalle]
\begin{enumerate}
\item Vérifier que l'intervalle $I$ proposé est bien inclus dans $D_f$ : $I \subset D_f$.
\item Conserver la formule de $f$ sans la modifier.
\item Préciser le nouvel ensemble de départ : $f_{|I}$ est définie sur $I$ (et non plus sur $D_f$).
\end{enumerate}
\end{methode}

\begin{attention}[Restriction et ensemble de définition]
On ne peut restreindre $f$ qu'à une partie de son ensemble de définition. Par exemple, pour $f(x) = \dfrac{1}{x}$ ($D_f = \mathbb{R}^*$), la restriction à $]0\,;\,+\infty[$ existe, mais parler de « restriction de $f$ à $[-1\,;\,1]$ » est incorrect, car $0 \in [-1\,;\,1]$ et $f(0)$ n'existe pas.
\end{attention}

\courssub{Égalité de deux fonctions}

\textbf{Intuition.} Deux machines peuvent avoir la même formule de calcul mais accepter des entrées différentes : ce ne sont alors pas les mêmes machines. Pour que deux fonctions soient \emph{égales}, il faut à la fois qu'elles acceptent les mêmes nombres et qu'elles donnent les mêmes résultats.

\begin{definition}[Égalité de deux fonctions]
Deux fonctions $f$ et $g$ sont \cle{égales}, et l'on note $f = g$, si et seulement si les deux conditions suivantes sont réunies :
\begin{enumerate}
\item elles ont le \textbf{même ensemble de définition} : $D_f = D_g$ ;
\item pour tout $x$ de cet ensemble commun, elles donnent la \textbf{même image} : $f(x) = g(x)$.
\end{enumerate}
\end{definition}

\begin{exemplebox}[Deux fonctions qui ne sont pas égales]
Considérons
\[ f(x) = \dfrac{x^2 - 1}{x - 1} \qquad \text{et} \qquad g(x) = x + 1. \]
Pour tout $x \neq 1$, on peut factoriser : $f(x) = \dfrac{(x-1)(x+1)}{x-1} = x + 1 = g(x)$. Les images coïncident donc partout où $f$ existe. Pourtant $f \neq g$, car $D_f = \mathbb{R} \setminus \{1\}$ alors que $D_g = \mathbb{R}$ : la valeur $1$ est interdite pour $f$ mais pas pour $g$. La première condition de la définition n'est pas satisfaite.
\end{exemplebox}

\begin{attention}[Simplifier ne suffit pas]
Simplifier l'écriture d'une fonction (factoriser, réduire un quotient) ne change pas son ensemble de définition. Avant toute simplification de $\dfrac{(x-1)(x+1)}{x-1}$ en $x+1$, pense à la condition $x \neq 1$, qui demeure même après simplification. C'est une faute grave d'affirmer que les deux fonctions de l'exemple précédent sont égales.
\end{attention}

\begin{exoresolu}[Images, antécédents et appartenance]
Soit $f$ la fonction définie par $f(x) = \dfrac{2x+3}{x+2}$.
\begin{enumerate}
\item Déterminer $D_f$.
\item Calculer l'image de $-1$.
\item Déterminer les antécédents éventuels de $2$.
\end{enumerate}
\tcblower
\textbf{1.} Le dénominateur doit être non nul : $x + 2 \neq 0 \iff x \neq -2$. Donc $D_f = \mathbb{R} \setminus \{-2\} = ]-\infty\,;\,-2[ \cup ]-2\,;\,+\infty[$.

\textbf{2.} $f(-1) = \dfrac{2\times(-1)+3}{-1+2} = \dfrac{1}{1} = 1$. L'image de $-1$ est $1$.

\textbf{3.} On résout dans $D_f$ l'équation $f(x) = 2$ :
\begin{align*}
\dfrac{2x+3}{x+2} = 2 &\iff 2x + 3 = 2(x+2) \quad (\text{car } x \neq -2)\\
&\iff 2x + 3 = 2x + 4 \iff 3 = 4,
\end{align*}
ce qui est impossible. Le nombre $2$ n'a donc \textbf{aucun antécédent} : il n'est atteint par aucune valeur de $x$. (Nous dirons plus tard que $f$ n'est pas surjective sur $\mathbb{R}$.)
\end{exoresolu}

\begin{savaistu}[D'où vient la notation $f(x)$ ?]
La notation $f(x)$ a été introduite par le mathématicien suisse \textbf{Leonhard Euler} (1707--1783) au XVIII\textsuperscript{e} siècle, dans son monumental traité \emph{Introductio in analysin infinitorum} (1748). Euler est l'un des mathématiciens les plus prolifiques de l'histoire : on lui doit aussi la notation $e$ pour le nombre d'Euler, $i$ pour l'imaginaire, et le symbole $\sum$ pour les sommes. Avant lui, on décrivait les fonctions par de longues phrases en latin ; sa notation a rendu le calcul fonctionnel aussi naturel que le calcul sur les nombres, et c'est elle que tu utilises chaque jour en classe !
\end{savaistu}

\begin{exercice}[À toi de jouer]
\begin{enumerate}
\item Déterminer les ensembles de définition des fonctions suivantes :
\[ f(x) = \dfrac{3x-1}{2x+6}, \qquad g(x) = \sqrt{4-2x}, \qquad h(x) = \dfrac{\sqrt{x+3}}{x-1}. \]
\item Soit $f(x) = x^2 - 5x + 6$ définie sur $\mathbb{R}$. Déterminer les antécédents de $0$, puis écrire la restriction de $f$ à l'intervalle $[3\,;\,+\infty[$.
\item Les fonctions $f(x) = \dfrac{x^2}{x}$ et $g(x) = x$ sont-elles égales ? Justifier.
\end{enumerate}
\end{exercice}

\begin{corrige}[Exercice n°1]
\textbf{1.} Pour $f$ : $2x + 6 \neq 0 \iff x \neq -3$, donc $D_f = \mathbb{R} \setminus \{-3\} = ]-\infty\,;\,-3[ \cup ]-3\,;\,+\infty[$.

Pour $g$ : $4 - 2x \geq 0 \iff x \leq 2$, donc $D_g = ]-\infty\,;\,2]$.

Pour $h$ : deux conditions simultanées : $x + 3 \geq 0 \iff x \geq -3$ et $x - 1 \neq 0 \iff x \neq 1$. Donc $D_h = [-3\,;\,1[ \cup ]1\,;\,+\infty[$.

\textbf{2.} $f(x) = 0 \iff x^2 - 5x + 6 = 0$. Le discriminant est $\Delta = 25 - 24 = 1$, d'où $x = 2$ ou $x = 3$ : les antécédents de $0$ sont $2$ et $3$. La restriction est $f_{|[3;+\infty[} : x \mapsto x^2 - 5x + 6$, définie uniquement sur $[3\,;\,+\infty[$.

\textbf{3.} Non. On a $D_f = \mathbb{R} \setminus \{0\}$ (dénominateur $x$) alors que $D_g = \mathbb{R}$. Même si $f(x) = g(x)$ pour tout $x \neq 0$, les ensembles de définition diffèrent, donc $f \neq g$.
\end{corrige}

\coursec{Applications injectives, surjectives, bijectives}

Dans la section précédente, nous avons appris à déterminer l'\cle{ensemble de définition} d'une fonction numérique et à restreindre une fonction à une partie de son ensemble de départ. Nous savons donc désormais \emph{où} une fonction agit. Il reste une question tout aussi fondamentale : \emph{comment} agit-elle ? Deux réels distincts peuvent-ils avoir la même image ? Tous les réels de l'ensemble d'arrivée sont-ils atteints ? Ces questions, qui paraissent abstraites, sont en réalité celles que se pose chaque jour un opérateur téléphonique : deux abonnés distincts doivent-ils avoir des numéros distincts (oui, sinon les appels se mélangent !), et tout numéro attribué correspond-il bien à un abonné ? C'est exactement le vocabulaire des applications \cle{injectives}, \cle{surjectives} et \cle{bijectives} que nous allons construire dans cette section.

\courssub{La notion d'application}

\begin{definition}[Application]
Soient $E$ et $F$ deux ensembles non vides. On appelle \cle{application} de $E$ vers $F$ toute correspondance $f$ qui, à \textbf{chaque} élément $x$ de $E$, associe \textbf{un seul} élément de $F$, noté $f(x)$ et appelé \emph{image} de $x$ par $f$. On note :
\[ f : E \longrightarrow F \qquad \text{ou} \qquad \begin{array}{rcl} f : E & \longrightarrow & F \\ x & \longmapsto & f(x) \end{array} \]
Si $y = f(x)$, on dit que $x$ est un \cle{antécédent} de $y$ par $f$.
\end{definition}

Retiens bien les deux exigences cachées dans cette définition :
\begin{itemize}
\item \textbf{Existence} : tout élément de $E$ a une image (aucun élément du départ n'est « orphelin ») ;
\item \textbf{Unicité} : cette image est unique (un élément ne peut pas « se dédoubler »).
\end{itemize}

En revanche, rien n'interdit \emph{a priori} que deux éléments de $E$ aient la même image, ni qu'un élément de $F$ n'ait aucun antécédent. Toute la richesse de cette section vient de l'étude de ces deux phénomènes.

\begin{attention}[Fonction ou application ?]
Une \emph{fonction} numérique peut n'être définie que sur une partie de $\mathbb{R}$ (son ensemble de définition). Parler d'\emph{application} de $E$ vers $F$ exige que $E$ soit entièrement couvert. Ainsi $x \mapsto \sqrt{x}$ n'est pas une application de $\mathbb{R}$ dans $\mathbb{R}$, mais c'est bien une application de $[0\,;\,+\infty[$ dans $\mathbb{R}$.
\end{attention}

\courssub{Applications injectives}

\textbf{Intuition.} Imagine la liste des numéros matricules des élèves d'une classe de Première C au lycée de Ngaoundéré. Deux élèves distincts ne peuvent pas porter le même matricule : sinon, impossible de savoir qui a composé à l'examen ! La correspondance « élève $\mapsto$ matricule » ne « confond » jamais deux élèves. On dit qu'elle est \emph{injective}.

\begin{definition}[Application injective]
Une application $f : E \to F$ est dite \cle{injective} lorsque deux éléments \textbf{distincts} de $E$ ont toujours des images \textbf{distinctes} :
\[ \text{pour tous } a, b \in E, \quad a \neq b \implies f(a) \neq f(b). \]
De manière équivalente (contraposée), $f$ est injective si et seulement si :
\[ \text{pour tous } a, b \in E, \quad f(a) = f(b) \implies a = b. \]
Autrement dit : \textbf{tout élément de $F$ a au plus un antécédent} (c'est-à-dire un seul, ou aucun).
\end{definition}

\begin{methode}[Justifier qu'une application est injective]
Pour montrer qu'une application $f : E \to F$ est injective :
\begin{enumerate}
\item On part de l'hypothèse $f(a) = f(b)$ avec $a, b \in E$ ;
\item on manipule cette égalité (calcul algébrique) jusqu'à obtenir $a = b$ ;
\item on conclut : « donc $f$ est injective ».
\end{enumerate}
Pour montrer qu'une application n'est \textbf{pas} injective, il suffit d'exhiber \emph{un seul} contre-exemple : deux éléments distincts $a \neq b$ tels que $f(a) = f(b)$.
\end{methode}

\begin{exemplebox}[Injectivité et non-injectivité]
\begin{itemize}
\item Soit $f : \mathbb{R} \to \mathbb{R}$ définie par $f(x) = 3x - 2$. Supposons $f(a) = f(b)$. Alors $3a - 2 = 3b - 2$, donc $3a = 3b$, donc $a = b$. L'application $f$ est \textbf{injective}.
\item Soit $g : \mathbb{R} \to \mathbb{R}$ définie par $g(x) = x^2$. On a $g(-2) = 4 = g(2)$ avec $-2 \neq 2$. L'application $g$ n'est \textbf{pas injective}.
\end{itemize}
\end{exemplebox}

\courssub{Applications surjectives}

\textbf{Intuition.} Reprenons notre classe de Première. La correspondance « élève $\mapsto$ matricule » est injective, mais il existe des numéros matricules (ceux des élèves de Terminale, par exemple) qui ne sont le matricule d'\emph{aucun} élève de cette classe : l'ensemble d'arrivée n'est pas entièrement atteint. À l'inverse, la correspondance « élève de la classe $\mapsto$ sa salle de composition » atteint toutes les salles prévues si chaque salle reçoit au moins un élève : on dit alors qu'elle est \emph{surjective}.

\begin{definition}[Application surjective]
Une application $f : E \to F$ est dite \cle{surjective} lorsque \textbf{tout} élément de $F$ possède \textbf{au moins un} antécédent dans $E$ :
\[ \text{pour tout } y \in F, \quad \text{il existe } x \in E \text{ tel que } f(x) = y. \]
Autrement dit, l'équation $f(x) = y$, d'inconnue $x \in E$, admet \textbf{au moins une solution} pour chaque $y \in F$.
\end{definition}

\begin{methode}[Justifier qu'une application est surjective]
Pour montrer que $f : E \to F$ est surjective :
\begin{enumerate}
\item On fixe un élément \textbf{quelconque} $y \in F$ ;
\item on résout l'équation $f(x) = y$ d'inconnue $x$ ;
\item on vérifie que la solution trouvée appartient bien à $E$ (c'est le point le plus souvent oublié !) ;
\item on conclut : « tout $y \in F$ admet un antécédent, donc $f$ est surjective ».
\end{enumerate}
Pour montrer que $f$ n'est \textbf{pas} surjective, il suffit d'exhiber \emph{un seul} élément de $F$ sans antécédent.
\end{methode}

\begin{exemplebox}[Surjectivité et non-surjectivité]
\begin{itemize}
\item Soit $f : \mathbb{R} \to \mathbb{R}$, $f(x) = 3x - 2$. Pour $y \in \mathbb{R}$ quelconque, $f(x) = y \iff 3x - 2 = y \iff x = \dfrac{y+2}{3}$, qui est bien un réel. Donc $f$ est \textbf{surjective}.
\item Soit $g : \mathbb{R} \to \mathbb{R}$, $g(x) = x^2$. L'équation $x^2 = -1$ n'a pas de solution réelle : $-1$ n'a aucun antécédent. Donc $g$ n'est \textbf{pas surjective} de $\mathbb{R}$ dans $\mathbb{R}$. En revanche, $g : \mathbb{R} \to [0\,;\,+\infty[$, $x \mapsto x^2$ \emph{est} surjective : tout $y \geq 0$ admet $\sqrt{y}$ comme antécédent.
\end{itemize}
\textbf{Leçon importante :} la surjectivité dépend de l'ensemble d'arrivée choisi !
\end{exemplebox}

\courssub{Applications bijectives}

\begin{definition}[Application bijective]
Une application $f : E \to F$ est dite \cle{bijective} lorsqu'elle est à la fois \textbf{injective} et \textbf{surjective}. Cela revient à dire que :
\[ \text{tout élément } y \in F \text{ possède un \textbf{unique} antécédent } x \in E. \]
On dit aussi que $f$ réalise une \cle{bijection} de $E$ sur $F$.
\end{definition}

Une bijection établit donc une correspondance parfaite, « un pour un », entre $E$ et $F$ : rien n'est confondu (injectivité), rien n'est oublié (surjectivité).

\begin{methode}[Justifier qu'une application est bijective — méthode fondamentale]
La méthode la plus efficace, exigible au Probatoire, consiste à résoudre directement l'équation $f(x) = y$ :
\begin{enumerate}
\item Fixer $y \in F$ quelconque et résoudre $f(x) = y$ d'inconnue $x \in E$ ;
\item montrer que l'équation admet une solution (cela prouve la \textbf{surjectivité}) ;
\item montrer que cette solution est \textbf{unique} et appartient à $E$ (cela prouve l'\textbf{injectivité}) ;
\item conclure : « pour tout $y \in F$, l'équation $f(x) = y$ admet une unique solution dans $E$, donc $f$ est bijective ».
\end{enumerate}
\end{methode}

\begin{exoresolu}[Bijectivité de $f(x) = 3x - 2$]
Montrer que l'application $f : \mathbb{R} \to \mathbb{R}$ définie par $f(x) = 3x - 2$ est bijective.
\tcblower
Soit $y \in \mathbb{R}$ quelconque. Résolvons dans $\mathbb{R}$ l'équation $f(x) = y$ :
\begin{align*}
f(x) = y &\iff 3x - 2 = y \\
&\iff 3x = y + 2 \\
&\iff x = \dfrac{y + 2}{3}.
\end{align*}
Pour tout réel $y$, le nombre $\dfrac{y+2}{3}$ existe, est un réel (donc appartient à l'ensemble de départ) et est \textbf{unique}. Ainsi, tout élément de l'ensemble d'arrivée possède un unique antécédent : l'application $f$ est \textbf{bijective} de $\mathbb{R}$ dans $\mathbb{R}$.
\end{exoresolu}

\begin{popfigure}
\begin{center}
\begin{tikzpicture}[scale=0.72, every node/.style={font=\small}]
% --- Diagramme 1 : injective non surjective ---
\begin{scope}
\draw[fill=popblueL, draw=popblue] (0,0) ellipse (0.9 and 1.5);
\draw[fill=poporangeL, draw=poporange] (3.2,0) ellipse (0.9 and 1.5);
\node[popblue] at (0,1.9) {$E$}; \node[poporange] at (3.2,1.9) {$F$};
\node at (0,0.8) {$a$}; \node at (0,0) {$b$}; \node at (0,-0.8) {$c$};
\node at (3.2,1.1) {$1$}; \node at (3.2,0.35) {$2$}; \node at (3.2,-0.35) {$3$}; \node at (3.2,-1.1) {$4$};
\draw[-{Stealth},popdark] (0.25,0.8) -- (2.95,1.1);
\draw[-{Stealth},popdark] (0.25,0) -- (2.95,0.35);
\draw[-{Stealth},popdark] (0.25,-0.8) -- (2.95,-0.35);
\node at (1.6,-2.1) {\textbf{Injective, non surjective}};
\end{scope}
% --- Diagramme 2 : surjective non injective ---
\begin{scope}[xshift=6.5cm]
\draw[fill=popblueL, draw=popblue] (0,0) ellipse (0.9 and 1.5);
\draw[fill=poporangeL, draw=poporange] (3.2,0) ellipse (0.9 and 1.5);
\node[popblue] at (0,1.9) {$E$}; \node[poporange] at (3.2,1.9) {$F$};
\node at (0,0.8) {$a$}; \node at (0,0) {$b$}; \node at (0,-0.8) {$c$};
\node at (3.2,0.5) {$1$}; \node at (3.2,-0.5) {$2$};
\draw[-{Stealth},popdark] (0.25,0.8) -- (2.95,0.55);
\draw[-{Stealth},popdark] (0.25,0) -- (2.95,0.45);
\draw[-{Stealth},popdark] (0.25,-0.8) -- (2.95,-0.5);
\node at (1.6,-2.1) {\textbf{Surjective, non injective}};
\end{scope}
% --- Diagramme 3 : bijective ---
\begin{scope}[yshift=-5.2cm]
\draw[fill=popblueL, draw=popblue] (0,0) ellipse (0.9 and 1.5);
\draw[fill=poporangeL, draw=poporange] (3.2,0) ellipse (0.9 and 1.5);
\node[popblue] at (0,1.9) {$E$}; \node[poporange] at (3.2,1.9) {$F$};
\node at (0,0.8) {$a$}; \node at (0,0) {$b$}; \node at (0,-0.8) {$c$};
\node at (3.2,0.8) {$1$}; \node at (3.2,0) {$2$}; \node at (3.2,-0.8) {$3$};
\draw[-{Stealth},popdark] (0.25,0.8) -- (2.95,0.8);
\draw[-{Stealth},popdark] (0.25,0) -- (2.95,0);
\draw[-{Stealth},popdark] (0.25,-0.8) -- (2.95,-0.8);
\node at (1.6,-2.1) {\textbf{Bijective}};
\end{scope}
% --- Diagramme 4 : non injective, non surjective ---
\begin{scope}[xshift=6.5cm, yshift=-5.2cm]
\draw[fill=popblueL, draw=popblue] (0,0) ellipse (0.9 and 1.5);
\draw[fill=poporangeL, draw=poporange] (3.2,0) ellipse (0.9 and 1.5);
\node[popblue] at (0,1.9) {$E$}; \node[poporange] at (3.2,1.9) {$F$};
\node at (0,0.8) {$a$}; \node at (0,0) {$b$}; \node at (0,-0.8) {$c$};
\node at (3.2,0.8) {$1$}; \node at (3.2,0) {$2$}; \node at (3.2,-0.8) {$3$};
\draw[-{Stealth},popdark] (0.25,0.8) -- (2.95,0.85);
\draw[-{Stealth},popdark] (0.25,0) -- (2.95,0.75);
\draw[-{Stealth},popdark] (0.25,-0.8) -- (2.95,-0.8);
\node at (1.6,-2.1) {\textbf{Ni injective, ni surjective}};
\end{scope}
\end{tikzpicture}
\end{center}
\legende{Quatre situations types : l'injectivité interdit que deux flèches arrivent au même point ; la surjectivité exige que chaque élément de $F$ soit atteint.}
\end{popfigure}

\courssub{Bijection réciproque}

\textbf{Intuition.} Si $f$ est bijective, chaque $y$ de $F$ provient d'un \emph{unique} $x$ de $E$. On peut donc « remonter les flèches » sans aucune ambiguïté : c'est cette remontée qui définit la \cle{bijection réciproque}.

\begin{definition}[Bijection réciproque]
Soit $f : E \to F$ une application bijective. L'application qui, à chaque $y \in F$, associe son \textbf{unique} antécédent par $f$ s'appelle la \cle{bijection réciproque} de $f$ et se note $f^{-1}$. C'est une application de $F$ vers $E$, caractérisée par :
\[ y = f(x) \iff x = f^{-1}(y). \]
On a alors :
\[ f \circ f^{-1} = \mathrm{id}_F \qquad \text{et} \qquad f^{-1} \circ f = \mathrm{id}_E, \]
où $\mathrm{id}_E$ et $\mathrm{id}_F$ désignent les applications identité ($\mathrm{id}_E(x) = x$ pour tout $x \in E$). La notation $f \circ g$ désigne la \cle{composée} des applications, que nous étudierons en détail dans la section suivante : $(f \circ g)(x) = f\bigl(g(x)\bigr)$.
\end{definition}

\begin{propriete}[La réciproque d'une bijection est une bijection]
Si $f : E \to F$ est bijective, alors $f^{-1} : F \to E$ est également bijective, et sa réciproque est $f$ : $(f^{-1})^{-1} = f$.
\emph{Cette propriété, admise, est évidente sur les définitions : la correspondance « un pour un » entre $E$ et $F$ fonctionne dans les deux sens.}
\end{propriete}

\begin{methode}[Expliciter la bijection réciproque]
Pour déterminer l'expression de $f^{-1}$ lorsque $f$ est bijective :
\begin{enumerate}
\item Poser $y = f(x)$ ;
\item \textbf{isoler $x$} dans cette égalité, c'est-à-dire exprimer $x$ en fonction de $y$ ;
\item vérifier que pour chaque $y$, la valeur obtenue est \textbf{unique} et appartient à l'ensemble de départ de $f$ ;
\item conclure en écrivant $f^{-1}(y) = \ldots$, puis, par habitude, renommer la variable : $f^{-1}(x) = \ldots$.
\end{enumerate}
\end{methode}

\begin{exoresolu}[Expliciter une bijection réciproque]
On a montré que $f : \mathbb{R} \to \mathbb{R}$, $f(x) = 3x - 2$ est bijective. Déterminer $f^{-1}$.
\tcblower
Posons $y = f(x)$, c'est-à-dire $y = 3x - 2$. Isolons $x$ :
\begin{align*}
y = 3x - 2 &\iff 3x = y + 2 \\
&\iff x = \dfrac{y + 2}{3}.
\end{align*}
Pour tout réel $y$, cette valeur existe et est unique. Donc :
\[ f^{-1}(y) = \dfrac{y + 2}{3}, \qquad \text{c'est-à-dire} \qquad \boxed{f^{-1}(x) = \dfrac{x + 2}{3}}. \]
\textbf{Vérification} (toujours profitable !) : pour tout réel $x$,
\[ f\left(f^{-1}(x)\right) = 3 \times \dfrac{x+2}{3} - 2 = x + 2 - 2 = x \quad \text{et} \quad f^{-1}\left(f(x)\right) = \dfrac{(3x-2)+2}{3} = x. \]
On retrouve bien $f \circ f^{-1} = f^{-1} \circ f = \mathrm{id}_{\mathbb{R}}$.
\end{exoresolu}

\begin{attention}[Ne pas confondre $f^{-1}(x)$ et $\dfrac{1}{f(x)}$]
C'est l'erreur la plus fréquente au Probatoire ! Le symbole $f^{-1}$ désigne la \textbf{bijection réciproque}, pas l'inverse. Pour $f(x) = 3x - 2$ :
\[ f^{-1}(x) = \dfrac{x+2}{3} \qquad \text{alors que} \qquad \dfrac{1}{f(x)} = \dfrac{1}{3x-2}. \]
Ces deux expressions n'ont rien à voir. De plus, $f^{-1}$ n'existe que si $f$ est \textbf{bijective} : écrire $f^{-1}$ pour une fonction non bijective est une faute grave.
\end{attention}

\begin{popfigure}
\begin{center}
\begin{tikzpicture}
\begin{axis}[
  width=0.85\linewidth, height=7.5cm,
  axis lines=middle, axis line style={-{Stealth}},
  xlabel={$x$}, ylabel={$y$},
  xmin=-4, xmax=4, ymin=-4, ymax=4,
  xtick={-3,-2,-1,1,2,3}, ytick={-3,-2,-1,1,2,3},
  grid=both, grid style={popline!40},
  legend style={at={(0.02,0.98)}, anchor=north west, font=\small}
]
\addplot[popcurve, domain=-1.7:1.2, samples=100] {2*x+1};
\addlegendentry{$y = f(x) = 2x+1$}
\addplot[popcurve, poporange, domain=-2.5:3.4, samples=100] {(x-1)/2};
\addlegendentry{$y = f^{-1}(x) = \dfrac{x-1}{2}$}
\addplot[popmuted, dashed, domain=-3.8:3.8, samples=2] {x};
\addlegendentry{$y = x$}
\end{axis}
\end{tikzpicture}
\end{center}
\legende{Les courbes de $f(x) = 2x+1$ et de sa réciproque $f^{-1}(x) = \dfrac{x-1}{2}$ sont symétriques par rapport à la première bissectrice $y = x$.}
\end{popfigure}

\begin{aretenir}[Symétrie des courbes de $f$ et $f^{-1}$]
Dans un repère orthonormé, les courbes représentatives d'une bijection $f$ et de sa réciproque $f^{-1}$ sont \textbf{symétriques par rapport à la droite d'équation $y = x$}. En effet, le point $(x\,;\,y)$ appartient à la courbe de $f$ si et seulement si le point $(y\,;\,x)$ appartient à celle de $f^{-1}$, et ces deux points sont symétriques par rapport à la première bissectrice.
\end{aretenir}

\courssub{Restriction et bijectivité : rendre une fonction bijective}

Nous avons vu que $g : \mathbb{R} \to \mathbb{R}$, $x \mapsto x^2$ n'est ni injective (car $g(-2) = g(2)$) ni surjective (car $-1$ n'a pas d'antécédent). Faut-il renoncer à « inverser » cette fonction ? Non : il suffit de \textbf{restreindre} les ensembles de départ et d'arrivée, notion vue dans la section précédente.

\begin{exemplebox}[La fonction carré devient bijective par restriction]
Considérons la restriction de la fonction carré :
\[ h : [0\,;\,+\infty[ \longrightarrow [0\,;\,+\infty[, \qquad x \longmapsto x^2. \]
Soit $y \in [0\,;\,+\infty[$. Résolvons $x^2 = y$ avec $x \geq 0$ : les solutions dans $\mathbb{R}$ sont $\sqrt{y}$ et $-\sqrt{y}$, mais seule $x = \sqrt{y}$ appartient à $[0\,;\,+\infty[$. L'équation admet donc une \textbf{unique} solution dans l'ensemble de départ : $h$ est bijective, et
\[ h^{-1} : [0\,;\,+\infty[ \longrightarrow [0\,;\,+\infty[, \qquad x \longmapsto \sqrt{x}. \]
C'est précisément ainsi qu'est « née » la fonction racine carrée ! Tu retrouveras ce procédé en Terminale pour construire les fonctions réciproques trigonométriques.
\end{exemplebox}

\begin{aretenir}[L'essentiel de la section]
\begin{itemize}
\item $f$ est \textbf{injective} : $f(a) = f(b) \implies a = b$ (chaque image a au plus un antécédent) ;
\item $f$ est \textbf{surjective} : tout $y \in F$ a au moins un antécédent (l'équation $f(x) = y$ a toujours une solution) ;
\item $f$ est \textbf{bijective} : injective \emph{et} surjective, c'est-à-dire que $f(x) = y$ a, pour chaque $y \in F$, une \textbf{unique} solution dans $E$ ;
\item si $f$ est bijective, sa réciproque $f^{-1}$ s'obtient en \textbf{isolant $x$} dans l'égalité $y = f(x)$, et les courbes de $f$ et $f^{-1}$ sont symétriques par rapport à la droite $y = x$.
\end{itemize}
\end{aretenir}

\begin{savaistu}[Pourquoi les bijections sont partout]
Les bijections sont l'outil préféré des mathématiciens pour dire que deux ensembles « se ressemblent ». Elles sont aussi au cœur de la vie quotidienne numérique : quand tu paies par Mobile Money (MTN MoMo ou Orange Money), ton numéro de téléphone est en bijection avec ton compte — c'est ce qui garantit que l'argent arrive chez le bon destinataire, et qu'on peut toujours « remonter » d'une transaction à son émetteur. En cryptographie, chiffrer un message, c'est appliquer une bijection ; déchiffrer, c'est appliquer sa réciproque !
\end{savaistu}

\begin{exercice}[À toi de jouer]
\begin{enumerate}
\item Montrer que l'application $f : \mathbb{R} \to \mathbb{R}$ définie par $f(x) = -2x + 5$ est bijective et expliciter $f^{-1}$.
\item L'application $g : \mathbb{R} \to \mathbb{R}$, $g(x) = x^2 + 1$ est-elle injective ? Surjective ? Justifier.
\item Déterminer une restriction de $g$ (ensemble de départ et ensemble d'arrivée) qui la rende bijective, et donner alors sa réciproque.
\end{enumerate}
\end{exercice}

\begin{corrige}[Exercice n°1]
\textbf{1.} Soit $y \in \mathbb{R}$. $f(x) = y \iff -2x + 5 = y \iff -2x = y - 5 \iff x = \dfrac{5 - y}{2}$. Pour tout $y \in \mathbb{R}$, cette solution existe, est unique et est réelle : $f$ est bijective et $f^{-1}(x) = \dfrac{5 - x}{2}$.

\textbf{2.} $g(-1) = 2 = g(1)$ avec $-1 \neq 1$ : $g$ n'est pas injective. L'équation $x^2 + 1 = 0$ n'a pas de solution réelle, donc $0$ n'a pas d'antécédent : $g$ n'est pas surjective de $\mathbb{R}$ dans $\mathbb{R}$.

\textbf{3.} Comme $x^2 \geq 0$, on a $g(x) \geq 1$ pour tout $x$. Prenons $h : [0\,;\,+\infty[ \to [1\,;\,+\infty[$, $h(x) = x^2 + 1$. Pour $y \geq 1$, $x^2 + 1 = y \iff x^2 = y - 1 \iff x = \sqrt{y-1}$ (on garde la solution positive, car $x \geq 0$). Solution unique dans $[0\,;\,+\infty[$ : $h$ est bijective et $h^{-1}(x) = \sqrt{x - 1}$.
\end{corrige}

\coursec{Composition des applications}

Après avoir étudié les propriétés internes d'une application (injectivité, surjectivité, bijectivité), nous nous intéressons maintenant à une opération fondamentale qui permet de \emph{chaîner} deux applications : la composition. C'est l'outil mathématique qui modélise l'enchaînement de deux processus : « d'abord faire $f$, puis faire $g$ sur le résultat ». On la retrouve partout : en algorithmique (fonctions imbriquées), en physique (changement de référentiel), en économie (coût de production puis prix de vente), ou simplement quand on calcule $\sqrt{3x+1}$ : on commence par multiplier par 3 et ajouter 1, puis on prend la racine carrée.

\courssub{Définition et notation de la composée}

Imagine deux machines $f$ et $g$. La machine $f$ prend une entrée $x$ dans $E$ et produit une sortie $f(x)$ dans $F$. La machine $g$ prend une entrée dans $F$ et produit une sortie dans $G$. Si on connecte la sortie de $f$ à l'entrée de $g$, on obtient une nouvelle machine qui envoie directement $x$ (de $E$) vers $g(f(x))$ (dans $G$). C'est la \textbf{composée de $g$ par $f$}.

\begin{definition}[Composée de deux applications]
Soient $E$, $F$, $G$ trois ensembles, $f : E \to F$ et $g : F \to G$ deux applications.
On appelle \textbf{composée de $g$ par $f$}, et on note $g \circ f$, l'application de $E$ dans $G$ qui associe à tout $x \in E$ l'élément $g(f(x)) \in G$ :
\[
g \circ f : E \longrightarrow G, \quad x \longmapsto g(f(x)).
\]
On lit « $g$ rond $f$ » ou « $g$ composé $f$ ».
\end{definition}

\noindent\textbf{Condition d'existence :} Pour que $g(f(x))$ ait un sens, il faut que l'image $f(x)$ appartienne à l'ensemble de définition de $g$ (noté $D_g$). Comme $f$ est une application de $E$ vers $F$, on a $f(x) \in F$. Si $g$ est définie sur tout $F$ (c'est-à-dire $D_g = F$), la composition est toujours possible. Si $g$ n'est définie que sur une partie de $F$ ($D_g \subsetneq F$), il faut vérifier que $\text{Im}(f) \subset D_g$.

\begin{popfigure}
\begin{tikzpicture}[
    node distance=2.5cm and 3cm,
    set/.style={draw=popdark, thick, ellipse, minimum width=2.5cm, minimum height=3cm},
    elem/.style={circle, fill=popblue, inner sep=2pt},
    arr/.style={-Stealth, thick, popdark},
    maparr/.style={-Stealth, thick, poporange, bend left=15},
    maparrrev/.style={-Stealth, thick, poppurple, bend right=15}
]
\node[set] (E) {$E$};
\node[set, right=of E] (F) {$F$};
\node[set, right=of F] (G) {$G$};

% Elements in E
\foreach \y/\name in {1/x_1, 0.5/x_2, 0/x_3, -0.5/x_4, -1/x_5}
    \node[elem, label=left:$\name$] (e\name) at ($(E.west)+(0,\y)$) {};

% Elements in F
\foreach \y/\name in {1/y_1, 0/y_2, -1/y_3}
    \node[elem, label=right:$\name$] (f\name) at ($(F.east)+(0,\y)$) {};

% Elements in G
\foreach \y/\name in {1/z_1, -1/z_2}
    \node[elem, label=right:$\name$] (g\name) at ($(G.east)+(0,\y)$) {};

% Arrows f: E -> F
\draw[arr] (ex_1) -- node[above, sloped, midway] {$f$} (fy_1);
\draw[arr] (ex_2) -- (fy_1);
\draw[arr] (ex_3) -- (fy_2);
\draw[arr] (ex_4) -- (fy_2);
\draw[arr] (ex_5) -- (fy_3);

% Arrows g: F -> G
\draw[arr] (fy_1) -- node[above, sloped, midway] {$g$} (gz_1);
\draw[arr] (fy_2) -- (gz_1);
\draw[arr] (fy_3) -- (gz_2);

% Composed arrow g o f: E -> G (dashed, different color)
\draw[popgreen, dashed, thick] (ex_1) to[bend left=30] node[above, sloped, midway] {$g \circ f$} (gz_1);
\draw[popgreen, dashed, thick] (ex_2) to[bend left=30] (gz_1);
\draw[popgreen, dashed, thick] (ex_3) to[bend left=30] (gz_1);
\draw[popgreen, dashed, thick] (ex_4) to[bend left=30] (gz_1);
\draw[popgreen, dashed, thick] (ex_5) to[bend right=30] (gz_2);

% Labels for functions
\node[above=0.3cm of E.north, popblue, font=\bfseries] {$f$};
\node[above=0.3cm of F.north, poporange, font=\bfseries] {$g$};
\node[above=0.3cm of G.north, popgreen, font=\bfseries] {$g \circ f$};
\end{tikzpicture}
\legende{Schéma sagittal (diagramme fléché) illustrant la composition $g \circ f$. Les flèches pleines représentent $f$ et $g$, les flèches vertes en pointillés représentent la composée $g \circ f : x \mapsto g(f(x))$.}
\end{popfigure}

\courssub{Ensemble de définition d'une composée de fonctions numériques}

En Première, nous manipulons des \emph{fonctions numériques} : $f : D_f \to \mathbb{R}$ et $g : D_g \to \mathbb{R}$. La composée $g \circ f$ n'est pas définie sur tout $D_f$ automatiquement : il faut que $f(x)$ « tombe » dans $D_g$.

\begin{propriete}[Ensemble de définition de $g \circ f$]
Soient $f$ et $g$ deux fonctions numériques d'ensembles de définition respectifs $D_f$ et $D_g$.
L'ensemble de définition de la composée $g \circ f$ est :
\[
D_{g \circ f} = \{ x \in D_f \mid f(x) \in D_g \}.
\]
En mots : on prend les $x$ qui appartiennent à $D_f$ \textbf{et} dont l'image par $f$ appartient à $D_g$.
\end{propriete}

\begin{methode}[Déterminer $D_{g \circ f}$ en deux étapes]
\begin{enumerate}
    \item \textbf{Condition 1 :} $x \in D_f$ (l'expression $f(x)$ existe).
    \item \textbf{Condition 2 :} $f(x) \in D_g$ (l'expression $g(f(x))$ existe).
    \item \textbf{Conclusion :} $D_{g \circ f}$ est l'intersection des solutions des deux conditions.
\end{enumerate}
\end{methode}

\begin{exoresolu}[Ensemble de définition d'une composée]
Soient $f(x) = \sqrt{x-2}$ et $g(x) = \dfrac{1}{x-3}$.
Déterminer $D_{g \circ f}$ et l'expression de $(g \circ f)(x)$.
\tcblower
\textbf{Solution.}
\underline{Étape 1 : Ensembles de définition de base.}
\begin{itemize}
    \item $f(x) = \sqrt{x-2}$ est définie si $x-2 \ge 0 \iff x \ge 2$. Donc $D_f = [2, +\infty[$.
    \item $g(x) = \dfrac{1}{x-3}$ est définie si $x-3 \neq 0 \iff x \neq 3$. Donc $D_g = \mathbb{R} \setminus \{3\}$.
\end{itemize}

\underline{Étape 2 : Condition pour la composée.}
On cherche $x \in D_f$ tel que $f(x) \in D_g$.
$x \in D_f \iff x \ge 2$.
$f(x) \in D_g \iff f(x) \neq 3 \iff \sqrt{x-2} \neq 3 \iff x-2 \neq 9 \iff x \neq 11$.

\underline{Étape 3 : Intersection.}
Il faut $x \ge 2$ \textbf{et} $x \neq 11$.
Donc $D_{g \circ f} = [2, 11[ \cup ]11, +\infty[$.

\underline{Étape 4 : Expression de la composée.}
Pour $x \in D_{g \circ f}$ :
\[
(g \circ f)(x) = g(f(x)) = \frac{1}{f(x) - 3} = \frac{1}{\sqrt{x-2} - 3}.
\]
\end{exoresolu}

\courssub{Non-commutativité et associativité de la composition}

C'est un point crucial : \textbf{l'ordre compte}. $g \circ f$ signifie « $f$ \emph{puis} $g$ ». En général, $g \circ f \neq f \circ g$.

\begin{attention}[Piège majeur : $g \circ f \neq f \circ g$]
Ne confonds jamais l'ordre ! $g \circ f$ s'applique en faisant \textbf{d'abord $f$}, ensuite $g$.
Un moyen mnémotechnique : dans $g \circ f$, le $f$ est « collé » au $x$ : $(g \circ f)(x) = g(f(x))$. C'est $f$ qui touche $x$ en premier.
\end{attention}

\begin{exemplebox}[Contre-exemple chiffré : $f(x)=x^2$, $g(x)=x-3$]
Soient $f(x)=x^2$ ($D_f=\mathbb{R}$) et $g(x)=x-3$ ($D_g=\mathbb{R}$).
\begin{align*}
(g \circ f)(x) &= g(f(x)) = g(x^2) = x^2 - 3. \\
(f \circ g)(x) &= f(g(x)) = f(x-3) = (x-3)^2 = x^2 - 6x + 9.
\end{align*}
Pour $x=0$ : $(g \circ f)(0) = -3$ alors que $(f \circ g)(0) = 9$.
Les deux fonctions sont \textbf{différentes}. La composition n'est \textbf{pas commutative}.
\end{exemplebox}

\begin{propriete}[Associativité de la composition (admise)]
Soient $f : E \to F$, $g : F \to G$, $h : G \to H$ trois applications telles que les compositions aient un sens.
Alors : $h \circ (g \circ f) = (h \circ g) \circ f$.
On note simplement $h \circ g \circ f$ sans parenthèses. Cela signifie qu'on applique $f$, puis $g$, puis $h$, peu importe comment on groupe les opérations.
\end{propriete}

\begin{exemplebox}[Vérification de l'associativité]
Soient $f(x)=x+1$, $g(x)=2x$, $h(x)=x^2$ (toutes définies sur $\mathbb{R}$).
\begin{itemize}
    \item $(g \circ f)(x) = 2(x+1) = 2x+2$. Donc $(h \circ (g \circ f))(x) = (2x+2)^2 = 4x^2 + 8x + 4$.
    \item $(h \circ g)(x) = (2x)^2 = 4x^2$. Donc $((h \circ g) \circ f)(x) = 4(x+1)^2 = 4(x^2+2x+1) = 4x^2+8x+4$.
\end{itemize}
On retrouve bien la même fonction.
\end{exemplebox}

\courssub{Composée de deux bijections et bijection réciproque}

C'est un résultat fondamental, souvent demandé en démonstration au Probatoire. Si $f$ et $g$ sont bijectives, leur composée l'est aussi, et la réciproque de la composée se calcule en \textbf{inversant l'ordre} des réciproques.

\begin{experience}[Démonstration guidée : $(g \circ f)^{-1} = f^{-1} \circ g^{-1}$]
Soient $E, F, G$ trois ensembles et $f : E \to F$, $g : F \to G$ deux \textbf{bijections}.
On veut prouver que $g \circ f$ est bijective et que $(g \circ f)^{-1} = f^{-1} \circ g^{-1}$.

\textbf{1. Montrer que $g \circ f$ est bijective.}
\begin{itemize}
    \item \textit{Injectivité :} Soient $x_1, x_2 \in E$ tels que $(g \circ f)(x_1) = (g \circ f)(x_2)$.
    Alors $g(f(x_1)) = g(f(x_2))$. Comme $g$ est injective, $f(x_1) = f(x_2)$.
    Comme $f$ est injective, $x_1 = x_2$. Donc $g \circ f$ est injective.
    \item \textit{Surjectivité :} Soit $z \in G$. Comme $g$ est surjective, $\exists y \in F$ tel que $g(y) = z$.
    Comme $f$ est surjective, $\exists x \in E$ tel que $f(x) = y$.
    Alors $(g \circ f)(x) = g(f(x)) = g(y) = z$. Donc $g \circ f$ est surjective.
\end{itemize}
$g \circ f$ est donc bijective de $E$ vers $G$. Sa réciproque $(g \circ f)^{-1}$ va de $G$ vers $E$.

\textbf{2. Identifier $(g \circ f)^{-1}$.}
On sait que $f^{-1} : F \to E$ et $g^{-1} : G \to F$ existent.
La composée $f^{-1} \circ g^{-1}$ va de $G$ vers $E$ (ordre inversé !).
Vérifions qu'elle est bien la réciproque de $g \circ f$ en composant dans les deux sens :
\begin{align*}
(g \circ f) \circ (f^{-1} \circ g^{-1}) &= g \circ (f \circ f^{-1}) \circ g^{-1} &&\text{(associativité)} \\
&= g \circ \text{Id}_F \circ g^{-1} &&\text{($f \circ f^{-1} = \text{Id}_F$)} \\
&= g \circ g^{-1} = \text{Id}_G. \\
(f^{-1} \circ g^{-1}) \circ (g \circ f) &= f^{-1} \circ (g^{-1} \circ g) \circ f &&\text{(associativité)} \\
&= f^{-1} \circ \text{Id}_F \circ f = f^{-1} \circ f = \text{Id}_E.
\end{align*}
Par définition de la bijection réciproque, on a bien \textbf{$(g \circ f)^{-1} = f^{-1} \circ g^{-1}$}.
\end{experience}

\begin{aretenir}[Règle pratique : « Chaussettes-Chaussures »]
Pour enlever ses chaussures et ses chaussettes, on enlève \textbf{d'abord les chaussures} (la dernière chose mise), \textbf{puis les chaussettes}.
Pour inverser une composée $g \circ f$ (fait $f$ puis $g$), on inverse l'ordre : on annule $g$ \textbf{d'abord} ($g^{-1}$), puis on annule $f$ ($f^{-1}$).
\[
(g \circ f)^{-1} = f^{-1} \circ g^{-1}
\]
\end{aretenir}

\begin{exoresolu}[Bijection réciproque d'une composée]
Soit $h(x) = 2\sqrt{x-1} + 3$ définie sur $D_h = [1, +\infty[$.
On peut écrire $h = g \circ f$ avec $f(x) = \sqrt{x-1}$ et $g(x) = 2x + 3$.
Montrer que $h$ est bijective de $[1, +\infty[$ vers $[3, +\infty[$ et déterminer $h^{-1}$.
\tcblower
\textbf{Solution.}
\underline{1. Vérification des bijections.}
\begin{itemize}
    \item $f : [1, +\infty[ \to [0, +\infty[$, $x \mapsto \sqrt{x-1}$.
    Strictement croissante (racine carrée et $x-1$ croissantes), $f(1)=0$, limite $+\infty$ en $+\infty$. \textbf{Bijection}.
    \item $g : [0, +\infty[ \to [3, +\infty[$, $x \mapsto 2x+3$.
    Affine, coefficient directeur $2>0$, $g(0)=3$. \textbf{Bijection}.
\end{itemize}
$h = g \circ f$ est donc bijective de $[1, +\infty[$ vers $[3, +\infty[$ (composition de deux bijections).

\underline{2. Réciproques des fonctions simples.}
\begin{itemize}
    \item $y = \sqrt{x-1} \iff y^2 = x-1 \ (y \ge 0) \iff x = y^2+1$. Donc $f^{-1}(y) = y^2+1$ (de $[0, +\infty[$ vers $[1, +\infty[$).
    \item $y = 2x+3 \iff x = \frac{y-3}{2}$. Donc $g^{-1}(y) = \frac{y-3}{2}$ (de $[3, +\infty[$ vers $[0, +\infty[$).
\end{itemize}

\underline{3. Réciproque de $h$ (inversion de l'ordre !).}
$h^{-1} = f^{-1} \circ g^{-1}$.
Pour $y \in [3, +\infty[$ :
\[
h^{-1}(y) = f^{-1}(g^{-1}(y)) = f^{-1}\left(\frac{y-3}{2}\right) = \left(\frac{y-3}{2}\right)^2 + 1 = \frac{(y-3)^2}{4} + 1.
\]
En changeant la variable en $x$ : $\boxed{h^{-1}(x) = \frac{(x-3)^2}{4} + 1}$, définie sur $[3, +\infty[$.
\end{exoresolu}

\courssub{Décomposition d'une fonction en composée de fonctions usuelles}

C'est l'exercice inverse : on te donne une fonction « complexe » $h$, et tu dois la « casser » en une suite de fonctions simples $f_1, f_2, \dots, f_n$ (affines, racine, carré, inverse, valeur absolue, partie entière, puissance entière $x \mapsto x^n$) telles que $h = f_n \circ \dots \circ f_2 \circ f_1$.
La méthode : \textbf{lis l'expression de la dernière opération effectuée}.

\begin{methode}[Décomposer $h$ en composée de fonctions usuelles]
\begin{enumerate}
    \item Identifie la \textbf{dernière opération} effectuée sur $x$ pour obtenir $h(x)$. C'est la fonction « extérieure » $g$ (la dernière appliquée).
    \item Ce qu'il y a « à l'intérieur » de cette dernière opération est l'image par la fonction « intérieure » $f$ (la première appliquée).
    \item Si l'intérieur est encore complexe, répète l'opération sur $f$.
    \item Vérifie : $h(x) = g(f(x))$ (ou $g(f_2(f_1(x)))$ etc.).
\end{enumerate}
\end{methode}

\begin{popfigure}
\begin{tikzpicture}[
    node distance=1.5cm and 2cm,
    box/.style={draw=popdark, thick, rounded corners, minimum width=2.2cm, minimum height=1.2cm, fill=popcream, align=center},
    arr/.style={-Stealth, very thick, popblue},
    labelstyle/.style={font=\small, text=popdark}
]
\node[box] (start) {$x$};
\node[box, right=of start] (f1) {$f_1 : x \mapsto 3x$};
\node[box, right=of f1] (f2) {$f_2 : u \mapsto u+1$};
\node[box, right=of f2] (f3) {$f_3 : v \mapsto \sqrt{v}$};
\node[box, right=of f3] (end) {$h(x) = \sqrt{3x+1}$};

\draw[arr] (start) -- node[above, labelstyle] {$x$} (f1);
\draw[arr] (f1) -- node[above, labelstyle] {$3x$} (f2);
\draw[arr] (f2) -- node[above, labelstyle] {$3x+1$} (f3);
\draw[arr] (f3) -- node[above, labelstyle] {$\sqrt{3x+1}$} (end);

% Annotations below using simple nodes instead of braces (library decorations.pathreplacing not guaranteed)
\node[below=1.0cm of start, poporange, font=\bfseries\small] {Étape 1};
\node[below=1.0cm of f1, poporange, font=\bfseries\small] {$\times 3$};
\node[below=1.0cm of f2, poporange, font=\bfseries\small] {Étape 2};
\node[below=1.0cm of f2, poporange, font=\bfseries\small, yshift=-0.5cm] {$+1$};
\node[below=1.0cm of f3, poporange, font=\bfseries\small] {Étape 3};
\node[below=1.0cm of f3, poporange, font=\bfseries\small, yshift=-0.5cm] {$\sqrt{\cdot}$ (dernière)};

\node[below=2.2cm of f2, poppurple, font=\bfseries\itshape] {Lecture de droite à gauche pour décomposer};
\end{tikzpicture}
\legende{Décomposition en chaîne de $h(x) = \sqrt{3x+1}$. La dernière opération est la racine carrée ($f_3$), avant cela l'addition de 1 ($f_2$), au début la multiplication par 3 ($f_1$). On a $h = f_3 \circ f_2 \circ f_1$.}
\end{popfigure}

\begin{exoresolu}[Décomposition de fonctions]
Décomposer les fonctions suivantes en composée de fonctions usuelles (affine, $x \mapsto x^2$, $x \mapsto \sqrt{x}$, $x \mapsto 1/x$, $x \mapsto |x|$, $x \mapsto E(x)$, $x \mapsto x^n$).
\begin{enumerate}
    \item $h_1(x) = (2x-5)^3$
    \item $h_2(x) = \dfrac{1}{\sqrt{x+4}}$
    \item $h_3(x) = |x^2 - 4|$
    \item $h_4(x) = E(\sqrt{x})$
\end{enumerate}
\tcblower
\textbf{Solution.}
On identifie la \textbf{dernière opération} (la plus « extérieure »).

\begin{enumerate}
    \item $h_1(x) = (2x-5)^3$. Dernière op. : élever au cube ($\varphi(u)=u^3$, fonction puissance usuelle). Intérieur : $2x-5$ (affine $f(x)=2x-5$).
    $h_1 = \varphi \circ f$ avec $f(x)=2x-5$, $\varphi(u)=u^3$.
    \item $h_2(x) = \dfrac{1}{\sqrt{x+4}}$. Dernière op. : inverse ($g(v)=1/v$). Intérieur : $\sqrt{x+4}$.
    Décomposition de l'intérieur : dernière op. racine ($f_2(u)=\sqrt{u}$), intérieur $x+4$ (affine $f_1(x)=x+4$).
    $h_2 = g \circ f_2 \circ f_1$ avec $f_1(x)=x+4$, $f_2(u)=\sqrt{u}$, $g(v)=1/v$.
    \item $h_3(x) = |x^2 - 4|$. Dernière op. : valeur absolue ($g(v)=|v|$). Intérieur : $x^2-4$.
    Décomposition intérieure : $x^2$ (carré $f_1(x)=x^2$) puis $-4$ (affine $f_2(u)=u-4$).
    Ordre : $x \xrightarrow{f_1} x^2 \xrightarrow{f_2} x^2-4 \xrightarrow{g} |x^2-4|$.
    $h_3 = g \circ f_2 \circ f_1$ avec $f_1(x)=x^2$, $f_2(u)=u-4$, $g(v)=|v|$.
    \item $h_4(x) = E(\sqrt{x})$. Dernière op. : partie entière ($g(v)=E(v)$). Intérieur : $\sqrt{x}$ ($f(x)=\sqrt{x}$).
    $h_4 = g \circ f$ avec $f(x)=\sqrt{x}$, $g(u)=E(u)$.
\end{enumerate}
\end{exoresolu}

\begin{savaistu}[Composition et modélisation : Le prix du taxi à Yaoundé]
Imaginons un trajet en taxi. Le prix dépend de la distance $d$ (en km).
\begin{itemize}
    \item Fonction $f$ : distance $\to$ temps de trajet. $f(d) = \frac{d}{v}$ ($v$ vitesse moyenne).
    \item Fonction $g$ : temps $\to$ coût. $g(t) = a + b \times t$ (forfait + horaire).
\end{itemize}
Le coût total en fonction de la distance est $C(d) = (g \circ f)(d) = a + b \cdot \frac{d}{v}$.
Si la vitesse change (embouteillages), $f$ change, donc $g \circ f$ change. C'est la puissance de la composition : modéliser des systèmes complexes en briques simples.
\end{savaistu}

\begin{exercice}[Entraînement : Composition et ensemble de définition]
Soient $f(x) = \sqrt{4-x}$ et $g(x) = \ln(x)$ (logarithme népérien).
\begin{enumerate}
    \item Déterminer $D_f$ et $D_g$.
    \item Déterminer $D_{g \circ f}$ et $D_{f \circ g}$.
    \item Calculer les expressions de $(g \circ f)(x)$ et $(f \circ g)(x)$ sur leurs ensembles de définition respectifs.
    \item $g \circ f$ et $f \circ g$ sont-elles égales ? Justifier.
\end{enumerate}
\end{exercice}

\begin{corrige}[Exercice n]
\begin{enumerate}
    \item $f(x)=\sqrt{4-x}$ définie si $4-x \ge 0 \iff x \le 4$. $D_f = ]-\infty, 4]$.
    $g(x)=\ln(x)$ définie si $x > 0$. $D_g = ]0, +\infty[$.
    \item $D_{g \circ f} = \{x \in D_f \mid f(x) \in D_g\} = \{x \le 4 \mid \sqrt{4-x} > 0\}$.
    $\sqrt{4-x} > 0 \iff 4-x > 0 \iff x < 4$.
    Donc $D_{g \circ f} = ]-\infty, 4[$.
    $D_{f \circ g} = \{x \in D_g \mid g(x) \in D_f\} = \{x > 0 \mid \ln(x) \le 4\}$.
    $\ln(x) \le 4 \iff x \le e^4$.
    Donc $D_{f \circ g} = ]0, e^4]$.
    \item $(g \circ f)(x) = \ln(\sqrt{4-x}) = \frac{1}{2}\ln(4-x)$ pour $x < 4$.
    $(f \circ g)(x) = \sqrt{4 - \ln(x)}$ pour $0 < x \le e^4$.
    \item Elles n'ont même pas le même ensemble de définition ($]-\infty, 4[$ vs $]0, e^4]$), donc elles sont \textbf{différentes}. La composition n'est pas commutative.
\end{enumerate}
\end{corrige}

\coursec{Opérations sur les polynômes et les fonctions}

Dans la section précédente, nous avons appris à \emph{composer} des applications, c'est-à-dire à les enchaîner : appliquer $f$, puis appliquer $g$ au résultat. Mais il existe une façon encore plus naturelle de combiner des fonctions : les \emph{additionner}, les \emph{multiplier}, les \emph{diviser}. Quand un commerçant de Mokolo calcule son bénéfice, il soustrait ses coûts de ses recettes ; quand Orange Cameroun facture un forfait, elle additionne un abonnement fixe et un coût variable. Derrière chacune de ces situations se cache une opération sur des fonctions. Cette section te donne les outils rigoureux pour effectuer ces opérations et, surtout, pour déterminer leurs \cle{ensembles de définition}, compétence incontournable au Probatoire.

\courssub{Rappels sur les polynômes}

\begin{definition}[Polynôme]
On appelle \cle{polynôme} (ou fonction polynôme) toute fonction $P$ définie sur $\mathbb{R}$ par une expression de la forme
\[ P(x) = a_n x^n + a_{n-1}x^{n-1} + \cdots + a_1 x + a_0, \]
où $n \in \mathbb{N}$ et les nombres $a_0, a_1, \ldots, a_n$ sont des réels appelés les \cle{coefficients} du polynôme.
\begin{itemize}
\item Si $a_n \neq 0$, on dit que $P$ est de \cle{degré} $n$, et on note $\deg P = n$. Le coefficient $a_n$ est appelé \emph{coefficient dominant}.
\item Un polynôme de degré $0$ est une constante non nulle.
\item Le \cle{polynôme nul} est celui dont tous les coefficients sont nuls ; par convention, il n'a pas de degré.
\end{itemize}
\end{definition}

\begin{exemplebox}[Reconnaître un polynôme et son degré]
\begin{itemize}
\item $P(x) = 2x^2 - 3x + 1$ est un polynôme de degré $2$ (trinôme du second degré).
\item $Q(x) = x + 4$ est un polynôme de degré $1$ (binôme du premier degré).
\item $R(x) = 7$ est un polynôme de degré $0$.
\item $S(x) = \sqrt{x} + 1$ n'est \textbf{pas} un polynôme, car $\sqrt{x}$ n'est pas une puissance entière de $x$.
\end{itemize}
\end{exemplebox}

\courssub{Somme et produit de deux polynômes}

L'idée est très simple : pour additionner deux polynômes, on regroupe les termes de même puissance de $x$ ; pour les multiplier, on développe et on réduit. Formalisons cela.

\begin{definition}[Somme et produit de deux polynômes]
Soient $P(x) = a_n x^n + \cdots + a_0$ et $Q(x) = b_m x^m + \cdots + b_0$ deux polynômes.
\begin{itemize}
\item Leur \cle{somme} est le polynôme $(P+Q)(x) = P(x) + Q(x)$, obtenu en additionnant les coefficients des termes de même degré.
\item Leur \cle{produit} est le polynôme $(P \times Q)(x) = P(x) \times Q(x)$, obtenu en distribuant chaque terme de $P$ sur chaque terme de $Q$, puis en réduisant.
\end{itemize}
\end{definition}

\begin{propriete}[Degrés de la somme et du produit]
Soient $P$ et $Q$ deux polynômes non nuls. Alors :
\[ \deg(P+Q) \leq \max(\deg P, \deg Q) \qquad \text{et} \qquad \deg(P \times Q) = \deg P + \deg Q. \]
\end{propriete}

\begin{experience}[Pourquoi le degré du produit est la somme des degrés]
Raisonnons sur les termes dominants. Écrivons $P(x) = a_n x^n + (\text{termes de degré} < n)$ avec $a_n \neq 0$, et $Q(x) = b_m x^m + (\text{termes de degré} < m)$ avec $b_m \neq 0$.
\begin{enumerate}
\item Quand on développe le produit $P \times Q$, le terme de plus haut degré provient du produit des deux termes dominants : $a_n x^n \times b_m x^m = a_n b_m\, x^{n+m}$.
\item Tous les autres produits de termes donnent des puissances de $x$ \emph{strictement inférieures} à $n+m$.
\item Comme $a_n \neq 0$ et $b_m \neq 0$, on a $a_n b_m \neq 0$ : le coefficient de $x^{n+m}$ ne peut pas disparaître.
\end{enumerate}
Conclusion : le terme dominant de $P \times Q$ est $a_n b_m x^{n+m}$, donc $\deg(P\times Q) = n + m = \deg P + \deg Q$. \hfill $\square$
\end{experience}

\begin{attention}[L'inégalité pour la somme n'est pas toujours une égalité]
Contrairement au produit, le degré de la somme peut \emph{chuter} ! Si les termes dominants s'annulent mutuellement, le degré diminue. Par exemple, avec $P(x) = 3x^2 + x$ et $Q(x) = -3x^2 + 5$ :
\[ (P+Q)(x) = x + 5, \]
qui est de degré $1$, alors que $\max(\deg P, \deg Q) = 2$. Retiens : $\deg(P+Q) \leq \max(\deg P, \deg Q)$, avec égalité dès que $\deg P \neq \deg Q$.
\end{attention}

\begin{exemplebox}[Calculs avec $P(x) = 2x^2 - 3x + 1$ et $Q(x) = x + 4$]
\textbf{Somme.} On regroupe les termes de même degré :
\begin{align*}
(P+Q)(x) &= (2x^2 - 3x + 1) + (x + 4) \\
&= 2x^2 + (-3+1)x + (1+4) \\
&= 2x^2 - 2x + 5.
\end{align*}
On vérifie : $\deg(P+Q) = 2 = \max(2, 1)$.

\textbf{Produit.} On distribue puis on réduit :
\begin{align*}
(P \times Q)(x) &= (2x^2 - 3x + 1)(x + 4) \\
&= 2x^3 + 8x^2 - 3x^2 - 12x + x + 4 \\
&= 2x^3 + 5x^2 - 11x + 4.
\end{align*}
On vérifie : $\deg(P \times Q) = 3 = 2 + 1$. Le terme dominant $2x^3$ est bien le produit des termes dominants $2x^2 \times x$.
\end{exemplebox}

\courssub{Quotient de deux polynômes : fractions rationnelles}

Diviser est plus délicat qu'additionner ou multiplier : le quotient de deux polynômes n'est en général \emph{pas} un polynôme. On obtient un nouvel objet, la \cle{fraction rationnelle}, et la question centrale devient : \emph{pour quelles valeurs de $x$ cette expression a-t-elle un sens ?}

\begin{definition}[Fraction rationnelle]
Soient $P$ et $Q$ deux polynômes, avec $Q$ \textbf{non nul}. On appelle \cle{fraction rationnelle} associée au quotient $\dfrac{P}{Q}$ la fonction définie par
\[ x \longmapsto \frac{P(x)}{Q(x)}. \]
\end{definition}

\begin{propriete}[Condition d'existence du quotient $P/Q$]
Le quotient $\dfrac{P(x)}{Q(x)}$ existe si et seulement si le dénominateur ne s'annule pas, c'est-à-dire si $Q(x) \neq 0$. L'ensemble de définition de la fraction rationnelle est donc
\[ D = \left\{ x \in \mathbb{R} \;\middle|\; Q(x) \neq 0 \right\} = \mathbb{R} \setminus \{\text{racines de } Q\}. \]
Les racines de $Q$ sont appelées les \cle{valeurs interdites}.
\end{propriete}

\begin{exemplebox}[Domaine du quotient $\dfrac{P}{Q}$ avec $P(x) = 2x^2 - 3x + 1$ et $Q(x) = x + 4$]
Le quotient $\dfrac{P(x)}{Q(x)} = \dfrac{2x^2 - 3x + 1}{x + 4}$ existe si et seulement si $x + 4 \neq 0$, c'est-à-dire $x \neq -4$. Donc
\[ D_{P/Q} = \mathbb{R} \setminus \{-4\} = \left]-\infty\,; -4\right[ \cup \left]-4\,; +\infty\right[. \]
La valeur $-4$ est interdite : en $x = -4$, le dénominateur vaut $0$ alors que le numérateur vaut $2(-4)^2 - 3(-4) + 1 = 32 + 12 + 1 = 45 \neq 0$ ; le quotient n'a aucun sens.
\end{exemplebox}

\begin{attention}[Simplifier avant de chercher le domaine : une erreur fatale]
Considère la fonction $h(x) = \dfrac{x^2 - 1}{x - 1}$. On est tenté de factoriser : $x^2 - 1 = (x-1)(x+1)$, puis de \og simplifier \fg{} : $h(x) = x + 1$. Erreur ! Cette simplification n'est valable que pour $x \neq 1$. La fonction $h$ \textbf{reste non définie en $1$}, même après simplification :
\[ D_h = \mathbb{R} \setminus \{1\}. \]
Règle d'or : on détermine \textbf{toujours} l'ensemble de définition \textbf{avant} toute simplification. Les valeurs interdites restent interdites, même si elles \og disparaissent \fg{} de l'écriture.
\end{attention}

\courssub{Somme, produit et quotient de deux fonctions numériques}

Tout ce que nous venons de faire avec les polynômes se généralise à des fonctions numériques quelconques. Le principe est \emph{ponctuel} : on combine les images en chaque point $x$.

\begin{definition}[Opérations ponctuelles sur les fonctions]
Soient $f$ et $g$ deux fonctions numériques définies respectivement sur $D_f$ et $D_g$.
\begin{itemize}
\item La \cle{somme} $f + g$ est définie par $(f+g)(x) = f(x) + g(x)$.
\item Le \cle{produit} $f \cdot g$ est définie par $(f \cdot g)(x) = f(x) \times g(x)$.
\item Le \cle{quotient} $\dfrac{f}{g}$ est défini par $\left(\dfrac{f}{g}\right)(x) = \dfrac{f(x)}{g(x)}$, là où $g(x) \neq 0$.
\end{itemize}
\end{definition}

\begin{propriete}[Ensembles de définition des opérations]
Avec les notations précédentes :
\[ D_{f+g} = D_{f \cdot g} = D_f \cap D_g \qquad \text{et} \qquad D_{f/g} = \left\{ x \in D_f \cap D_g \;\middle|\; g(x) \neq 0 \right\}. \]
\end{propriete}

L'intuition est limpide : pour calculer $f(x) + g(x)$, il faut que \emph{les deux} nombres $f(x)$ et $g(x)$ existent, d'où l'intersection. Pour le quotient, on ajoute la condition vitale : le dénominateur ne doit pas s'annuler.

\begin{methode}[Déterminer l'ensemble de définition d'un quotient de fonctions]
Pour déterminer $D_{f/g}$, procède toujours dans cet ordre :
\begin{enumerate}
\item Détermine séparément $D_f$ et $D_g$ (conditions d'existence de chaque fonction : dénominateurs non nuls, quantités sous les racines positives ou nulles, etc.).
\item Forme l'intersection $D_f \cap D_g$.
\item Résous l'équation $g(x) = 0$ et \textbf{exclus} toutes ses solutions de l'intersection.
\item Conclus en écrivant $D_{f/g}$ comme une réunion d'intervalles, \emph{avant} toute simplification éventuelle de l'expression.
\end{enumerate}
\end{methode}

\begin{exoresolu}[Trois ensembles de définition complets]
Soient $f(x) = \dfrac{1}{x-2}$ et $g(x) = \sqrt{x+1}$. Déterminer les ensembles de définition de $f + g$, de $f \cdot g$ et de $\dfrac{f}{g}$.
\tcblower
\textbf{Étape 1 : domaines de $f$ et $g$.}
\begin{itemize}
\item $f$ existe ssi $x - 2 \neq 0$, donc $D_f = \mathbb{R} \setminus \{2\}$.
\item $g$ existe ssi $x + 1 \geq 0$, c'est-à-dire $x \geq -1$, donc $D_g = [-1\,; +\infty[$.
\end{itemize}
\textbf{Étape 2 : intersection.}
\[ D_f \cap D_g = [-1\,; +\infty[ \setminus \{2\} = [-1\,; 2[ \cup ]2\,; +\infty[. \]
Donc $D_{f+g} = D_{f \cdot g} = [-1\,; 2[ \cup ]2\,; +\infty[$.

\textbf{Étape 3 : zéros de $g$.} $g(x) = 0 \iff \sqrt{x+1} = 0 \iff x = -1$. Or $-1 \in D_f \cap D_g$ : on doit l'exclure pour le quotient.

\textbf{Conclusion :}
\[ D_{f/g} = \left]-1\,; 2\right[ \cup \left]2\,; +\infty\right[. \]
Remarque bien la différence : $-1$ est autorisé pour $f+g$ et $f \cdot g$, mais interdit pour $f/g$ car il annule le dénominateur.
\end{exoresolu}

\begin{center}
\begin{tabularx}{\linewidth}{|l|X|X|}
\hline
\rowcolor{popblueL} \textbf{Opération} & \textbf{Expression} & \textbf{Ensemble de définition} \\
\hline
Somme & $(f+g)(x) = f(x)+g(x)$ & $D_f \cap D_g$ \\
\hline
Produit & $(f\cdot g)(x) = f(x)\,g(x)$ & $D_f \cap D_g$ \\
\hline
Quotient & $\left(\dfrac{f}{g}\right)(x) = \dfrac{f(x)}{g(x)}$ & $\{x \in D_f \cap D_g \mid g(x) \neq 0\}$ \\
\hline
Composée (rappel) & $(g \circ f)(x) = g(f(x))$ & $\{x \in D_f \mid f(x) \in D_g\}$ \\
\hline
\end{tabularx}
\end{center}

\courssub{Lecture graphique : construire $f + g$ point par point}

Comment visualiser la somme de deux fonctions ? Pour chaque abscisse $x$, l'ordonnée de $f+g$ est la \emph{somme des ordonnées} de $f$ et de $g$. Graphiquement, on \og empile \fg{} les hauteurs algébriques (en tenant compte des signes).

\begin{popfigure}
\begin{center}
\begin{tikzpicture}
\begin{axis}[
  width=0.9\linewidth, height=7cm,
  axis lines=middle, xlabel=$x$, ylabel=$y$,
  xmin=-3.2, xmax=3.2, ymin=-4.5, ymax=6.5,
  xtick={-3,-2,-1,1,2,3}, ytick={-4,-2,2,4,6},
  grid=both, grid style={popline!40},
  legend style={at={(0.02,0.98)}, anchor=north west, font=\small}
]
\addplot[popcurve, popblue, domain=-3:3, samples=200]{x + 1};
\addlegendentry{$g(x)=x+1$}
\addplot[popcurve, poporange, domain=-3:3, samples=200]{0.5*x^2 - 2};
\addlegendentry{$f(x)=0{,}5x^2-2$}
\addplot[popcurve, poppurple, very thick, domain=-3:3, samples=200]{0.5*x^2 + x - 1};
\addlegendentry{$(f+g)(x)=0{,}5x^2+x-1$}
\draw[popgreen, dashed] (axis cs:2,0) -- (axis cs:2,3);
\fill[poporange] (axis cs:2,0) circle (2pt);
\fill[popblue] (axis cs:2,3) circle (2pt);
\fill[poppurple] (axis cs:2,3) circle (2.6pt);
\node[popdark, anchor=west, font=\small] at (axis cs:2.1,1.5) {$f(2)+g(2)=0+3=3$};
\end{axis}
\end{tikzpicture}
\end{center}
\legende{En chaque abscisse $x$, l'ordonnée de $f+g$ (courbe violette) est la somme des ordonnées de $f$ (orange) et de $g$ (bleue). Exemple en $x=2$ : $f(2)=0$ et $g(2)=3$, donc $(f+g)(2)=3$ : les points bleu et violet se confondent.}
\end{popfigure}

\begin{attention}[Le domaine de $f+g$ n'est pas une réunion]
Une erreur classique consiste à écrire $D_{f+g} = D_f \cup D_g$. Faux ! Si $x$ n'appartient qu'à l'un des deux domaines, l'une des deux images n'existe pas et la somme n'a pas de sens. C'est bien l'\textbf{intersection} qui intervient : les deux fonctions doivent être définies \emph{simultanément}.
\end{attention}

\begin{savaistu}[Des polynômes dans ton téléphone]
Les fractions rationnelles — quotients de polynômes — ne sont pas qu'un exercice scolaire. En traitement du signal, les filtres qui nettoient le son de tes appels MTN ou Orange sont décrits par des quotients de polynômes (leur \emph{fonction de transfert}). Les racines du dénominateur, appelées \emph{pôles}, déterminent les fréquences que le filtre amplifie ou rejette : exactement les \og valeurs interdites \fg{} que tu viens d'apprendre à traquer !
\end{savaistu}

\begin{exercice}[À toi de jouer]
Soient $P(x) = x^2 - 5x + 6$ et $Q(x) = x^2 - 4$.
\begin{enumerate}
\item Calculer $(P+Q)(x)$ et $(P \times Q)(x)$, en précisant leurs degrés.
\item Déterminer l'ensemble de définition de la fraction rationnelle $R(x) = \dfrac{P(x)}{Q(x)}$.
\item Factoriser $P(x)$ et $Q(x)$, simplifier $R(x)$, et vérifier que les valeurs interdites n'ont pas changé.
\end{enumerate}
\end{exercice}

\begin{corrige}[Exercice : opérations sur $P$ et $Q$]
\textbf{1.} Somme : $(P+Q)(x) = 2x^2 - 5x + 2$, de degré $2 = \max(2,2)$ (les termes dominants ne s'annulent pas). Produit :
\begin{align*}
(P \times Q)(x) &= (x^2 - 5x + 6)(x^2 - 4) \\
&= x^4 - 5x^3 + 6x^2 - 4x^2 + 20x - 24 \\
&= x^4 - 5x^3 + 2x^2 + 20x - 24,
\end{align*}
de degré $4 = 2 + 2$, comme prévu.

\textbf{2.} $Q(x) = 0 \iff x^2 = 4 \iff x = 2$ ou $x = -2$. Donc
\[ D_R = \mathbb{R} \setminus \{-2\,; 2\}. \]

\textbf{3.} $P(x) = (x-2)(x-3)$ et $Q(x) = (x-2)(x+2)$. Pour $x \neq 2$ et $x \neq -2$ :
\[ R(x) = \frac{(x-2)(x-3)}{(x-2)(x+2)} = \frac{x-3}{x+2}. \]
La simplification ne change \textbf{rien} au domaine : $x = 2$ reste une valeur interdite, même si le facteur $(x-2)$ a disparu de l'écriture. On a bien $D_R = \mathbb{R} \setminus \{-2\,; 2\}$.
\end{corrige}

En résumé, additionner, multiplier ou diviser des fonctions est une opération \emph{ponctuelle} simple ; toute la rigueur réside dans la détermination des ensembles de définition : intersection des domaines pour la somme et le produit, intersection privée des zéros du dénominateur pour le quotient, et jamais de simplification avant l'étude du domaine. Ces réflexes te serviront dans toute la suite du cours, notamment lors de l'étude des fonctions associées et des courbes.

\coursec{Parité, périodicité et éléments de symétrie}

Dans la section précédente, nous avons appris à construire de nouvelles fonctions par somme, produit, quotient et composition, et à déterminer leurs ensembles de définition. Nous allons maintenant étudier les \cle{propriétés géométriques} que certaines fonctions possèdent : la \cle{parité}, la \cle{périodicité} et les \cle{éléments de symétrie} de leurs courbes. L'idée est simple et puissante : si une courbe possède une symétrie, il suffit de l'étudier sur la moitié (voire le quart) de son ensemble de définition, le reste s'en déduit gratuitement. C'est un gain de temps considérable, au Probatoire comme dans la vie du mathématicien.

\courssub{Fonctions paires}

\begin{definition}[Fonction paire]
Soit $f$ une fonction numérique d'ensemble de définition $D_f$. On dit que $f$ est \cle{paire} lorsque les deux conditions suivantes sont satisfaites :
\begin{enumerate}
\item $D_f$ est \textbf{symétrique par rapport à $0$} : pour tout $x \in D_f$, on a $-x \in D_f$ ;
\item pour tout $x \in D_f$, $f(-x) = f(x)$.
\end{enumerate}
\end{definition}

L'intuition est la suivante : une fonction paire « répond de la même façon » à une valeur et à son opposée. La fonction $f(x) = x^2$ en est le modèle : $(-3)^2 = 3^2 = 9$.

\begin{propriete}[Conséquence graphique de la parité]
Si $f$ est paire, alors sa courbe représentative $\mathcal{C}_f$ admet l'axe des ordonnées $(Oy)$ comme \cle{axe de symétrie}. En effet, les points $M(x\,;\,f(x))$ et $M'(-x\,;\,f(-x)) = M'(-x\,;\,f(x))$ sont symétriques par rapport à $(Oy)$.
\end{propriete}

\begin{exemplebox}[Exemples de fonctions paires]
\begin{itemize}
\item $f(x) = x^2$ sur $\mathbb{R}$ : $f(-x) = (-x)^2 = x^2 = f(x)$, donc $f$ est paire.
\item $g(x) = |x|$ sur $\mathbb{R}$ : $g(-x) = |-x| = |x| = g(x)$, donc $g$ est paire.
\item $h(x) = \dfrac{1}{x^2}$ sur $\mathbb{R}^*$ : $\mathbb{R}^*$ est symétrique par rapport à $0$ et $h(-x) = \dfrac{1}{(-x)^2} = \dfrac{1}{x^2} = h(x)$, donc $h$ est paire.
\end{itemize}
\end{exemplebox}

\courssub{Fonctions impaires}

\begin{definition}[Fonction impaire]
Soit $f$ une fonction numérique d'ensemble de définition $D_f$. On dit que $f$ est \cle{impaire} lorsque :
\begin{enumerate}
\item $D_f$ est symétrique par rapport à $0$ ;
\item pour tout $x \in D_f$, $f(-x) = -f(x)$.
\end{enumerate}
\end{definition}

\begin{propriete}[Conséquence graphique de l'imparité]
Si $f$ est impaire, alors sa courbe $\mathcal{C}_f$ admet l'origine $O$ du repère comme \cle{centre de symétrie}. En effet, les points $M(x\,;\,f(x))$ et $M'(-x\,;\,-f(x))$ sont symétriques par rapport à $O$.
\end{propriete}

\begin{exemplebox}[Exemples de fonctions impaires]
\begin{itemize}
\item $f(x) = x^3$ sur $\mathbb{R}$ : $f(-x) = (-x)^3 = -x^3 = -f(x)$, donc $f$ est impaire.
\item $g(x) = \dfrac{1}{x}$ sur $\mathbb{R}^*$ : $g(-x) = -\dfrac{1}{x} = -g(x)$, donc $g$ est impaire.
\end{itemize}
\end{exemplebox}

\begin{popfigure}
\begin{tikzpicture}
\begin{axis}[width=0.48\linewidth, axis lines=middle, xlabel={$x$}, ylabel={$y$}, xmin=-3, xmax=3, ymin=-1, ymax=5, xtick={-2,-1,1,2}, ytick={1,2,3,4}, tick label style={font=\scriptsize}, title={$f(x)=x^2$ : paire}, title style={font=\small}]
\addplot[popcurve,domain=-2.1:2.1,samples=200]{x^2};
\addplot[poporange,dashed,domain=-1:4.5,samples=2] coordinates {(0,-1) (0,4.5)};
\addplot[only marks,mark=*,popblue] coordinates {(1.5,2.25) (-1.5,2.25)};
\node[popblue,font=\scriptsize] at (axis cs:1.7,2.9) {$M$};
\node[popblue,font=\scriptsize] at (axis cs:-1.7,2.9) {$M'$};
\end{axis}
\end{tikzpicture}\hfill
\begin{tikzpicture}
\begin{axis}[width=0.48\linewidth, axis lines=middle, xlabel={$x$}, ylabel={$y$}, xmin=-2.5, xmax=2.5, ymin=-4, ymax=4, xtick={-2,-1,1,2}, ytick={-3,-2,-1,1,2,3}, tick label style={font=\scriptsize}, title={$g(x)=x^3$ : impaire}, title style={font=\small}]
\addplot[popcurve,domain=-1.55:1.55,samples=200]{x^3};
\addplot[only marks,mark=*,popgreen] coordinates {(1.2,1.728) (-1.2,-1.728)};
\addplot[only marks,mark=*,poporange] coordinates {(0,0)};
\node[popgreen,font=\scriptsize] at (axis cs:1.35,2.5) {$M$};
\node[popgreen,font=\scriptsize] at (axis cs:-1.4,-2.6) {$M'$};
\node[poporange,font=\scriptsize] at (axis cs:0.3,-0.6) {$O$};
\draw[popmuted,dashed] (axis cs:-1.2,-1.728) -- (axis cs:1.2,1.728);
\end{axis}
\end{tikzpicture}
\legende{À gauche, une fonction paire : la courbe est symétrique par rapport à $(Oy)$. À droite, une fonction impaire : la courbe est symétrique par rapport à l'origine $O$.}
\end{popfigure}

\courssub{Méthode d'étude de la parité}

\begin{methode}[Étudier la parité d'une fonction en trois étapes]
\begin{enumerate}
\item \textbf{Symétrie du domaine.} Déterminer $D_f$ et vérifier qu'il est symétrique par rapport à $0$ : si $x \in D_f$ alors $-x \in D_f$. Si ce n'est pas le cas, conclure immédiatement : $f$ n'est ni paire ni impaire.
\item \textbf{Calcul de $f(-x)$.} Pour $x \in D_f$, calculer $f(-x)$ et le simplifier.
\item \textbf{Comparaison et conclusion.} Si $f(-x) = f(x)$ pour tout $x$, alors $f$ est paire. Si $f(-x) = -f(x)$ pour tout $x$, alors $f$ est impaire. Sinon, $f$ n'est ni paire ni impaire.
\end{enumerate}
\end{methode}

\begin{attention}[Le piège classique du domaine]
Ne conclus \textbf{jamais} sur la parité sans avoir vérifié d'abord que $D_f$ est symétrique par rapport à $0$. Contre-exemple : soit $f(x) = x^2$ définie sur $D_f = [0\,;\,+\infty[$. On a bien $(-x)^2 = x^2$ au niveau du calcul, mais $-x$ n'appartient pas à $D_f$ : la fonction n'est \emph{ni paire ni impaire}. De même, une fonction définie sur $[-2\,;\,3]$ ne peut être ni paire ni impaire, car $3 \in D_f$ mais $-3 \notin D_f$.
\end{attention}

\begin{exoresolu}[Étude de parité]
Étudier la parité des fonctions suivantes :
\[
f(x) = \frac{x}{x^2 + 1} \qquad ; \qquad g(x) = x^2 + x \qquad ; \qquad h(x) = \sqrt{x^2 - 4}.
\]
\tcblower
\textbf{Fonction $f$.} $x^2 + 1 \geq 1 > 0$ pour tout réel $x$, donc $D_f = \mathbb{R}$, symétrique par rapport à $0$. Pour tout $x$ :
\[
f(-x) = \frac{-x}{(-x)^2 + 1} = \frac{-x}{x^2 + 1} = -f(x).
\]
Donc $f$ est \textbf{impaire} : sa courbe admet $O$ comme centre de symétrie.

\textbf{Fonction $g$.} $D_g = \mathbb{R}$, symétrique par rapport à $0$. On calcule :
\[
g(-x) = (-x)^2 + (-x) = x^2 - x.
\]
Or $g(-x) \neq g(x)$ (par exemple $g(1) = 2$ et $g(-1) = 0$) et $g(-x) \neq -g(x)$ (car $-g(1) = -2 \neq 0 = g(-1)$). Donc $g$ n'est \textbf{ni paire ni impaire}.

\textbf{Fonction $h$.} $h(x)$ existe si et seulement si $x^2 - 4 \geq 0$, c'est-à-dire $x \leq -2$ ou $x \geq 2$. Donc $D_h = \left]-\infty\,;\,-2\right] \cup \left[2\,;\,+\infty\right[$, qui est symétrique par rapport à $0$. De plus :
\[
h(-x) = \sqrt{(-x)^2 - 4} = \sqrt{x^2 - 4} = h(x).
\]
Donc $h$ est \textbf{paire}.
\end{exoresolu}

\courssub{Fonctions périodiques}

Certaines situations se répètent à intervalles réguliers : les marées, les saisons agricoles, le passage du car Express Union devant le lycée toutes les 30 minutes, la tension du courant alternatif distribué par ENEO (50 Hz, soit une période de 0,02 s). Les mathématiques modélisent cela par les fonctions périodiques.

\begin{definition}[Fonction périodique]
Soit $f$ une fonction d'ensemble de définition $D_f$ et $T$ un réel \textbf{strictement positif}. On dit que $f$ est \cle{périodique} de période $T$ lorsque pour tout $x \in D_f$ :
\[
x + T \in D_f \qquad \text{et} \qquad f(x + T) = f(x).
\]
La \textbf{période} de $f$ est, lorsqu'elle existe, la \cle{plus petite} valeur strictement positive $T$ vérifiant cette propriété.
\end{definition}

\begin{aretenir}[Conséquence graphique de la périodicité]
Si $f$ est périodique de période $T$, sa courbe est \textbf{invariante par la translation de vecteur} $T\vec{\imath}$ : il suffit de tracer la courbe sur un intervalle de longueur $T$ (une « bande » de largeur $T$), puis de la reproduire par translations successives. Par exemple, la fonction sinus est périodique de période $2\pi$ : $\sin(x + 2\pi) = \sin x$ pour tout réel $x$.
\end{aretenir}

\begin{savaistu}[Périodicité et signal électrique]
Au Cameroun, le courant alternatif distribué par ENEO a une fréquence de 50 Hz. La tension aux bornes d'une prise suit une loi sinusoïdale de période $T = \frac{1}{50} = 0,02$ s. C'est grâce à la théorie de Fourier (décomposition de signaux périodiques en somme de sinusoïdes) que l'on peut analyser et transmettre les signaux radio (MTN, Orange, CRTV) ou compresser des images (JPEG) et du son (MP3).
\end{savaistu}

\begin{popfigure}
\begin{tikzpicture}
\begin{axis}[width=0.9\linewidth, height=5cm, axis lines=middle, xlabel={$x$}, ylabel={$y$}, xmin=-0.5, xmax=6.8, ymin=-1.6, ymax=1.6, xtick={1.57,3.14,4.71,6.28}, xticklabels={$\frac{\pi}{2}$,$\pi$,$\frac{3\pi}{2}$,$2\pi$}, ytick={-1,1}, tick label style={font=\scriptsize}]
\addplot[popcurve,domain=0:6.5,samples=300]{sin(deg(x))};
\draw[poporange,thick,<->] (axis cs:0,1.25) -- (axis cs:6.283,1.25) node[midway,above,font=\small]{$T = 2\pi$};
\addplot[only marks,mark=*,popgreen] coordinates {(0.9,0.783) (7.183-6.283,0.783)};
\end{axis}
\end{tikzpicture}
\legende{La fonction sinus est périodique de période $T = 2\pi$ : le motif se répète identiquement à lui-même sur chaque intervalle de longueur $2\pi$.}
\end{popfigure}

\begin{exemplebox}[Démontrer qu'une fonction est périodique]
Soit $f$ la fonction définie sur $\mathbb{R}$ par $f(x) = x - E(x)$, où $E(x)$ désigne la partie entière de $x$ (le plus grand entier inférieur ou égal à $x$). Montrons que $f$ est périodique de période $1$. Pour tout réel $x$, $E(x + 1) = E(x) + 1$, donc :
\[
f(x + 1) = (x + 1) - E(x + 1) = x + 1 - E(x) - 1 = x - E(x) = f(x).
\]
Ainsi $f$ est périodique de période $1$ (c'est la fonction « en dents de scie »). On vérifie qu'aucun réel $T \in \left]0\,;\,1\right[$ ne convient, donc $1$ est bien la plus petite période.
\end{exemplebox}

\courssub{Centre de symétrie d'une courbe}

Généralisons l'imparité : une courbe peut admettre un centre de symétrie autre que l'origine. Par exemple, l'hyperbole d'équation $y = 1 + \dfrac{2}{x-1}$ « tourne autour » du point $\Omega(1\,;\,1)$.

\begin{definition}[Centre de symétrie]
Soit $f$ une fonction de courbe $\mathcal{C}_f$ et $\Omega(a\,;\,b)$ un point. On dit que $\Omega$ est \cle{centre de symétrie} de $\mathcal{C}_f$ lorsque pour tout réel $h$ tel que $a + h \in D_f$, on a $a - h \in D_f$ et :
\[
f(a + h) + f(a - h) = 2b.
\]
\end{definition}

\begin{experience}[Démonstration : pourquoi $f(a+h) + f(a-h) = 2b$ ?]
Cette démonstration est formatrice ; elle repose sur les coordonnées du milieu d'un segment.
\begin{enumerate}
\item Soit $h$ un réel tel que $a + h$ et $a - h$ appartiennent à $D_f$. Considérons les deux points de la courbe :
\[
M\bigl(a + h\,;\,f(a+h)\bigr) \qquad \text{et} \qquad M'\bigl(a - h\,;\,f(a-h)\bigr).
\]
\item Dire que $\Omega(a\,;\,b)$ est le milieu du segment $[MM']$, c'est dire que :
\[
\frac{(a+h) + (a-h)}{2} = a \qquad \text{et} \qquad \frac{f(a+h) + f(a-h)}{2} = b.
\]
\item La première égalité est toujours vraie. La seconde s'écrit exactement :
\[
f(a+h) + f(a-h) = 2b.
\]
\item Ainsi, $\Omega$ est centre de symétrie de $\mathcal{C}_f$ si et seulement si, pour tout $h$ admissible, $f(a+h) + f(a-h) = 2b$. \hfill$\blacksquare$
\end{enumerate}
\end{experience}

\begin{exemplebox}[Exemple chiffré : un centre de symétrie]
Montrons que la courbe de $f(x) = 1 + \dfrac{2}{x - 1}$ admet le point $\Omega(1\,;\,1)$ comme centre de symétrie.

$D_f = \mathbb{R} \setminus \{1\}$ : si $1 + h \in D_f$ (c'est-à-dire $h \neq 0$), alors $1 - h \in D_f$. Pour tout $h \neq 0$ :
\begin{align*}
f(1 + h) + f(1 - h) &= \left(1 + \frac{2}{h}\right) + \left(1 + \frac{2}{-h}\right) \\
&= 1 + \frac{2}{h} + 1 - \frac{2}{h} = 2 = 2 \times 1.
\end{align*}
La condition $f(a+h) + f(a-h) = 2b$ est satisfaite avec $a = 1$ et $b = 1$ : $\Omega(1\,;\,1)$ est bien centre de symétrie de $\mathcal{C}_f$.
\end{exemplebox}

\begin{popfigure}
\begin{tikzpicture}
\begin{axis}[width=0.75\linewidth, height=7cm, axis lines=middle, xlabel={$x$}, ylabel={$y$}, xmin=-3.5, xmax=5.5, ymin=-2.5, ymax=4.5, xtick={1}, ytick={1}, tick label style={font=\scriptsize}]
\addplot[popcurve,domain=1.45:5.3,samples=200]{1 + 2/(x-1)};
\addplot[popcurve,domain=-3.3:0.55,samples=200]{1 + 2/(x-1)};
\addplot[only marks,mark=*,poporange] coordinates {(1,1)};
\node[poporange,font=\small] at (axis cs:1.25,0.55) {$\Omega(1;1)$};
\addplot[only marks,mark=*,popblue] coordinates {(3,2) (-1,0)};
\node[popblue,font=\small] at (axis cs:3.15,2.3) {$M(3;2)$};
\node[popblue,font=\small] at (axis cs:-1.35,-0.35) {$M'(-1;0)$};
\draw[popmuted,dashed] (axis cs:-1,0) -- (axis cs:3,2);
\draw[popgreen,dashed] (axis cs:1,-2.5) -- (axis cs:1,4.5);
\draw[popgreen,dashed] (axis cs:-3.5,1) -- (axis cs:5.5,1);
\end{axis}
\end{tikzpicture}
\legende{Hyperbole d'équation $y = 1 + \dfrac{2}{x-1}$ : le point $\Omega(1\,;\,1)$ est centre de symétrie ; $M$ et $M'$ sont symétriques par rapport à $\Omega$.}
\end{popfigure}

\courssub{Axe de symétrie d'une courbe}

De même, une courbe peut admettre un axe de symétrie vertical autre que $(Oy)$ : la parabole d'équation $y = (x-2)^2$ est symétrique par rapport à la droite d'équation $x = 2$.

\begin{definition}[Axe de symétrie vertical]
Soit $f$ une fonction de courbe $\mathcal{C}_f$ et $a$ un réel. La droite d'équation $x = a$ est \cle{axe de symétrie} de $\mathcal{C}_f$ lorsque pour tout réel $h$ tel que $a + h \in D_f$, on a $a - h \in D_f$ et :
\[
f(a + h) = f(a - h).
\]
\end{definition}

\begin{experience}[Démonstration guidée : la condition $f(a+h) = f(a-h)$]
Montrons que la droite $\Delta$ d'équation $x = a$ est axe de symétrie de $\mathcal{C}_f$ si et seulement si $f(a+h) = f(a-h)$ pour tout $h$ admissible.
\begin{enumerate}
\item Rappel : le symétrique du point $M(x_0\,;\,y_0)$ par rapport à la droite verticale d'équation $x = a$ est le point $M'(2a - x_0\,;\,y_0)$ (même ordonnée, abscisse « reflétée »).
\item Prenons un point de la courbe d'abscisse $a + h$ : $M\bigl(a+h\,;\,f(a+h)\bigr)$.
\item Son symétrique par rapport à $\Delta$ est $M'\bigl(2a - (a+h)\,;\,f(a+h)\bigr) = M'\bigl(a - h\,;\,f(a+h)\bigr)$.
\item Dire que $\Delta$ est axe de symétrie, c'est dire que $M'$ appartient aussi à $\mathcal{C}_f$, c'est-à-dire que son ordonnée vaut $f(a - h)$ :
\[
f(a - h) = f(a + h).
\]
\item Réciproquement, si cette égalité vaut pour tout $h$ admissible, le symétrique de tout point de la courbe est sur la courbe : $\Delta$ est axe de symétrie. \hfill$\blacksquare$
\end{enumerate}
\end{experience}

\begin{exoresolu}[Justifier un axe de symétrie]
Soit $f$ la fonction définie sur $\mathbb{R}$ par $f(x) = x^2 - 4x + 5$. Montrer que la droite d'équation $x = 2$ est axe de symétrie de $\mathcal{C}_f$.
\tcblower
$D_f = \mathbb{R}$ : pour tout réel $h$, $2 + h$ et $2 - h$ appartiennent à $D_f$. Calculons :
\begin{align*}
f(2 + h) &= (2+h)^2 - 4(2+h) + 5 = 4 + 4h + h^2 - 8 - 4h + 5 = h^2 + 1, \\
f(2 - h) &= (2-h)^2 - 4(2-h) + 5 = 4 - 4h + h^2 - 8 + 4h + 5 = h^2 + 1.
\end{align*}
On a donc $f(2+h) = f(2-h)$ pour tout réel $h$ : la droite d'équation $x = 2$ est bien axe de symétrie de $\mathcal{C}_f$. (On retrouve le sommet de la parabole en $S(2\,;\,1)$.)
\end{exoresolu}

\courssub{Appartenance d'un point à une courbe}

\begin{methode}[Montrer qu'un point appartient à une courbe]
Pour justifier que le point $A(x_0\,;\,y_0)$ appartient à la courbe $\mathcal{C}_f$, il suffit de :
\begin{enumerate}
\item vérifier que $x_0 \in D_f$ ;
\item calculer $f(x_0)$ et vérifier que $f(x_0) = y_0$.
\end{enumerate}
Par exemple, pour $f(x) = 1 + \dfrac{2}{x-1}$ définie sur $\mathbb{R}\setminus\{1\}$ :
\begin{itemize}
\item Le point $A(3\,;\,2)$ appartient à $\mathcal{C}_f$ car $3 \in D_f$ et $f(3) = 1 + \dfrac{2}{2} = 2$.
\item Le point $B(2\,;\,3)$ appartient aussi à $\mathcal{C}_f$ car $2 \in D_f$ et $f(2) = 1 + \dfrac{2}{2-1} = 3$.
\end{itemize}
Vérifie toujours par le calcul, jamais à l'œil sur un tracé approximatif.
\end{methode}

\begin{attention}[Ce que le programme attend de toi]
La \textbf{recherche systématique} des éléments de symétrie d'une courbe est \emph{hors programme}. Au Probatoire, l'énoncé te \emph{donnera} le point $\Omega(a\,;\,b)$ ou la droite $x = a$, et te demandera de \emph{justifier} qu'il s'agit bien d'un centre ou d'un axe de symétrie, en utilisant les conditions $f(a+h) + f(a-h) = 2b$ ou $f(a+h) = f(a-h)$. À l'inverse, la lecture graphique (conjecturer une symétrie, une période sur une courbe donnée) est un savoir-faire exigible.
\end{attention}

\begin{exercice}[Pour s'entraîner]
\begin{enumerate}
\item Étudier la parité de $f(x) = \dfrac{x^3 - x}{x^2 + 4}$ et de $g(x) = \dfrac{x + 1}{x - 1}$.
\item Soit $f(x) = \dfrac{x^2 - 2x + 3}{x - 1}$ définie sur $\mathbb{R} \setminus \{1\}$. Montrer que le point $\Omega(1\,;\,0)$ est centre de symétrie de $\mathcal{C}_f$.
\item Montrer que la droite d'équation $x = -1$ est axe de symétrie de la courbe de $h(x) = x^2 + 2x - 3$.
\end{enumerate}
\end{exercice}

\begin{corrige}[Exercice n°1]
\textbf{1.} $D_f = \mathbb{R}$ (car $x^2 + 4 > 0$), symétrique par rapport à $0$, et $f(-x) = \dfrac{-x^3 + x}{x^2 + 4} = -f(x)$ : $f$ est impaire. Pour $g$ : $D_g = \mathbb{R} \setminus \{1\}$ n'est pas symétrique par rapport à $0$ ($-1 \in D_g$ mais $1 \notin D_g$), donc $g$ n'est ni paire ni impaire.

\textbf{2.} Pour $h \neq 0$ : $f(1+h) = \dfrac{(1+h)^2 - 2(1+h) + 3}{h} = \dfrac{h^2 + 2}{h} = h + \dfrac{2}{h}$ et $f(1-h) = \dfrac{h^2 + 2}{-h} = -h - \dfrac{2}{h}$. Donc $f(1+h) + f(1-h) = 0 = 2 \times 0$ : $\Omega(1\,;\,0)$ est centre de symétrie.

\textbf{3.} $h(-1+t) = (-1+t)^2 + 2(-1+t) - 3 = t^2 - 4$ et $h(-1-t) = t^2 - 4$. L'égalité $h(-1+t) = h(-1-t)$ pour tout $t$ prouve que la droite $x = -1$ est axe de symétrie.
\end{corrige}

\coursec{Courbe d'une fonction et fonctions associées}

Dans la section précédente, nous avons appris à repérer, \emph{par le calcul}, la parité, la périodicité et les éléments de symétrie d'une fonction. Toutes ces propriétés ont un point commun : elles se \emph{voient} sur un dessin. Une fonction paire donne une courbe symétrique par rapport à l'axe des ordonnées ; une fonction périodique donne un motif qui se répète. Il est donc temps de donner à nos fonctions un visage : leur \cle{courbe représentative}. C'est ce visage qui permet, d'un seul coup d'œil, de conjecturer le domaine, les variations, les asymptotes — avant de tout démontrer rigoureusement par le calcul.

\courssub{Courbe représentative d'une fonction}

\textbf{Intuition.} Quand un médecin regarde un électrocardiogramme, il lit une courbe : à chaque instant $t$ correspond une tension électrique $f(t)$. La courbe, c'est simplement la trace laissée par tous les points de coordonnées $(t\,;\,f(t))$. En mathématiques, l'idée est identique.

\begin{definition}[Courbe représentative]
Soit $f$ une fonction numérique définie sur un ensemble $D_f$, et $(O\,;\vec{\imath},\vec{\jmath})$ un repère du plan. La \cle{courbe représentative} de $f$, notée $\mathcal{C}_f$, est l'ensemble des points du plan dont l'abscisse $x$ parcourt $D_f$ et dont l'ordonnée est l'image $f(x)$ :
\[ \mathcal{C}_f = \left\{ M\big(x\,;\,f(x)\big) \;\middle|\; x \in D_f \right\}. \]
On dit que la courbe $\mathcal{C}_f$ a pour \cle{équation} $y = f(x)$ dans ce repère.
\end{definition}

\begin{attention}[Une courbe de fonction ne « revient » jamais en arrière]
Une fonction associe à chaque $x \in D_f$ \textbf{une seule} image. Par conséquent, toute droite verticale d'équation $x = a$ (avec $a \in D_f$) coupe $\mathcal{C}_f$ en \textbf{au plus un point}. Un cercle complet, par exemple, n'est pas la courbe d'une fonction : pour une même abscisse, il y aurait deux ordonnées.
\end{attention}

\courssub{Appartenance d'un point à une courbe}

Puisque $\mathcal{C}_f$ est l'ensemble des points $M(x\,;\,f(x))$, dire qu'un point est sur la courbe revient à vérifier une égalité.

\begin{propriete}[Condition d'appartenance]
Soit $f$ une fonction de courbe $\mathcal{C}_f$ et $M(x_0\,;\,y_0)$ un point du plan. Alors :
\[ M(x_0\,;\,y_0) \in \mathcal{C}_f \iff x_0 \in D_f \ \text{et}\ y_0 = f(x_0). \]
\end{propriete}

\begin{methode}[Justifier qu'un point appartient (ou non) à une courbe]
\begin{enumerate}
\item Vérifier que l'abscisse $x_0$ appartient à l'ensemble de définition $D_f$.
\item Calculer $f(x_0)$.
\item Comparer avec $y_0$ : si $f(x_0) = y_0$, le point est sur la courbe ; sinon, il n'y est pas.
\end{enumerate}
\end{methode}

\begin{exoresolu}[Points sur la courbe d'une fonction homographique]
Soit $f$ la fonction définie par $f(x) = \dfrac{2x+1}{x-1}$ et $\mathcal{C}_f$ sa courbe. Les points $A(2\,;\,5)$, $B(0\,;\,1)$ et $C(1\,;\,3)$ appartiennent-ils à $\mathcal{C}_f$ ?
\tcblower
\textbf{Solution.} $f$ est définie pour $x - 1 \neq 0$, donc $D_f = \mathbb{R}\setminus\{1\}$.
\begin{itemize}
\item Pour $A(2\,;\,5)$ : $2 \in D_f$ et $f(2) = \dfrac{2\times 2 + 1}{2-1} = \dfrac{5}{1} = 5 = y_A$. Donc $A \in \mathcal{C}_f$.
\item Pour $B(0\,;\,1)$ : $0 \in D_f$ et $f(0) = \dfrac{1}{-1} = -1 \neq 1$. Donc $B \notin \mathcal{C}_f$.
\item Pour $C(1\,;\,3)$ : $1 \notin D_f$, donc $C \notin \mathcal{C}_f$ (inutile de calculer !).
\end{itemize}
\end{exoresolu}

\courssub{Intersection de la courbe avec les axes de coordonnées}

Deux questions naturelles se posent lorsqu'on trace une courbe : où coupe-t-elle l'axe des ordonnées ? Où coupe-t-elle l'axe des abscisses ?

\begin{propriete}[Intersections avec les axes]
Soit $f$ une fonction de courbe $\mathcal{C}_f$.
\begin{itemize}
\item \textbf{Axe des ordonnées} : si $0 \in D_f$, alors $\mathcal{C}_f$ coupe l'axe $(Oy)$ en l'\textbf{unique} point $\big(0\,;\,f(0)\big)$. Si $0 \notin D_f$, il n'y a pas d'intersection.
\item \textbf{Axe des abscisses} : les points d'intersection de $\mathcal{C}_f$ avec $(Ox)$ ont pour abscisses les solutions de l'équation $f(x) = 0$ (les \cle{zéros} de $f$). Il peut y en avoir zéro, un, ou plusieurs.
\end{itemize}
\end{propriete}

\begin{exemplebox}[Avec un trinôme]
Soit $f(x) = x^2 - x - 6$, définie sur $\mathbb{R}$.
\begin{itemize}
\item Intersection avec $(Oy)$ : $f(0) = -6$, donc le point $(0\,;\,-6)$.
\item Intersection avec $(Ox)$ : on résout $x^2 - x - 6 = 0$. Le discriminant est $\Delta = (-1)^2 - 4\times 1 \times (-6) = 1 + 24 = 25$, d'où $x = \dfrac{1-5}{2} = -2$ ou $x = \dfrac{1+5}{2} = 3$. La courbe coupe $(Ox)$ en $(-2\,;\,0)$ et $(3\,;\,0)$.
\end{itemize}
\end{exemplebox}

\begin{attention}[Asymétrie des deux situations]
Avec $(Oy)$, il y a \textbf{au plus un} point d'intersection (car une fonction a une seule image par $x$). Avec $(Ox)$, il peut y en avoir \textbf{plusieurs} (une équation peut avoir plusieurs solutions). Ne cherche donc jamais « le » point d'intersection avec $(Ox)$ avant d'avoir résolu l'équation.
\end{attention}

\courssub{Fonctions associées à une fonction donnée}

\textbf{Intuition.} Le réseau MTN affiche une tarification « de base » $f(x)$ en FCFA pour $x$ minutes d'appel. Lors d'une promotion, l'opérateur peut : ajouter un forfait fixe de $200$ FCFA ($f(x) + 200$), décaler le début de la facturation ($f(x-1)$), ou appliquer un avoir ($-f(x)$). À chaque fois, la \emph{nouvelle} tarification se déduit de l'ancienne par une transformation simple. Mathématiquement, ces nouvelles fonctions sont dites \cle{associées} à $f$, et leurs courbes se déduisent de $\mathcal{C}_f$ par des déplacements ou des symétries.

\begin{definition}[Fonctions associées]
Soit $f$ une fonction définie sur $D_f$, $k$ et $a$ deux réels. Les fonctions suivantes sont dites \cle{associées} à $f$ :
\begin{itemize}
\item $g : x \mapsto f(x) + k$, définie sur $D_f$ ;
\item $h : x \mapsto f(x - a)$, définie sur $\{x \in \mathbb{R} \mid x - a \in D_f\}$ ;
\item $p : x \mapsto -f(x)$, définie sur $D_f$ ;
\item $q : x \mapsto f(-x)$, définie sur $\{x \in \mathbb{R} \mid -x \in D_f\}$ ;
\item $r : x \mapsto |f(x)|$, définie sur $D_f$.
\end{itemize}
\end{definition}

\begin{propriete}[Effets géométriques sur la courbe — admise]
Soit $\mathcal{C}_f$ la courbe de $f$ dans un repère $(O\,;\vec{\imath},\vec{\jmath})$. On obtient les courbes des fonctions associées ainsi :
\begin{itemize}
\item $y = f(x) + k$ : \textbf{translation} de $\mathcal{C}_f$ de vecteur $k\,\vec{\jmath}$ (verticale, vers le haut si $k>0$) ;
\item $y = f(x - a)$ : \textbf{translation} de $\mathcal{C}_f$ de vecteur $a\,\vec{\imath}$ (horizontale, vers la droite si $a>0$) ;
\item $y = -f(x)$ : \textbf{symétrie} de $\mathcal{C}_f$ par rapport à l'axe $(Ox)$ ;
\item $y = f(-x)$ : \textbf{symétrie} de $\mathcal{C}_f$ par rapport à l'axe $(Oy)$ ;
\item $y = |f(x)|$ : on conserve les parties de $\mathcal{C}_f$ situées \textbf{au-dessus} de $(Ox)$, et on \textbf{réfléchit} (symétrie par rapport à $(Ox)$) les parties situées en dessous.
\end{itemize}
\end{propriete}

\begin{experience}[Pourquoi $f(x-a)$ translate-t-il vers la droite ?]
Prenons un point $M\big(x_0\,;\,f(x_0)\big)$ de $\mathcal{C}_f$ et posons $h(x) = f(x-1)$. Pour que $h$ prenne la valeur $f(x_0)$, il faut $x - 1 = x_0$, c'est-à-dire $x = x_0 + 1$. Le point de la courbe de $h$ ayant la même ordonnée que $M$ est donc $M'(x_0+1\,;\,f(x_0))$ : il s'obtient en poussant $M$ d'une unité \textbf{vers la droite}. Chaque point subit le même sort, donc toute la courbe est translatée de $\vec{\imath}$. Le même raisonnement montre que $f(x)+k$ translate de $k\,\vec{\jmath}$.
\end{experience}

\begin{popfigure}
\begin{center}
\begin{tikzpicture}
\begin{axis}[
  width=0.95\linewidth, height=7.5cm,
  axis lines=middle, xlabel={$x$}, ylabel={$y$},
  xmin=-3.2, xmax=4.2, ymin=-3.5, ymax=5.5,
  xtick={-3,-2,-1,1,2,3,4}, ytick={-3,-2,-1,1,2,3,4,5},
  legend style={at={(0.02,0.98)}, anchor=north west, font=\small},
  grid=both, grid style={popline!40},
]
\addplot[popcurve, domain=-2.4:2.4, samples=200] {0.5*x^2 - 1};
\addlegendentry{$y=f(x)$}
\addplot[popcurve, poporange, domain=-2.4:2.4, samples=200] {0.5*x^2 + 1};
\addlegendentry{$y=f(x)+2$}
\addplot[popcurve, popgreen, domain=-1.4:3.4, samples=200] {0.5*(x-1)^2 - 1};
\addlegendentry{$y=f(x-1)$}
\addplot[popcurve, poppurple, domain=-2.4:2.4, samples=200] {-0.5*x^2 + 1};
\addlegendentry{$y=-f(x)$}
\draw[-{Stealth}, thick, poporange] (axis cs:1.6,0.28) -- (axis cs:1.6,2.28);
\draw[-{Stealth}, thick, popgreen] (axis cs:-1.5,0.125) -- (axis cs:-0.5,0.125);
\draw[-{Stealth}, thick, poppurple] (axis cs:2,1) -- (axis cs:2,-1);
\end{axis}
\end{tikzpicture}
\legende{La parabole $\mathcal{C}_f$ d'équation $y=\tfrac{1}{2}x^2-1$ (bleu) et trois courbes associées : translation verticale de $2\vec{\jmath}$ (orange), translation horizontale de $\vec{\imath}$ (vert), symétrie par rapport à $(Ox)$ (violet).}
\end{center}
\end{popfigure}

\begin{attention}[Piège classique : $f(x)+k$ contre $f(x+k)$]
C'est \textbf{l'erreur la plus fréquente} au Probatoire sur ce chapitre.
\begin{itemize}
\item $f(x) + k$ agit sur l'\textbf{ordonnée} : déplacement \textbf{vertical} de $k$ unités (vers le haut si $k>0$).
\item $f(x + k)$ agit sur l'\textbf{abscisse} : déplacement \textbf{horizontal} de $k$ unités, mais \textbf{vers la gauche} si $k>0$ (car $f(x+k) = f\big(x-(-k)\big)$).
\end{itemize}
Moyen mnémotechnique : ce qui est « à l'extérieur » de $f$ bouge la courbe verticalement dans le sens attendu ; ce qui est « à l'intérieur » la bouge horizontalement, \textbf{en sens inverse} du signe.
\end{attention}

\begin{exoresolu}[Construire une courbe associée avec une valeur absolue]
Soit $f(x) = x^2 - 2x$ et $g(x) = |f(x)|$. Déterminer les zéros de $f$, puis décrire la construction de $\mathcal{C}_g$ à partir de $\mathcal{C}_f$.
\tcblower
\textbf{Solution.} $f(x) = x(x-2)$, donc $f(x)=0$ pour $x=0$ ou $x=2$. Le trinôme $f$ est du signe de $a=1>0$ à l'extérieur de ses racines et négatif entre elles : $f(x) < 0$ sur $]0\,;\,2[$, et $f(x) \geq 0$ ailleurs. Pour tracer $\mathcal{C}_g$ :
\begin{itemize}
\item sur $]-\infty\,;\,0]$ et $[2\,;\,+\infty[$, $g(x) = f(x)$ : on garde $\mathcal{C}_f$ telle quelle ;
\item sur $]0\,;\,2[$, $g(x) = -f(x)$ : on retourne la partie de la parabole située sous l'axe $(Ox)$ par symétrie d'axe $(Ox)$ (le creux de sommet $(1\,;\,-1)$ devient une bosse de sommet $(1\,;\,1)$).
\end{itemize}
La courbe de $g$ est donc entièrement située au-dessus de $(Ox)$, ce qui est normal : une valeur absolue est toujours positive.
\end{exoresolu}

\courssub{Lecture graphique : conjecturer avant de démontrer}

Face à une courbe tracée (à la calculatrice, par un logiciel, ou donnée dans un énoncé), on peut \emph{lire} de précieuses informations. Attention au vocabulaire : une lecture graphique fournit une \cle{conjecture}, jamais une preuve.

\begin{methode}[Que peut-on conjecturer sur un graphique ?]
\begin{center}
\begin{tabularx}{\linewidth}{|l|X|}
\hline
\rowcolor{popblueL} \textbf{Information} & \textbf{Comment la lire sur la courbe} \\
\hline
Ensemble de définition & Les valeurs de $x$ pour lesquelles la courbe existe (projection de la courbe sur l'axe $(Ox)$). Un « trou » ou une asymptote verticale $x=a$ signale souvent une valeur interdite. \\
\hline
Sens de variation & La courbe « monte » (fonction croissante) ou « descend » (fonction décroissante) quand $x$ augmente. \\
\hline
Asymptote verticale $x = a$ & La courbe s'approche indéfiniment d'une droite verticale sans jamais la toucher. \\
\hline
Asymptote horizontale $y = b$ & Quand $x$ devient très grand (ou très petit), la courbe se plaque contre une droite horizontale. \\
\hline
Éléments de symétrie & Axe vertical $x = a$ (axe de symétrie) ou point $\Omega$ (centre de symétrie) visible sur le dessin. \\
\hline
\end{tabularx}
\end{center}
\end{methode}

\begin{popfigure}
\begin{center}
\begin{tikzpicture}
\begin{axis}[
  width=0.95\linewidth, height=7.5cm,
  axis lines=middle, xlabel={$x$}, ylabel={$y$},
  xmin=-6.5, xmax=6.5, ymin=-3.5, ymax=6.5,
  xtick={-6,-4,-2,2,4,6}, ytick={-3,-2,-1,1,2,3,4,5,6},
  grid=both, grid style={popline!40},
]
\addplot[popcurve, domain=0.5:6.3, samples=200] {2 + 3/x};
\addplot[popcurve, domain=-6.3:-0.5, samples=200] {2 + 3/x};
\addplot[dashed, thick, poporange, domain=-6.3:6.3] {2};
\addplot[dashed, thick, poppurple] coordinates {(0,-3.4) (0,6.4)};
\node[poporange, anchor=south west, font=\small] at (axis cs:3.5,2.15) {$y=2$};
\node[poppurple, anchor=west, font=\small, rotate=90] at (axis cs:0.25,-2.6) {$x=0$};
\end{axis}
\end{tikzpicture}
\legende{Courbe de $f(x)=2+\dfrac{3}{x}$ : on conjecture $D_f=\mathbb{R}^*$, une asymptote verticale $x=0$, une asymptote horizontale $y=2$, une décroissance sur chaque intervalle du domaine, et un centre de symétrie $\Omega(0\,;\,2)$.}
\end{center}
\end{popfigure}

\begin{exemplebox}[Lecture graphique complète, puis preuve]
Sur la figure ci-dessus, on conjecture pour $f(x) = 2 + \dfrac{3}{x}$ :
\begin{itemize}
\item $D_f = \mathbb{R}\setminus\{0\}$ (la courbe n'existe pas en $0$) ;
\item $f$ est décroissante sur $]-\infty\,;\,0[$ et sur $]0\,;\,+\infty[$ ;
\item asymptote verticale d'équation $x = 0$ et asymptote horizontale d'équation $y = 2$ ;
\item le point $\Omega(0\,;\,2)$ semble centre de symétrie.
\end{itemize}
\textbf{Vérification par le calcul} (la preuve !) : pour tout $x \neq 0$, on a $-x \neq 0$ (le domaine est symétrique par rapport à $0$) et
\[ \frac{f(-x)+f(x)}{2} = \frac{\left(2-\dfrac{3}{x}\right)+\left(2+\dfrac{3}{x}\right)}{2} = \frac{4}{2} = 2. \]
Les points d'abscisses $-x$ et $x$ sont symétriques par rapport à $\Omega(0\,;\,2)$ : $\Omega$ est bien centre de symétrie de $\mathcal{C}_f$.
\end{exemplebox}

\begin{attention}[Conjecture n'est pas preuve]
Une courbe dessinée est toujours \emph{partielle} : on ne voit qu'une fenêtre finie. Une fonction peut sembler croissante sur $[-5\,;\,5]$ et redescendre ensuite ; une asymptote apparente peut être un simple effet d'échelle. Au Probatoire, la rédaction correcte est : « \emph{graphiquement, on conjecture que...} » puis « \emph{démontrons par le calcul que...} ». Cette distinction deviendra cruciale en Terminale avec l'étude complète des fonctions et des limites.
\end{attention}

\begin{savaistu}[Une technique d'expert très rentable]
Les fonctions associées sont l'un des outils les plus rentables du programme. Plutôt que d'étudier de zéro la fonction $g(x) = (x-3)^2 + 2$, un bon élève reconnaît immédiatement la parabole de référence $x \mapsto x^2$, translatée de $3$ unités vers la droite puis de $2$ vers le haut : sommet en $(3\,;\,2)$, minimum égal à $2$, aucune intersection avec $(Ox)$... Tout cela sans aucun calcul de dérivée ! Cette technique, qui consiste à ramener une fonction compliquée à une \emph{fonction de référence} transformée, est au cœur des sujets de Probatoire et de Baccalauréat, et les mathématiciens l'utilisent partout : en traitement du signal (décalages temporels), en économie (courbes d'offre translatées par une taxe), ou en cartographie, où les courbes de niveau du relief du mont Cameroun se déduisent les unes des autres par simple translation verticale.
\end{savaistu}

\begin{exercice}[À toi de jouer]
Soit $f$ la fonction définie par $f(x) = \dfrac{1}{x}$ et $\mathcal{C}_f$ sa courbe.
\begin{enumerate}
\item Donner l'ensemble de définition de chacune des fonctions $g(x) = f(x) - 3$ et $h(x) = f(x+2)$.
\item Décrire précisément la transformation qui envoie $\mathcal{C}_f$ sur la courbe de $g$, puis sur celle de $h$.
\item Le point $A(4\,;\,1)$ appartient-il à la courbe de $h$ ?
\end{enumerate}
\end{exercice}

\begin{corrige}[Exercice 1]
\begin{enumerate}
\item $g(x) = \dfrac{1}{x} - 3$ est définie pour $x \neq 0$ : $D_g = \mathbb{R}^*$. $h(x) = \dfrac{1}{x+2}$ est définie pour $x + 2 \neq 0$, soit $D_h = \mathbb{R}\setminus\{-2\}$.
\item La courbe de $g$ s'obtient par translation de $\mathcal{C}_f$ de vecteur $-3\,\vec{\jmath}$ (3 unités vers le bas). Celle de $h$ s'obtient par translation de vecteur $-2\,\vec{\imath}$ (2 unités vers la gauche), car $h(x) = f\big(x-(-2)\big)$.
\item $4 \in D_h$ et $h(4) = \dfrac{1}{4+2} = \dfrac{1}{6} \neq 1$, donc $A \notin \mathcal{C}_h$.
\end{enumerate}
\end{corrige}

\newpage

\coursec{Activité d'intégration}

\courssub{Objectifs de l'activité}

À l'issue de cette activité d'intégration, tu seras capable de mobiliser \emph{simultanément} plusieurs savoir-faire de la leçon pour résoudre un problème concret de la vie quotidienne camerounaise. Plus précisément, tu sauras :

\begin{itemize}
\item justifier qu'une formule définit bien une \cle{application} et déterminer son \cle{ensemble image} ;
\item démontrer qu'une application est \cle{bijective} et \cle{expliciter sa bijection réciproque}, puis interpréter concrètement le résultat ;
\item construire et étudier une \cle{composée} de deux applications, en précisant son ensemble de définition ;
\item effectuer des \cle{opérations sur les fonctions} (différence) et exploiter le signe du résultat pour comparer deux tarifs ;
\item \cle{conjecturer par lecture graphique} un ensemble de définition, un sens de variation et des éléments de symétrie, puis \cle{vérifier par le calcul} qu'un point appartient à une courbe.
\end{itemize}

\courssub{Situation}

À Bafoussam, chef-lieu de la région de l'Ouest, la coopérative de transport urbain \og Nkap Express \fg{} assure la liaison entre le marché central et les quartiers périphériques (Famla, Tougang, Djeleng). Pour éviter les disputes entre chauffeurs et clients sur les prix, la coopérative a adopté une grille tarifaire unique : le prix d'une course dépend uniquement de la distance parcourue.

Pour une distance $d$ exprimée en kilomètres, le prix $p(d)$ en francs CFA est donné par :
\[ p(d) = 200 + 150\,d \qquad \text{pour } d \in [0\,;\,50]. \]
Le terme $200$ FCFA correspond à la \emph{prise en charge} (montant fixe payé dès la montée dans le véhicule) et $150$ FCFA est le prix par kilomètre parcouru. La coopérative n'accepte aucune course dépassant $50$ km.

L'État prélève sur chaque course une taxe de transport égale à $2\,\%$ du prix facturé. Pour un montant $x$ en francs CFA, cette taxe est donc donnée par :
\[ t(x) = 0{,}02\,x. \]

Une société concurrente, \og Wouri Trans \fg, vient de s'installer en ville et propose la tarification :
\[ q(d) = 100 + 180\,d \qquad \text{pour } d \in [0\,;\,50]. \]

Enfin, le mécanicien de la coopérative a tracé la courbe ci-dessous, qui représente le \emph{coût hebdomadaire d'entretien} $c(d)$ d'un véhicule (en milliers de FCFA) en fonction d'une variable $d$ (en dizaines de km) :
\[ c(d) = 0{,}1\,d^{2} + 5 \qquad \text{sur } [-30\,;\,30]. \]

\begin{popfigure}
\begin{center}
\begin{tikzpicture}
\begin{axis}[width=0.85\linewidth, height=6.5cm, axis lines=middle,
xlabel={$d$ (dizaines de km)}, ylabel={$c(d)$ (milliers FCFA)},
xmin=-33, xmax=33, ymin=0, ymax=100,
xtick={-30,-20,-10,0,10,20,30}, ytick={20,40,60,80},
grid=both, grid style={popline!40}, tick label style={font=\small}]
\addplot[popcurve,domain=-30:30,samples=100]{0.1*x^2+5};
\addplot[only marks, mark=*, mark size=1.5pt, color=poporange] coordinates {(10,15)};
\node[poporange, anchor=south west, font=\small] at (axis cs:10,15) {$A(10\,;\,15)$};
\end{axis}
\end{tikzpicture}
\end{center}
\legende{Courbe du coût hebdomadaire d'entretien $c(d) = 0{,}1d^2 + 5$ sur $[-30\,;\,30]$.}
\end{popfigure}

Le président de la coopérative te sollicite, toi l'élève de Première, pour répondre rigoureusement à ses questions de gestion.

\courssub{Consignes}

\begin{enumerate}
\item \textbf{Application et ensemble image.} Justifier que la formule $p(d) = 200 + 150d$ définit bien une application de $[0\,;\,50]$ dans $\mathbb{R}$. Déterminer l'ensemble image $p([0\,;\,50])$, c'est-à-dire l'ensemble de tous les prix possibles.
\item \textbf{Bijection et bijection réciproque.} Montrer que $p$ est une bijection de $[0\,;\,50]$ sur $p([0\,;\,50])$. Expliciter la bijection réciproque $p^{-1}$. Expliquer concrètement à quoi elle peut servir au chauffeur.
\item \textbf{Composition.} Déterminer la composée $t \circ p$, préciser son ensemble de définition et donner l'expression de $(t \circ p)(d)$. Interpréter le résultat : que représente $(t \circ p)(20)$ ? Le calculer.
\item \textbf{Comparaison de tarifs.} Préciser l'ensemble de définition commun à $p$ et $q$, calculer $(q - p)(d)$, étudier son signe et en déduire pour quelles distances la coopérative \og Nkap Express \fg{} est moins chère que sa concurrente.
\item \textbf{Lecture graphique et vérification.} Par lecture de la courbe de $c$ : conjecturer son ensemble de définition, son sens de variation et ses éléments de symétrie éventuels. Vérifier ensuite \emph{par le calcul} que le point $A(10\,;\,15)$ appartient à la courbe, et confirmer la symétrie observée en démontrant que $c$ est paire.
\end{enumerate}

\courssub{Ressources à mobiliser}

\begin{aretenir}[Boîte à outils de la leçon]
\begin{itemize}
\item \textbf{Application :} $f$ est une application de $E$ dans $F$ si tout élément de $E$ a \emph{exactement une} image dans $F$.
\item \textbf{Ensemble image :} $f(E) = \{\, f(x) \mid x \in E \,\}$. Pour une fonction affine strictement monotone sur $[a\,;\,b]$, l'image est l'intervalle $[f(a)\,;\,f(b)]$ si $f$ est croissante, et $[f(b)\,;\,f(a)]$ si $f$ est décroissante.
\item \textbf{Bijection :} $f : E \to F$ est bijective si tout élément de $F$ admet \emph{un unique} antécédent dans $E$. La bijection réciproque $f^{-1}$ s'obtient en résolvant $y = f(x)$ d'inconnue $x$.
\item \textbf{Composée :} $(t \circ p)(d) = t\bigl(p(d)\bigr)$ ; $d$ doit appartenir à l'ensemble de définition de $p$ \emph{et} $p(d)$ à celui de $t$.
\item \textbf{Différence de fonctions :} $(q-p)(d) = q(d) - p(d)$, définie sur l'\emph{intersection} des ensembles de définition de $q$ et de $p$.
\item \textbf{Parité :} $c$ est paire si son ensemble de définition est centré en $0$ et si $c(-d) = c(d)$ pour tout $d$ ; sa courbe est alors symétrique par rapport à l'axe des ordonnées.
\item \textbf{Appartenance à une courbe :} $A(x_A\,;\,y_A) \in \mathcal{C}_c$ si et seulement si $x_A \in D_c$ et $c(x_A) = y_A$.
\end{itemize}
\end{aretenir}

\courssub{Démarche et corrigé commenté}

\begin{exoresolu}[La coopérative de transport de Bafoussam]
Énoncé : répondre aux cinq questions du président de la coopérative \og Nkap Express \fg{} (voir les consignes ci-dessus).
\tcblower
\textbf{Question 1 : application et ensemble image.}

Pour tout $d \in [0\,;\,50]$, le nombre $200 + 150d$ existe et est unique (somme et produit de nombres réels). Donc $p$ définit bien une \cle{application} de $[0\,;\,50]$ dans $\mathbb{R}$.

La fonction $p$ est affine de coefficient directeur $150 > 0$ : elle est strictement croissante sur $[0\,;\,50]$. On calcule les images des bornes :
\[ p(0) = 200 + 150 \times 0 = 200 \qquad ; \qquad p(50) = 200 + 150 \times 50 = 200 + 7500 = 7700. \]
Par stricte croissance, l'ensemble image est :
\[ p\bigl([0\,;\,50]\bigr) = [200\,;\,7700]. \]
\emph{Interprétation :} toute course coûte entre $200$ FCFA (prise en charge seule) et $7700$ FCFA (course maximale de $50$ km).

\medskip
\textbf{Question 2 : bijection et bijection réciproque.}

Soit $y \in [200\,;\,7700]$. Résolvons l'équation $p(d) = y$ d'inconnue $d \in [0\,;\,50]$ :
\begin{align*}
200 + 150\,d &= y \\
150\,d &= y - 200 \\
d &= \frac{y - 200}{150}.
\end{align*}
Cette solution existe et est \emph{unique}. Vérifions qu'elle appartient bien à $[0\,;\,50]$ : comme $200 \leq y \leq 7700$, on a $0 \leq y - 200 \leq 7500$, donc en divisant par $150 > 0$ :
\[ 0 \leq \frac{y-200}{150} \leq \frac{7500}{150} = 50. \]

Ainsi, tout prix $y \in [200\,;\,7700]$ admet un unique antécédent dans $[0\,;\,50]$ : $p$ est une \cle{bijection} de $[0\,;\,50]$ sur $[200\,;\,7700]$, et sa bijection réciproque est :
\[ p^{-1}(y) = \frac{y - 200}{150}, \qquad y \in [200\,;\,7700]. \]
\emph{Interprétation concrète :} la bijection réciproque permet au chauffeur de \textbf{retrouver la distance à partir du prix}. Par exemple, si un client a payé $1700$ FCFA, la course mesurait :
\[ p^{-1}(1700) = \frac{1700 - 200}{150} = \frac{1500}{150} = 10 \text{ km}. \]
C'est l'outil idéal pour vérifier les réclamations des clients !

\medskip
\textbf{Question 3 : composée $t \circ p$.}

La fonction taxe $t(x) = 0{,}02x$ est définie sur $\mathbb{R}$ tout entier. Pour tout $d \in [0\,;\,50]$, $p(d)$ existe et appartient à $\mathbb{R}$, donc $t\bigl(p(d)\bigr)$ existe. L'ensemble de définition de $t \circ p$ est donc $[0\,;\,50]$, et :
\[ (t \circ p)(d) = t\bigl(p(d)\bigr) = 0{,}02\,(200 + 150\,d) = 4 + 3\,d. \]
\emph{Interprétation :} $(t \circ p)(d) = 4 + 3d$ est le \textbf{montant de la taxe} (en FCFA) prélevée sur une course de $d$ km. Pour une course de $20$ km :
\[ (t \circ p)(20) = 4 + 3 \times 20 = 64 \text{ FCFA}. \]
Vérification directe : $p(20) = 200 + 150 \times 20 = 200 + 3000 = 3200$ FCFA, et $2\,\%$ de $3200$ vaut bien $0{,}02 \times 3200 = 64$ FCFA. Les deux calculs concordent.

\medskip
\textbf{Question 4 : comparaison des tarifs.}

$p$ et $q$ sont toutes deux définies sur $[0\,;\,50]$ : l'ensemble de définition commun est $[0\,;\,50]$. Pour tout $d$ de cet intervalle :
\[ (q - p)(d) = q(d) - p(d) = (100 + 180\,d) - (200 + 150\,d) = 30\,d - 100. \]
Étudions le signe de cette différence :
\[ (q-p)(d) > 0 \iff 30\,d - 100 > 0 \iff d > \frac{100}{30} = \frac{10}{3} \approx 3{,}33. \]
On dresse le tableau de signes :
\begin{center}
\begin{tabular}{|c|ccccc|}\hline
$d$ & $0$ & & $\frac{10}{3}$ & & $50$ \\ \hline
$(q-p)(d)$ & & $-$ & $0$ & $+$ & \\ \hline
\end{tabular}
\end{center}
Conclusion :
\begin{itemize}
\item si $d > \dfrac{10}{3}$ km, alors $(q-p)(d) > 0$, c'est-à-dire $q(d) > p(d)$ : la coopérative \og Nkap Express \fg{} est \textbf{moins chère} ;
\item si $0 \leq d < \dfrac{10}{3}$ km, alors $(q-p)(d) < 0$, c'est-à-dire $q(d) < p(d)$ : c'est \og Wouri Trans \fg{} qui est moins chère (grâce à sa prise en charge plus faible) ;
\item si $d = \dfrac{10}{3}$ km $\approx 3{,}33$ km, les deux tarifs sont égaux : $p\!\left(\tfrac{10}{3}\right) = 200 + 150 \times \tfrac{10}{3} = 200 + 500 = 700$ FCFA et $q\!\left(\tfrac{10}{3}\right) = 100 + 180 \times \tfrac{10}{3} = 100 + 600 = 700$ FCFA.
\end{itemize}
\emph{Conseil au président :} la coopérative est compétitive sur les courses de plus de $3{,}33$ km, c'est-à-dire l'immense majorité des liaisons vers les quartiers périphériques.

\medskip
\textbf{Question 5 : lecture graphique et vérification.}

\emph{Conjectures par lecture graphique :}
\begin{itemize}
\item \textbf{Ensemble de définition :} la courbe s'étend de $d = -30$ à $d = 30$, donc $D_c = [-30\,;\,30]$ ;
\item \textbf{Variations :} la courbe \emph{descend} sur $[-30\,;\,0]$ puis \emph{monte} sur $[0\,;\,30]$ : $c$ semble décroissante sur $[-30\,;\,0]$ puis croissante sur $[0\,;\,30]$, avec un minimum en $0$ valant $c(0) = 5$ ;
\item \textbf{Symétrie :} les deux branches se superposent par pliage le long de l'axe des ordonnées : cet axe semble être \cle{axe de symétrie} de la courbe, ce qui laisse penser que $c$ est \cle{paire}.
\end{itemize}

\emph{Vérification par le calcul.} Pour le point $A(10\,;\,15)$ : on a $10 \in [-30\,;\,30]$ et
\[ c(10) = 0{,}1 \times 10^{2} + 5 = 0{,}1 \times 100 + 5 = 10 + 5 = 15. \]
Donc $c(10) = 15 = y_A$ : le point $A$ \textbf{appartient bien} à la courbe $\mathcal{C}_c$.

\emph{Confirmation de la parité :} l'ensemble $[-30\,;\,30]$ est centré en $0$ (si $d \in [-30\,;\,30]$, alors $-d \in [-30\,;\,30]$), et pour tout $d \in [-30\,;\,30]$ :
\[ c(-d) = 0{,}1\,(-d)^{2} + 5 = 0{,}1\,d^{2} + 5 = c(d). \]
La fonction $c$ est donc \cle{paire} : l'axe des ordonnées est bien axe de symétrie de sa courbe, ce qui confirme la conjecture graphique.
\end{exoresolu}

\begin{attention}[Ne pas confondre $t \circ p$ et $p \circ t$]
La composition n'est \textbf{pas commutative} : en général, $t \circ p \neq p \circ t$. Ici, $(t \circ p)(d) = 4 + 3d$ est la taxe sur une course de $d$ km, tandis que $(p \circ t)(x) = 200 + 150 \times 0{,}02x = 200 + 3x$ n'a aucune signification concrète dans la situation. Retiens l'ordre de lecture : dans $t \circ p$, c'est $p$ qui agit \emph{en premier} (on calcule d'abord le prix, puis la taxe).
\end{attention}

\courssub{Bilan de l'activité}

Cette activité t'a permis de constater que les notions de la leçon ne sont pas des objets abstraits, mais de véritables \textbf{outils de décision} :

\begin{itemize}
\item la notion d'\cle{application} et d'\cle{ensemble image} a donné un sens précis à la grille tarifaire et à l'ensemble des prix possibles ;
\item la \cle{bijection réciproque} a fourni un outil pratique : retrouver la distance à partir du prix payé, ce qui illustre le rôle \og inverseur \fg{} de $p^{-1}$ ;
\item la \cle{composition} a permis d'enchaîner deux traitements (tarification puis taxation) en une seule formule $4 + 3d$, en veillant à la compatibilité des ensembles de définition ;
\item les \cle{opérations sur les fonctions} (ici la différence $q - p$) ont transformé une question commerciale en une simple étude de signe, aboutissant au seuil décisif de $\dfrac{10}{3}$ km ;
\item enfin, la \cle{lecture graphique} a fourni des conjectures (variations, symétrie) que seul le \cle{calcul} (parité, appartenance d'un point) a permis de valider rigoureusement : \emph{le graphique suggère, le calcul prouve}.
\end{itemize}

Tu as ainsi mobilisé la quasi-totalité des savoir-faire de la leçon dans un contexte authentique de gestion d'une coopérative de transport camerounaise : c'est exactement l'esprit de l'approche par les compétences. Retiens la démarche générale : \textbf{modéliser} (traduire la situation en fonctions), \textbf{traiter} (calculer avec les outils du cours), \textbf{interpréter} (redonner sens au résultat dans la réalité).

\newpage

\coursec{Méthodes et exercices résolus}

\courssub{Méthode 1 : Déterminer par calculs l'ensemble de définition d'une fonction numérique}

\begin{methode}[Trouver l'ensemble de définition d'une fonction]
L'\cle{ensemble de définition} d'une fonction $f$, noté $D_f$, est l'ensemble des réels $x$ pour lesquels $f(x)$ existe (est calculable). Pour le déterminer \emph{par le calcul} :
\begin{enumerate}
\item \textbf{Repérer les écritures « à risque »} dans l'expression de $f(x)$ : dénominateurs, racines carrées (plus généralement racines d'indice pair).
\item \textbf{Traduire chaque contrainte} :
\begin{itemize}
\item un \textbf{dénominateur} doit être \emph{non nul} : $g(x)\neq 0$ ;
\item une quantité sous une \textbf{racine carrée} doit être \emph{positive ou nulle} : $g(x)\geq 0$.
\end{itemize}
\item \textbf{Résoudre} chaque condition (équation, inéquation, tableau de signes si produit ou quotient).
\item \textbf{Conclure} en écrivant $D_f$ sous forme d'intervalle ou de réunion d'intervalles, et \emph{toujours} en précisant les valeurs exclues.
\end{enumerate}
\textbf{Réflexes :} un polynôme est défini sur $\mathbb{R}$ tout entier ; pour une fonction quotient, on commence \emph{toujours} par chercher les valeurs qui annulent le dénominateur ; quand plusieurs conditions coexistent, elles doivent être vérifiées \emph{simultanément} (on prend l'intersection des ensembles solutions).
\end{methode}

\begin{exoresolu}[Ensemble de définition d'une fonction avec racine et dénominateur]
On considère la fonction $f$ définie par $f(x)=\dfrac{\sqrt{x+2}}{x-3}$.
\begin{enumerate}
\item Déterminer l'ensemble de définition $D_f$ de $f$.
\item Le réel $7$ appartient-il à $D_f$ ? Si oui, calculer $f(7)$.
\end{enumerate}
\tcblower
\textbf{1. Détermination de $D_f$.}

L'expression de $f(x)$ contient deux écritures « à risque » : une racine carrée au numérateur et un dénominateur.

\textbf{Condition 1 (racine carrée) :} la quantité sous le radical doit être positive ou nulle :
\[x+2\geq 0 \iff x\geq -2.\]

\textbf{Condition 2 (dénominateur) :} le dénominateur doit être non nul :
\[x-3\neq 0 \iff x\neq 3.\]

Les deux conditions doivent être vérifiées \emph{simultanément}. On doit donc avoir $x\geq -2$ \textbf{et} $x\neq 3$. Comme $3\geq -2$, la valeur $3$ appartient à la zone $x\geq -2$ et doit en être retirée.

\textbf{Conclusion :}
\[D_f=[-2\,;\,3[\ \cup\ ]3\,;\,+\infty[.\]

\textbf{2. Étude du réel $7$.}

$7\in D_f$ car $7>3$. On peut donc calculer :
\[f(7)=\dfrac{\sqrt{7+2}}{7-3}=\dfrac{\sqrt{9}}{4}=\dfrac{3}{4}.\]
\textbf{Conclusion :} $7\in D_f$ et $f(7)=\dfrac{3}{4}$.
\end{exoresolu}

\begin{aretenir}[Restriction d'une fonction]
Soit $f$ une fonction d'ensemble de définition $D_f$ et $I$ un intervalle (ou une partie) contenu dans $D_f$. La \cle{restriction} de $f$ à $I$, souvent notée $f_{|I}$, est la fonction définie sur $I$ par la même formule : pour tout $x\in I$, $f_{|I}(x)=f(x)$. On change \emph{l'ensemble de départ}, pas la règle de calcul. Par exemple, la restriction de $x\mapsto x^2$ à $[0\,;\,+\infty[$ est strictement croissante, alors que $x\mapsto x^2$ ne l'est pas sur $\mathbb{R}$ : restreindre le domaine peut donner à une fonction des propriétés nouvelles (idée qui servira pour construire des bijections).
\end{aretenir}

\courssub{Méthode 2 : Justifier qu'une application est injective, surjective, bijective}

\begin{methode}[Injectivité, surjectivité, bijectivité]
Soit $f:E\to F$ une \emph{application} (tout élément de $E$ a une image dans $F$).
\begin{enumerate}
\item \textbf{Injectivité} : montrer que deux éléments distincts ont des images distinctes. En pratique, on part de l'égalité $f(a)=f(b)$ avec $a,b\in E$ et on démontre que $a=b$.
\item \textbf{Surjectivité} : montrer que tout élément $y\in F$ possède au moins un antécédent dans $E$. En pratique, on résout l'équation $f(x)=y$ d'inconnue $x\in E$ et on vérifie que la solution trouvée \emph{appartient bien à} $E$.
\item \textbf{Bijectivité} : montrer que $f$ est \emph{à la fois} injective et surjective, ou montrer directement que l'équation $f(x)=y$ admet, pour tout $y\in F$, une \textbf{unique} solution dans $E$. La bijection réciproque $f^{-1}$ est alors définie par : $f^{-1}(y)$ est l'unique solution de $f(x)=y$.
\end{enumerate}
\textbf{Réflexe :} pour prouver qu'une application n'est \emph{pas} injective, il suffit d'exhiber un contre-exemple : deux réels distincts ayant la même image. Pour prouver qu'elle n'est \emph{pas} surjective, il suffit d'exhiber un élément de $F$ sans antécédent.
\end{methode}

\begin{exoresolu}[Bijectivité d'une application affine]
On considère l'application $f:\mathbb{R}\to\mathbb{R}$ définie par $f(x)=3x-5$.
\begin{enumerate}
\item Montrer que $f$ est injective.
\item Montrer que $f$ est surjective.
\item Conclure et expliciter la bijection réciproque $f^{-1}$.
\end{enumerate}
\tcblower
\textbf{1. Injectivité.}

Soient $a$ et $b$ deux réels tels que $f(a)=f(b)$. Alors :
\begin{align*}
3a-5 &= 3b-5\\
3a &= 3b \quad \text{(en ajoutant } 5 \text{ aux deux membres)}\\
a &= b \quad \text{(en divisant par } 3\neq 0\text{)}.
\end{align*}
Deux réels ayant la même image sont donc égaux : $f$ est \textbf{injective}.

\textbf{2. Surjectivité.}

Soit $y$ un réel quelconque. Cherchons $x\in\mathbb{R}$ tel que $f(x)=y$ :
\begin{align*}
3x-5 &= y\\
3x &= y+5\\
x &= \dfrac{y+5}{3}.
\end{align*}
Pour tout $y\in\mathbb{R}$, le réel $x=\dfrac{y+5}{3}$ existe, appartient bien à $\mathbb{R}$ et vérifie $f(x)=y$. Donc tout réel admet un antécédent : $f$ est \textbf{surjective}.

\textbf{3. Conclusion.}

$f$ est injective et surjective, donc $f$ est \textbf{bijective} de $\mathbb{R}$ vers $\mathbb{R}$. La résolution précédente donne directement l'antécédent unique de $y$ : la bijection réciproque est
\[f^{-1}:\mathbb{R}\to\mathbb{R},\qquad f^{-1}(y)=\dfrac{y+5}{3}.\]
\textbf{Vérification :} $f\circ f^{-1}(y)=3\times\dfrac{y+5}{3}-5=y$ et $f^{-1}\circ f(x)=\dfrac{(3x-5)+5}{3}=x$. La formule est correcte.
\end{exoresolu}

\begin{exemplebox}[Une application non injective et non surjective]
Soit $g:\mathbb{R}\to\mathbb{R}$ définie par $g(x)=x^2$.
\begin{itemize}
\item $g$ n'est \textbf{pas injective} : $g(2)=g(-2)=4$ alors que $2\neq -2$.
\item $g$ n'est \textbf{pas surjective} : le réel $-1$ n'a pas d'antécédent, car l'équation $x^2=-1$ n'a pas de solution réelle.
\end{itemize}
En revanche, l'application $h:[0\,;\,+\infty[\ \to\ [0\,;\,+\infty[$ définie par $h(x)=x^2$ est bijective : c'est la restriction de $g$ qui « répare » les deux défauts, et sa bijection réciproque est $h^{-1}(y)=\sqrt{y}$.
\end{exemplebox}

\courssub{Méthode 3 : Déterminer la composée de deux applications}

\begin{methode}[Composer deux applications]
Soient $f:E\to F$ et $g:F\to G$ deux applications. La \cle{composée} de $f$ par $g$, notée $g\circ f$, est définie par :
\[(g\circ f)(x)=g\bigl(f(x)\bigr).\]
\begin{enumerate}
\item \textbf{Identifier l'ordre} : dans $g\circ f$, c'est $f$ qui agit \emph{en premier} (on lit de droite à gauche).
\item \textbf{Vérifier la compatibilité} : les valeurs prises par $f(x)$ doivent appartenir à l'ensemble de départ de $g$.
\item \textbf{Calculer} en remplaçant, dans l'expression de $g$, la variable par l'expression complète de $f(x)$, entre parenthèses.
\item \textbf{Simplifier} et donner l'ensemble de définition de la composée : $D_{g\circ f}=\{x\in D_f \mid f(x)\in D_g\}$.
\end{enumerate}
\textbf{Attention :} en général $g\circ f\neq f\circ g$ : la composition n'est \emph{pas commutative}.
\end{methode}

\begin{exoresolu}[Composées de deux fonctions]
On considère les fonctions $f$ et $g$ définies sur $\mathbb{R}$ par $f(x)=2x+1$ et $g(x)=x^2-3$.
\begin{enumerate}
\item Déterminer l'expression de $(g\circ f)(x)$.
\item Déterminer l'expression de $(f\circ g)(x)$. A-t-on $g\circ f=f\circ g$ ?
\item Calculer $(g\circ f)(2)$.
\end{enumerate}
\tcblower
\textbf{1. Calcul de $g\circ f$.}

Dans $g\circ f$, on applique d'abord $f$, puis $g$. On remplace donc $x$ par $f(x)=2x+1$ dans l'expression de $g$ :
\begin{align*}
(g\circ f)(x) &= g\bigl(f(x)\bigr) = g(2x+1)\\
&= (2x+1)^2-3\\
&= 4x^2+4x+1-3\\
&= 4x^2+4x-2.
\end{align*}
Donc $(g\circ f)(x)=4x^2+4x-2$, définie sur $\mathbb{R}$ (polynôme).

\textbf{2. Calcul de $f\circ g$.}

On applique d'abord $g$, puis $f$ :
\begin{align*}
(f\circ g)(x) &= f\bigl(g(x)\bigr) = f(x^2-3)\\
&= 2(x^2-3)+1\\
&= 2x^2-6+1\\
&= 2x^2-5.
\end{align*}
Donc $(f\circ g)(x)=2x^2-5$. Or $4x^2+4x-2\neq 2x^2-5$ (par exemple en $x=0$ : $(g\circ f)(0)=-2$ et $(f\circ g)(0)=-5$).

\textbf{Conclusion :} $g\circ f\neq f\circ g$ : la composition n'est pas commutative.

\textbf{3. Calcul de $(g\circ f)(2)$.}

Deux méthodes : directement, $(g\circ f)(2)=4(2)^2+4(2)-2=16+8-2=22$ ; ou étape par étape : $f(2)=2\times 2+1=5$, puis $g(5)=5^2-3=25-3=22$.

\textbf{Conclusion :} $(g\circ f)(2)=22$.
\end{exoresolu}

\begin{exoresolu}[Ensemble de définition d'une composée]
On considère les fonctions $f(x)=\sqrt{x}$ et $g(x)=\dfrac{1}{x-4}$. Déterminer l'ensemble de définition de $g\circ f$.
\tcblower
On a $D_f=[0\,;\,+\infty[$ et $D_g=\mathbb{R}\setminus\{4\}$.

D'après la formule $D_{g\circ f}=\{x\in D_f \mid f(x)\in D_g\}$, il faut :
\begin{itemize}
\item $x\in D_f$, c'est-à-dire $x\geq 0$ ;
\item $f(x)\in D_g$, c'est-à-dire $\sqrt{x}\neq 4$, soit $x\neq 16$.
\end{itemize}
\textbf{Conclusion :} $D_{g\circ f}=[0\,;\,16[\ \cup\ ]16\,;\,+\infty[$, et pour tout $x$ de cet ensemble, $(g\circ f)(x)=\dfrac{1}{\sqrt{x}-4}$.
\end{exoresolu}

\courssub{Méthode 4 : Opérations sur les polynômes et les fonctions}

\begin{methode}[Somme, produit, quotient de polynômes et de fonctions]
\textbf{Sur les polynômes :}
\begin{enumerate}
\item \textbf{Somme} : on additionne les coefficients des termes de \emph{même degré}.
\item \textbf{Produit} : on distribue chaque terme de l'un sur chaque terme de l'autre, puis on réduit. Degré du produit : $\deg(PQ)=\deg P+\deg Q$.
\item \textbf{Quotient} $\dfrac{P}{Q}$ : il n'existe que pour les valeurs qui n'annulent pas $Q$ ; condition d'existence : $Q(x)\neq 0$.
\end{enumerate}
\textbf{Sur les fonctions} $f$ et $g$ d'ensembles de définition $D_f$ et $D_g$ :
\begin{enumerate}
\item $D_{f+g}=D_{fg}=D_f\cap D_g$ ;
\item $D_{f/g}=\{x\in D_f\cap D_g \mid g(x)\neq 0\}$ ;
\item $D_{g\circ f}=\{x\in D_f \mid f(x)\in D_g\}$.
\end{enumerate}
\textbf{Réflexe :} pour un quotient, on écrit toujours \emph{deux} conditions : appartenir aux deux domaines \emph{et} ne pas annuler le dénominateur.
\end{methode}

\begin{exoresolu}[Somme, produit et quotient]
On considère les polynômes $P(x)=2x^2-3x+1$ et $Q(x)=x-2$, et les fonctions $f(x)=\sqrt{x}$ et $g(x)=\dfrac{1}{x-1}$.
\begin{enumerate}
\item Calculer $P(x)+Q(x)$ et $P(x)\times Q(x)$.
\item Donner la condition d'existence du quotient $\dfrac{P(x)}{Q(x)}$.
\item Déterminer l'ensemble de définition de la fonction $h=f+g$.
\end{enumerate}
\tcblower
\textbf{1. Somme et produit des polynômes.}

\textbf{Somme :} on regroupe les termes de même degré :
\[P(x)+Q(x)=2x^2-3x+1+x-2=2x^2-2x-1.\]

\textbf{Produit :} on distribue :
\begin{align*}
P(x)\times Q(x) &= (2x^2-3x+1)(x-2)\\
&= 2x^3-4x^2-3x^2+6x+x-2\\
&= 2x^3-7x^2+7x-2.
\end{align*}
On vérifie le degré : $\deg P+\deg Q=2+1=3$, cohérent avec le résultat.

\textbf{2. Condition d'existence du quotient.}

Le quotient $\dfrac{P(x)}{Q(x)}$ existe si et seulement si le dénominateur est non nul :
\[Q(x)\neq 0 \iff x-2\neq 0 \iff x\neq 2.\]
Le quotient est donc défini sur $\mathbb{R}\setminus\{2\}=]-\infty\,;\,2[\ \cup\ ]2\,;\,+\infty[$.

\textbf{3. Ensemble de définition de $h=f+g$.}

$D_f=[0\,;\,+\infty[$ (racine carrée) et $D_g=\mathbb{R}\setminus\{1\}$ (dénominateur). Donc :
\[D_h=D_f\cap D_g=[0\,;\,+\infty[\ \setminus\ \{1\}=[0\,;\,1[\ \cup\ ]1\,;\,+\infty[.\]
\textbf{Conclusion :} $h$ est définie sur $[0\,;\,1[\ \cup\ ]1\,;\,+\infty[$.
\end{exoresolu}

\courssub{Méthode 5 : Montrer qu'une fonction est paire, impaire ou périodique}

\begin{methode}[Parité et périodicité]
\textbf{Parité — étapes obligatoires :}
\begin{enumerate}
\item \textbf{Vérifier que $D_f$ est symétrique par rapport à $0$} : pour tout $x\in D_f$, on a $-x\in D_f$. Si ce n'est pas le cas, la fonction n'est ni paire ni impaire, et on s'arrête là.
\item \textbf{Calculer $f(-x)$} et le comparer à $f(x)$ :
\begin{itemize}
\item si $f(-x)=f(x)$ pour tout $x\in D_f$ : $f$ est \cle{paire} (courbe symétrique par rapport à l'axe des ordonnées) ;
\item si $f(-x)=-f(x)$ pour tout $x\in D_f$ : $f$ est \cle{impaire} (courbe symétrique par rapport à l'origine).
\end{itemize}
\end{enumerate}
\textbf{Périodicité :} pour montrer que $f$ est périodique de période $T$ ($T>0$), on vérifie que pour tout $x\in D_f$, $x+T\in D_f$ et $f(x+T)=f(x)$. La courbe se reproduit alors à l'identique sur tout intervalle d'amplitude $T$ : il suffit de l'étudier sur un tel intervalle, puis de compléter par translations de vecteurs $kT\vec{i}$ ($k\in\mathbb{Z}$).
\end{methode}

\begin{exoresolu}[Parité d'une fonction rationnelle]
On considère la fonction $f$ définie par $f(x)=\dfrac{x}{x^2-4}$.
\begin{enumerate}
\item Déterminer l'ensemble de définition $D_f$ et vérifier qu'il est symétrique par rapport à $0$.
\item Étudier la parité de $f$.
\item En déduire un élément de symétrie de la courbe de $f$.
\end{enumerate}
\tcblower
\textbf{1. Ensemble de définition.}

Le dénominateur doit être non nul :
\[x^2-4\neq 0 \iff (x-2)(x+2)\neq 0 \iff x\neq 2 \text{ et } x\neq -2.\]
Donc $D_f=\mathbb{R}\setminus\{-2\,;\,2\}$.

\textbf{Symétrie de $D_f$ :} si $x\in D_f$, alors $x\neq 2$ et $x\neq -2$, donc $-x\neq -2$ et $-x\neq 2$, c'est-à-dire $-x\in D_f$. L'ensemble $D_f$ est bien symétrique par rapport à $0$ : l'étude de parité a un sens.

\textbf{2. Calcul de $f(-x)$.}

Pour tout $x\in D_f$ :
\[f(-x)=\dfrac{-x}{(-x)^2-4}=\dfrac{-x}{x^2-4}=-\dfrac{x}{x^2-4}=-f(x).\]
(On a utilisé $(-x)^2=x^2$.)

\textbf{Conclusion :} pour tout $x\in D_f$, $f(-x)=-f(x)$ : la fonction $f$ est \textbf{impaire}.

\textbf{3. Élément de symétrie.}

Puisque $f$ est impaire, sa courbe représentative est symétrique par rapport à l'\textbf{origine $O$ du repère}. Il suffit donc d'étudier $f$ sur $[0\,;\,+\infty[\ \cap\ D_f$ et de compléter par symétrie centrale de centre $O$.
\end{exoresolu}

\begin{exemplebox}[Périodicité]
Soit $f$ la fonction définie sur $\mathbb{R}$ par $f(x)=x-\lfloor x\rfloor$ (partie fractionnaire de $x$). Pour tout réel $x$, $\lfloor x+1\rfloor=\lfloor x\rfloor+1$, donc :
\[f(x+1)=(x+1)-\lfloor x+1\rfloor=x+1-\lfloor x\rfloor-1=x-\lfloor x\rfloor=f(x).\]
La fonction $f$ est donc \textbf{périodique de période $1$} : sa courbe est formée de segments identiques répétés sur chaque intervalle $[n\,;\,n+1[$. On retrouve la même idée avec les fonctions trigonométriques : $\cos$ et $\sin$ sont périodiques de période $2\pi$.
\end{exemplebox}

\courssub{Méthode 6 : Justifier un axe ou un centre de symétrie d'une courbe}

\begin{methode}[Éléments de symétrie d'une courbe]
Soit $f$ une fonction de courbe $\mathcal{C}_f$ dans un repère.
\begin{enumerate}
\item \textbf{Axe de symétrie d'équation $x=a$} : montrer que pour tout réel $h$ tel que $a+h$ et $a-h$ appartiennent à $D_f$ :
\[f(a+h)=f(a-h).\]
\item \textbf{Centre de symétrie $\Omega(a\,;\,b)$} : montrer que pour tout réel $h$ tel que $a+h$ et $a-h$ appartiennent à $D_f$ :
\[f(a+h)+f(a-h)=2b.\]
\item \textbf{Conclure} en citant la propriété : « donc la droite d'équation $x=a$ est axe de symétrie de $\mathcal{C}_f$ » ou « donc $\Omega(a\,;\,b)$ est centre de symétrie de $\mathcal{C}_f$ ».
\end{enumerate}
\textbf{Réflexe :} penser à $a$ comme le « milieu » : $a+h$ et $a-h$ sont deux points symétriques par rapport à $a$. Remarque : la parité est un cas particulier — $f$ paire correspond à l'axe $x=0$, $f$ impaire au centre $O(0\,;\,0)$.
\end{methode}

\begin{exoresolu}[Centre de symétrie d'une courbe]
Soit $f$ la fonction définie sur $\mathbb{R}\setminus\{1\}$ par $f(x)=\dfrac{2x+1}{x-1}$, de courbe $\mathcal{C}_f$.
\begin{enumerate}
\item Montrer que pour tout $x\in D_f$, $f(x)=2+\dfrac{3}{x-1}$.
\item Démontrer que le point $\Omega(1\,;\,2)$ est centre de symétrie de $\mathcal{C}_f$.
\end{enumerate}
\tcblower
\textbf{1. Réécriture de $f(x)$.}

Pour $x\neq 1$ :
\[2+\dfrac{3}{x-1}=\dfrac{2(x-1)+3}{x-1}=\dfrac{2x-2+3}{x-1}=\dfrac{2x+1}{x-1}=f(x).\]
L'égalité est démontrée.

\textbf{2. Centre de symétrie.}

Soit $h$ un réel non nul. Les réels $1+h$ et $1-h$ sont différents de $1$, donc appartiennent à $D_f$. Calculons avec la forme obtenue à la question 1 :
\begin{align*}
f(1+h) &= 2+\dfrac{3}{(1+h)-1}=2+\dfrac{3}{h},\\[2pt]
f(1-h) &= 2+\dfrac{3}{(1-h)-1}=2-\dfrac{3}{h}.
\end{align*}
En additionnant :
\[f(1+h)+f(1-h)=\left(2+\dfrac{3}{h}\right)+\left(2-\dfrac{3}{h}\right)=4=2\times 2.\]
On a donc, pour tout $h$ convenable, $f(1+h)+f(1-h)=2b$ avec $a=1$ et $b=2$.

\textbf{Conclusion :} le point $\Omega(1\,;\,2)$ est \textbf{centre de symétrie} de la courbe $\mathcal{C}_f$.
\end{exoresolu}

\courssub{Méthode 7 : Courbe d'une fonction, appartenance d'un point et lecture graphique}

\begin{methode}[Appartenance d'un point à une courbe et fonctions associées]
\begin{enumerate}
\item \textbf{Appartenance :} le point $M(x_0\,;\,y_0)$ appartient à la courbe $\mathcal{C}_f$ si et seulement si $x_0\in D_f$ \emph{et} $f(x_0)=y_0$. On calcule donc $f(x_0)$ et on compare à $y_0$.
\item \textbf{Lecture graphique :} l'ensemble de définition se lit sur l'axe des abscisses (projection de la courbe) ; le sens de variation se lit de gauche à droite (la courbe « monte » ou « descend ») ; une asymptote verticale se devine quand la courbe s'approche d'une droite verticale sans la toucher ; les symétries se conjecturent en pliant mentalement la feuille. Toute conjecture graphique doit ensuite être \emph{démontrée par le calcul}.
\item \textbf{Fonctions associées :} à partir de la courbe de $f$, on obtient celle de :
\begin{itemize}
\item $x\mapsto f(x)+k$ : translation de vecteur $k\vec{j}$ ;
\item $x\mapsto f(x-a)$ : translation de vecteur $a\vec{i}$ ;
\item $x\mapsto -f(x)$ : symétrie par rapport à l'axe des abscisses ;
\item $x\mapsto f(-x)$ : symétrie par rapport à l'axe des ordonnées.
\end{itemize}
\end{enumerate}
\end{methode}

\begin{exoresolu}[Appartenance à une courbe et courbe associée]
Soit $f$ la fonction définie sur $\mathbb{R}$ par $f(x)=x^2-2x+3$, de courbe $\mathcal{C}_f$.
\begin{enumerate}
\item Les points $A(2\,;\,3)$ et $B(-1\,;\,5)$ appartiennent-ils à $\mathcal{C}_f$ ?
\item Montrer que la droite d'équation $x=1$ est axe de symétrie de $\mathcal{C}_f$.
\item On note $g$ la fonction définie par $g(x)=f(x)-4$. Comment obtient-on la courbe de $g$ à partir de $\mathcal{C}_f$ ?
\end{enumerate}
\tcblower
\textbf{1. Appartenance des points.}

\textbf{Point $A$ :} $2\in D_f=\mathbb{R}$ et $f(2)=2^2-2\times 2+3=4-4+3=3$. Or $3$ est bien l'ordonnée de $A$. Donc $A(2\,;\,3)\in\mathcal{C}_f$.

\textbf{Point $B$ :} $f(-1)=(-1)^2-2\times(-1)+3=1+2+3=6$. Or l'ordonnée de $B$ est $5\neq 6$. Donc $B(-1\,;\,5)\notin\mathcal{C}_f$.

\textbf{2. Axe de symétrie.}

Soit $h$ un réel quelconque. Calculons :
\begin{align*}
f(1+h) &= (1+h)^2-2(1+h)+3 = 1+2h+h^2-2-2h+3 = h^2+2,\\
f(1-h) &= (1-h)^2-2(1-h)+3 = 1-2h+h^2-2+2h+3 = h^2+2.
\end{align*}
Pour tout réel $h$, $f(1+h)=f(1-h)$.

\textbf{Conclusion :} la droite d'équation $x=1$ est \textbf{axe de symétrie} de $\mathcal{C}_f$. (On reconnaît l'axe de la parabole, d'abscisse du sommet $x=\dfrac{-(-2)}{2\times 1}=1$.)

\textbf{3. Courbe de $g$.}

Pour tout réel $x$, $g(x)=f(x)+(-4)$ : la courbe de $g$ s'obtient à partir de $\mathcal{C}_f$ par la \textbf{translation de vecteur $-4\vec{j}$} (glissement de 4 unités vers le bas).
\end{exoresolu}

\begin{popfigure}
\begin{center}
\begin{tikzpicture}[scale=0.85]
\draw[->,popmuted] (-2.5,0) -- (4.5,0) node[below left] {$x$};
\draw[->,popmuted] (0,-4.5) -- (0,5.5) node[below left] {$y$};
\draw[popblue,thick,domain=-1:3,samples=100] plot (\x,{\x*\x-2*\x+3});
\draw[poporange,thick,domain=-1:3,samples=100] plot (\x,{\x*\x-2*\x-1});
\draw[popline,dashed] (1,-4.5) -- (1,5.5);
\node[popblue,right] at (2.6,4.4) {$\mathcal{C}_f$};
\node[poporange,right] at (2.6,-0.6) {$\mathcal{C}_g$};
\node[popmuted,right] at (1.05,5.2) {$x=1$};
\fill[popdark] (2,3) circle (1.5pt) node[above right,popdark] {$A$};
\foreach \p in {(1,0),(0,1)} \fill[popmuted] \p circle (1pt);
\node[below,popmuted] at (1,0) {\small $1$};
\node[left,popmuted] at (0,1) {\small $1$};
\end{tikzpicture}
\end{center}
\legende{La parabole $\mathcal{C}_f$ d'axe de symétrie $x=1$, et $\mathcal{C}_g$ obtenue par translation de vecteur $-4\vec{j}$.}
\end{popfigure}

\begin{savaistu}[Pourquoi les bijections sont partout]
Quand tu composes un numéro Orange ou MTN, le réseau associe à ton numéro un identifiant unique, puis le retrouve sans ambiguïté : c'est une \emph{bijection} entre numéros et abonnés. De même, le numéro matricule d'un candidat au Probatoire code un et un seul candidat : sans injectivité, deux élèves recevraient le même relevé de notes ! En mathématiques, la bijection réciproque $f^{-1}$ permet de « remonter » du résultat à la donnée : c'est exactement ce que fait un décodeur.
\end{savaistu}

\begin{attention}[Les 5 erreurs qui coûtent des points au Probatoire]
\begin{enumerate}
\item \textbf{Oublier une condition dans l'ensemble de définition.} Pour $f(x)=\dfrac{\sqrt{x+2}}{x-3}$, écrire seulement « $x\neq 3$ » en oubliant $x\geq -2$ (ou l'inverse) fait perdre l'essentiel des points. Recense \emph{toutes} les écritures à risque avant de conclure, et écris le résultat sous forme d'intervalles.
\item \textbf{Étudier la parité sans vérifier la symétrie de $D_f$.} Calculer $f(-x)$ n'a de sens que si $D_f$ est symétrique par rapport à $0$. Cette vérification, rédigée en une ligne, est exigée dans la copie.
\item \textbf{Confondre $g\circ f$ et $f\circ g$.} Dans $g\circ f$, c'est $f$ qui agit en premier. Erreur classique : écrire $(g\circ f)(x)=f\bigl(g(x)\bigr)$. Retiens : on lit la composition \emph{de droite à gauche}.
\item \textbf{Conclure « bijective » après avoir seulement résolu $f(x)=y$.} Trouver un antécédent prouve la surjectivité ; il faut encore justifier son \emph{unicité} (injectivité). Rédige les deux étapes séparément, ou montre directement que la solution est unique.
\item \textbf{Affirmer qu'un point appartient à une courbe sans calcul.} Dire « $A(2\,;\,3)\in\mathcal{C}_f$ car ça se voit » n'est pas une justification : il faut calculer $f(2)$, obtenir $3$, et conclure par l'égalité $f(2)=3$. De même, pour un centre de symétrie, la relation $f(a+h)+f(a-h)=2b$ doit être \emph{démontrée}, pas annoncée.
\end{enumerate}
\end{attention}

\newpage

\coursec{Exercices}

\courssub{Je vérifie mes connaissances}

\begin{exercice}[QCM : une seule bonne réponse]
Pour chaque question, coche la case correspondant à la bonne réponse.
\begin{enumerate}
\item L'ensemble de définition de $f(x)=\dfrac{1}{x-5}$ est :
\begin{itemize}
\item[\textbf{a)}] $\mathbb{R}$
\item[\textbf{b)}] $\mathbb{R}\setminus\{5\}$
\item[\textbf{c)}] $\mathbb{R}\setminus\{-5\}$
\item[\textbf{d)}] $]5\,;\,+\infty[$
\end{itemize}

\item L'ensemble de définition de $g(x)=\sqrt{3x-6}$ est :
\begin{itemize}
\item[\textbf{a)}] $[2\,;\,+\infty[$
\item[\textbf{b)}] $]-\infty\,;\,2]$
\item[\textbf{c)}] $\mathbb{R}$
\item[\textbf{d)}] $]2\,;\,+\infty[$
\end{itemize}

\item Une application $f:E\to F$ est bijective si et seulement si :
\begin{itemize}
\item[\textbf{a)}] tout élément de $F$ a au plus un antécédent ;
\item[\textbf{b)}] tout élément de $F$ a au moins un antécédent ;
\item[\textbf{c)}] tout élément de $F$ a exactement un antécédent ;
\item[\textbf{d)}] deux éléments distincts de $E$ ont des images distinctes.
\end{itemize}

\item Si $f$ est une fonction impaire, alors sa courbe admet :
\begin{itemize}
\item[\textbf{a)}] l'axe des ordonnées comme axe de symétrie ;
\item[\textbf{b)}] l'origine du repère comme centre de symétrie ;
\item[\textbf{c)}] l'axe des abscisses comme axe de symétrie ;
\item[\textbf{d)}] aucun élément de symétrie.
\end{itemize}
\end{enumerate}
\end{exercice}

\begin{exercice}[Vrai ou faux, en justifiant]
Réponds par \textbf{Vrai} ou \textbf{Faux} et justifie chaque réponse par une démonstration courte ou un contre-exemple.
\begin{enumerate}
\item La fonction $f$ définie sur $\mathbb{R}$ par $f(x)=x^2$ est injective.
\item La fonction $f$ définie sur $\mathbb{R}$ par $f(x)=2x-3$ est bijective.
\item Si $f$ est paire et impaire à la fois, alors $f$ est la fonction nulle.
\item Pour toutes fonctions $f$ et $g$, on a $f\circ g = g\circ f$.
\item La fonction $x\mapsto \dfrac{1}{x}$ est définie en $0$.
\end{enumerate}
\end{exercice}

\begin{exercice}[Définitions à compléter]
Recopie et complète les phrases suivantes.
\begin{enumerate}
\item On appelle \cle{ensemble de définition} d'une fonction numérique $f$ l'ensemble des réels $x$ pour lesquels $f(x)$ \emph{existe} (ou \emph{a un sens}).
\item Une application $f:E\to F$ est \cle{injective} lorsque deux éléments distincts de $E$ ont \emph{des images distinctes}.
\item Une application $f:E\to F$ est \cle{surjective} lorsque tout élément de $F$ \emph{a au moins un antécédent dans $E$}.
\item Une fonction $f$ est \cle{paire} lorsque son ensemble de définition est symétrique par rapport à $0$ et que pour tout $x$ de cet ensemble, $f(-x)=f(x)$.
\item Une fonction $f$ est \cle{périodique} de période $T$ (avec $T>0$) lorsque pour tout $x$ de son ensemble de définition, $x+T$ appartient à cet ensemble et $f(x+T)=f(x)$.
\end{enumerate}
\end{exercice}

\begin{exercice}[Calculs directs sur les polynômes]
On considère les polynômes $P(x)=x^3-2x+1$ et $Q(x)=x^2+x$.
\begin{enumerate}
\item Calculer $P(x)+Q(x)$ et préciser son degré.
\item Calculer $P(x)\times Q(x)$ et préciser son degré.
\item Déterminer l'ensemble de définition de la fonction quotient $x \mapsto \dfrac{P(x)}{Q(x)}$.
\end{enumerate}
\end{exercice}

\courssub{Je m'entraîne}

\begin{exercice}[Ensembles de définition et restriction]
Déterminer par le calcul l'ensemble de définition de chacune des fonctions suivantes :
\[ f(x)=\frac{3x-1}{x^2-4}\;;\qquad g(x)=\sqrt{5-2x}\;;\qquad h(x)=\frac{\sqrt{x+1}}{x-3}. \]
Donner ensuite la restriction de $f$ à l'intervalle $]2\,;\,+\infty[$.
\end{exercice}

\begin{exercice}[Bijection et bijection réciproque]
Soit $f$ l'application définie par $f(x)=\dfrac{2x+3}{x-1}$.
\begin{enumerate}
\item Vérifier que $f$ est bien définie sur $\mathbb{R}\setminus\{1\}$ et que pour tout $x\neq 1$, $f(x)\neq 2$.
\item Montrer que $f$ est une bijection de $\mathbb{R}\setminus\{1\}$ sur $\mathbb{R}\setminus\{2\}$.
\item Expliciter sa bijection réciproque $f^{-1}$.
\end{enumerate}
\end{exercice}

\begin{exercice}[Composées de fonctions]
Soient $f$ et $g$ les fonctions définies par $f(x)=x^2-1$ et $g(x)=\sqrt{x}$.
\begin{enumerate}
\item Déterminer l'expression de $g\circ f(x)$ et l'ensemble de définition de $g\circ f$.
\item Déterminer l'expression de $f\circ g(x)$ et l'ensemble de définition de $f\circ g$.
\item Comparer $g\circ f$ et $f\circ g$. Que remarque-t-on ?
\end{enumerate}
\end{exercice}

\begin{exercice}[Parité de fonctions]
Pour chacune des fonctions suivantes, déterminer d'abord l'ensemble de définition, vérifier qu'il est symétrique par rapport à $0$, puis étudier la parité :
\[ f(x)=x^3-3x\;;\qquad g(x)=\frac{x^2}{x^2+1}\;;\qquad h(x)=x+\frac{1}{x}. \]
Préciser, pour chaque fonction paire ou impaire, la conséquence géométrique sur sa courbe représentative.
\end{exercice}

\courssub{Je me prépare au Probatoire}

\begin{exercice}[Étude complète d'une fonction homographique]
Un opérateur de téléphonie mobile propose à ses abonnés de Yaoundé une formule dont le coût mensuel, en milliers de francs CFA, est modélisé par la fonction
\[ f(x)=\frac{3x+1}{x+2}, \]
où $x$ désigne le nombre de gigaoctets (Go) consommés dans le mois ($x\geq 0$).
\begin{enumerate}
\item Déterminer l'ensemble de définition $D_f$ de $f$ sur $\mathbb{R}$.
\item Montrer que pour tout $x\in D_f$, on a $f(x)=3-\dfrac{5}{x+2}$.
\item Montrer que le point $\Omega(-2\,;\,3)$ est centre de symétrie de la courbe $(\mathcal{C}_f)$ de $f$.
\item Montrer que $f$ réalise une bijection de $\mathbb{R}\setminus\{-2\}$ sur $\mathbb{R}\setminus\{3\}$, et expliciter sa bijection réciproque $f^{-1}$.
\item Les points $A(0\,;\,\tfrac{1}{2})$ et $B(3\,;\,2)$ appartiennent-ils à $(\mathcal{C}_f)$ ? Justifier par le calcul.
\item Application pratique : un abonné a payé \num{2,6} milliers de francs CFA. En utilisant $f^{-1}$, déterminer le nombre de Go qu'il a consommés.
\end{enumerate}
\end{exercice}

\begin{exercice}[Lecture graphique et confirmation par le calcul]
\emph{(Exercice type Probatoire.)} On donne ci-dessous la courbe représentative $(\mathcal{C})$ d'une fonction $f$ sur l'intervalle $[-4\,;\,4]$.

\begin{popfigure}
\begin{center}
\begin{tikzpicture}[scale=0.9]
\draw[popline, very thin, step=1] (-4.4,-3.4) grid (4.4,4.4);
\draw[->, popdark, thick] (-4.5,0) -- (4.5,0) node[below right] {$x$};
\draw[->, popdark, thick] (0,-3.5) -- (0,4.5) node[above left] {$y$};
\draw[popdark] (0,0) node[below left] {$O$};
\foreach \x in {-4,-3,-2,-1,1,2,3,4} \draw[popdark] (\x,0.08) -- (\x,-0.08) node[below] {\footnotesize $\x$};
\foreach \y in {-3,-2,-1,1,2,3,4} \draw[popdark] (0.08,\y) -- (-0.08,\y) node[left] {\footnotesize $\y$};
\draw[popblue, very thick, domain=-4:4, samples=200] plot (\x,{0.25*\x*\x - 2});
\draw[popblue, fill=popblue] (-4,2) circle (1.5pt);
\draw[popblue, fill=popblue] (4,2) circle (1.5pt);
\end{tikzpicture}
\end{center}
\legende{Courbe représentative $(\mathcal{C})$ de la fonction $f$ sur $[-4\,;\,4]$}
\end{popfigure}

\textbf{Partie A : lecture graphique.}
\begin{enumerate}
\item Conjecturer l'ensemble de définition de $f$.
\item Déterminer graphiquement l'image de $2$ par $f$, puis les antécédents éventuels de $0$.
\item Dresser le tableau de variations conjecturé de $f$ sur $[-4\,;\,4]$.
\item Conjecturer les éléments de symétrie éventuels de $(\mathcal{C})$.
\end{enumerate}

\textbf{Partie B : confirmation par le calcul.} On admet que $f(x)=\dfrac{x^2}{4}-2$.
\begin{enumerate}
\item[5.] Vérifier par le calcul que les points $M(2\,;\,-1)$ et $N(-4\,;\,2)$ appartiennent à $(\mathcal{C})$.
\item[6.] Démontrer que la droite d'équation $x=0$ est axe de symétrie de $(\mathcal{C})$.
\item[7.] Résoudre par le calcul l'équation $f(x)=0$ et comparer avec la lecture graphique de la question 2.
\end{enumerate}
\end{exercice}

\begin{exercice}[Fonctions associées, symétries et périodicité]
Soit $f$ la fonction définie sur $\mathbb{R}$ par $f(x)=x^2-4x+5$, et $g$ la fonction définie par $g(x)=2+\dfrac{3}{x-1}$.
\begin{enumerate}
\item Montrer que pour tout réel $x$, $f(x)=(x-2)^2+1$.
\item En déduire que la droite d'équation $x=2$ est axe de symétrie de la courbe $(\mathcal{C}_f)$.
\item Déterminer l'ensemble de définition de $g$, puis montrer que le point $\Omega(1\,;\,2)$ est centre de symétrie de la courbe $(\mathcal{C}_g)$.
\item On considère la fonction $h$ définie sur $\mathbb{R}$ par $h(x)=\sin(2x)$. Montrer que $h$ est périodique de période $\pi$ et étudier sa parité.
\item À partir de la courbe de $f$, décrire précisément les transformations permettant de construire les courbes des fonctions $x\mapsto f(x)-2$, $x\mapsto f(x+1)$ et $x\mapsto |f(x)|$. La courbe de $x\mapsto |f(x)|$ diffère-t-elle de celle de $f$ ? Justifier à l'aide de la question 1.
\end{enumerate}
\end{exercice}

\newpage

\coursec{Corrigés des exercices}

\begin{corrige}[Exercice 1 — QCM : une seule bonne réponse]
\begin{enumerate}
\item \textbf{Réponse b)} : $\mathbb{R}\setminus\{5\}$. En effet, $f(x)=\dfrac{1}{x-5}$ existe si et seulement si $x-5\neq 0$, c'est-à-dire $x\neq 5$.
\item \textbf{Réponse a)} : $[2\,;\,+\infty[$. En effet, $\sqrt{3x-6}$ existe si et seulement si $3x-6\geq 0$, c'est-à-dire $x\geq 2$.
\item \textbf{Réponse c)} : une application est bijective si et seulement si elle est à la fois injective (au plus un antécédent) et surjective (au moins un antécédent), donc si et seulement si tout élément de $F$ a \textbf{exactement un} antécédent dans $E$.
\item \textbf{Réponse b)} : si $f$ est impaire, alors $f(-x)=-f(x)$, ce qui se traduit géométriquement par : l'\textbf{origine du repère est centre de symétrie} de la courbe.
\end{enumerate}
\end{corrige}

\begin{corrige}[Exercice 2 — Vrai ou faux, en justifiant]
\begin{enumerate}
\item \textbf{Faux.} Contre-exemple : $f(-2)=(-2)^2=4$ et $f(2)=4$ : deux réels distincts $-2$ et $2$ ont la même image, donc $f$ n'est pas injective sur $\mathbb{R}$.
\item \textbf{Vrai.} Soit $y\in\mathbb{R}$. L'équation $2x-3=y$ équivaut à $x=\dfrac{y+3}{2}$ : tout réel $y$ admet un antécédent et un seul. Donc $f$ est bijective de $\mathbb{R}$ sur $\mathbb{R}$ (sa bijection réciproque est $f^{-1}(y)=\dfrac{y+3}{2}$).
\item \textbf{Vrai.} Si $f$ est paire et impaire, alors pour tout $x$ de son ensemble de définition : $f(-x)=f(x)$ (parité) et $f(-x)=-f(x)$ (imparité). D'où $f(x)=-f(x)$, soit $2f(x)=0$, donc $f(x)=0$ pour tout $x$ : $f$ est la fonction nulle.
\item \textbf{Faux.} Contre-exemple : $f(x)=x^2$ et $g(x)=x+1$. Alors $f\circ g(x)=(x+1)^2$ tandis que $g\circ f(x)=x^2+1$. En $x=1$ : $f\circ g(1)=4$ et $g\circ f(1)=2$ : les fonctions sont différentes. La composition n'est pas commutative en général.
\item \textbf{Faux.} La fonction inverse n'est pas définie en $0$ car on ne peut pas diviser par $0$ : son ensemble de définition est $\mathbb{R}\setminus\{0\}=\mathbb{R}^*$.
\end{enumerate}
\end{corrige}

\begin{corrige}[Exercice 3 — Définitions à compléter]
Les réponses sont les mots en gras ci-dessous :
\begin{enumerate}
\item $f(x)$ \textbf{existe} (ou \textbf{a un sens}) ;
\item des \textbf{images distinctes} ;
\item \textbf{a au moins un antécédent} dans $E$ ;
\item $f(-x)=f(x)$ ;
\item $f(x+T)=f(x)$.
\end{enumerate}
\end{corrige}

\begin{corrige}[Exercice 4 — Calculs directs sur les polynômes]
\begin{enumerate}
\item Somme :
\[ P(x)+Q(x)=x^3-2x+1+x^2+x=x^3+x^2-x+1. \]
Le terme de plus haut degré est $x^3$ : le degré de $P+Q$ est $\boxed{3}$.
\item Produit :
\begin{align*}
P(x)\times Q(x)&=(x^3-2x+1)(x^2+x)\\
&=x^5+x^4-2x^3-2x^2+x^2+x\\
&=x^5+x^4-2x^3-x^2+x.
\end{align*}
Le degré du produit est la somme des degrés : $3+2=\boxed{5}$.
\item Le quotient $\dfrac{P(x)}{Q(x)}$ existe si et seulement si $Q(x)\neq 0$. Or $Q(x)=x^2+x=x(x+1)$ s'annule en $0$ et en $-1$. L'ensemble de définition est donc :
\[ D=\mathbb{R}\setminus\{-1\,;\,0\}. \]
\end{enumerate}
\end{corrige}

\begin{corrige}[Exercice 5 — Ensembles de définition et restriction]
\textbf{Fonction $f$.} $f(x)=\dfrac{3x-1}{x^2-4}$ existe si et seulement si $x^2-4\neq 0$. Or $x^2-4=(x-2)(x+2)$ s'annule en $2$ et $-2$. Donc :
\[ D_f=\mathbb{R}\setminus\{-2\,;\,2\}. \]

\textbf{Fonction $g$.} $\sqrt{5-2x}$ existe si et seulement si $5-2x\geq 0$, c'est-à-dire $x\leq \dfrac{5}{2}$. Donc :
\[ D_g=\left]-\infty\,;\,\tfrac{5}{2}\right]. \]

\textbf{Fonction $h$.} Deux conditions simultanées : la quantité sous le radical doit être positive, $x+1\geq 0$ soit $x\geq -1$, et le dénominateur non nul, $x-3\neq 0$ soit $x\neq 3$. Donc :
\[ D_h=[-1\,;\,+\infty[\setminus\{3\}=[-1\,;\,3[\;\cup\;]3\,;\,+\infty[. \]

\textbf{Restriction de $f$ à $]2\,;\,+\infty[$.} C'est la fonction $\tilde{f}$ définie sur $]2\,;\,+\infty[$ par $\tilde{f}(x)=\dfrac{3x-1}{x^2-4}$ : même expression que $f$, mais ensemble de départ réduit à $]2\,;\,+\infty[$ (sur lequel $f$ est bien définie puisque $x>2$ implique $x\neq \pm 2$).
\end{corrige}

\begin{corrige}[Exercice 6 — Bijection et bijection réciproque]
\begin{enumerate}
\item Le dénominateur $x-1$ s'annule uniquement en $x=1$ : $f$ est bien définie sur $\mathbb{R}\setminus\{1\}$. Supposons $f(x)=2$ : alors $2x+3=2(x-1)$, soit $2x+3=2x-2$, c'est-à-dire $3=-2$, ce qui est absurde. Donc pour tout $x\neq 1$, $f(x)\neq 2$ : $f$ prend ses valeurs dans $\mathbb{R}\setminus\{2\}$.
\item Soit $y\in\mathbb{R}\setminus\{2\}$. Résolvons dans $\mathbb{R}\setminus\{1\}$ l'équation $f(x)=y$ :
\begin{align*}
\frac{2x+3}{x-1}=y &\iff 2x+3=y(x-1) \quad (\text{car } x\neq 1)\\
&\iff 2x+3=yx-y\\
&\iff x(2-y)=-y-3\\
&\iff x=\frac{y+3}{y-2} \quad (\text{licite car } y\neq 2).
\end{align*}
Vérifions que cet antécédent n'est pas égal à $1$ : $\dfrac{y+3}{y-2}=1$ donnerait $y+3=y-2$, soit $3=-2$, absurde. Donc tout $y\in\mathbb{R}\setminus\{2\}$ admet un antécédent et un seul dans $\mathbb{R}\setminus\{1\}$ : $f$ est une \textbf{bijection} de $\mathbb{R}\setminus\{1\}$ sur $\mathbb{R}\setminus\{2\}$.
\item D'après le calcul précédent, la bijection réciproque est :
\[ f^{-1}(y)=\frac{y+3}{y-2},\qquad y\in\mathbb{R}\setminus\{2\}. \]
\end{enumerate}
\end{corrige}

\begin{corrige}[Exercice 7 — Composées de fonctions]
\begin{enumerate}
\item $g\circ f(x)=g\bigl(f(x)\bigr)=\sqrt{x^2-1}$. Cette expression existe si et seulement si $x^2-1\geq 0$, c'est-à-dire $(x-1)(x+1)\geq 0$, soit $x\leq -1$ ou $x\geq 1$. Donc :
\[ D_{g\circ f}=\left]-\infty\,;\,-1\right]\cup\left[1\,;\,+\infty\right[. \]
\item $f\circ g(x)=f\bigl(g(x)\bigr)=\bigl(\sqrt{x}\bigr)^2-1=x-1$. Pour que le calcul ait un sens, il faut d'abord que $g(x)$ existe, c'est-à-dire $x\geq 0$ ; ensuite $f$ est définie partout. Donc :
\[ D_{f\circ g}=[0\,;\,+\infty[. \]
\item On a $g\circ f(x)=\sqrt{x^2-1}$ et $f\circ g(x)=x-1$ : les expressions sont différentes, et même les ensembles de définition sont différents. On remarque que \textbf{la composition n'est pas commutative} : en général $g\circ f\neq f\circ g$.
\end{enumerate}
\end{corrige}

\begin{attention}[Piège classique sur l'ensemble de définition d'une composée]
Même si l'expression de $f\circ g$ se simplifie en $x-1$, son ensemble de définition n'est \textbf{pas} $\mathbb{R}$ : la condition $x\geq 0$ imposée par $g$ subsiste. Pour déterminer $D_{f\circ g}$, on exige toujours : $x\in D_g$ \emph{et} $g(x)\in D_f$.
\end{attention}

\begin{corrige}[Exercice 8 — Parité de fonctions]
\textbf{Fonction $f(x)=x^3-3x$.} $D_f=\mathbb{R}$, symétrique par rapport à $0$. Pour tout réel $x$ :
\[ f(-x)=(-x)^3-3(-x)=-x^3+3x=-(x^3-3x)=-f(x). \]
$f$ est \textbf{impaire} : sa courbe admet l'origine $O$ du repère comme centre de symétrie.

\textbf{Fonction $g(x)=\dfrac{x^2}{x^2+1}$.} Pour tout réel $x$, $x^2+1\geq 1>0$ : le dénominateur ne s'annule jamais, donc $D_g=\mathbb{R}$, symétrique par rapport à $0$. De plus :
\[ g(-x)=\frac{(-x)^2}{(-x)^2+1}=\frac{x^2}{x^2+1}=g(x). \]
$g$ est \textbf{paire} : sa courbe admet l'axe des ordonnées (droite $x=0$) comme axe de symétrie.

\textbf{Fonction $h(x)=x+\dfrac{1}{x}$.} $h(x)$ existe si et seulement si $x\neq 0$ : $D_h=\mathbb{R}^*$, symétrique par rapport à $0$. Pour tout $x\neq 0$ :
\[ h(-x)=-x+\frac{1}{-x}=-x-\frac{1}{x}=-\left(x+\frac{1}{x}\right)=-h(x). \]
$h$ est \textbf{impaire} : sa courbe admet l'origine $O$ comme centre de symétrie.
\end{corrige}

\begin{corrige}[Exercice 9 — Étude complète d'une fonction homographique]
\emph{Barème indicatif (sur 10 pts) : question 1 : 1 pt ; question 2 : 1,5 pt ; question 3 : 2 pts ; question 4 : 2,5 pts ; question 5 : 1,5 pt ; question 6 : 1,5 pt.}

\begin{enumerate}
\item $f(x)=\dfrac{3x+1}{x+2}$ existe si et seulement si $x+2\neq 0$, c'est-à-dire $x\neq -2$. Donc :
\[ D_f=\mathbb{R}\setminus\{-2\}. \]
\item Pour tout $x\in D_f$ :
\[ 3-\frac{5}{x+2}=\frac{3(x+2)-5}{x+2}=\frac{3x+6-5}{x+2}=\frac{3x+1}{x+2}=f(x). \]
\item Pour montrer que $\Omega(-2\,;\,3)$ est centre de symétrie de $(\mathcal{C}_f)$, on vérifie que $D_f$ est symétrique par rapport à $-2$ (c'est le cas : si $x\neq -2$ alors $-4-x\neq -2$) et que pour tout $x\in D_f$ :
\[ f(-4-x)+f(x)=2\times 3=6. \]
Calculons avec la forme de la question 2 :
\begin{align*}
f(-4-x)+f(x)&=\left(3-\frac{5}{-4-x+2}\right)+\left(3-\frac{5}{x+2}\right)\\
&=6-\frac{5}{-x-2}-\frac{5}{x+2}=6+\frac{5}{x+2}-\frac{5}{x+2}=6.
\end{align*}
Donc $\dfrac{f(-4-x)+f(x)}{2}=3$ : le point $\Omega(-2\,;\,3)$ est bien \textbf{centre de symétrie} de $(\mathcal{C}_f)$.
\item Soit $y\in\mathbb{R}\setminus\{3\}$. Résolvons $f(x)=y$ dans $\mathbb{R}\setminus\{-2\}$ :
\begin{align*}
\frac{3x+1}{x+2}=y &\iff 3x+1=y(x+2)\\
&\iff 3x-yx=2y-1\\
&\iff x(3-y)=2y-1\\
&\iff x=\frac{2y-1}{3-y} \quad (\text{licite car } y\neq 3).
\end{align*}
Cet antécédent est différent de $-2$ : $\dfrac{2y-1}{3-y}=-2$ donnerait $2y-1=-6+2y$, soit $-1=-6$, absurde. Ainsi tout $y\in\mathbb{R}\setminus\{3\}$ a un antécédent unique : $f$ est une \textbf{bijection} de $\mathbb{R}\setminus\{-2\}$ sur $\mathbb{R}\setminus\{3\}$, et :
\[ f^{-1}(y)=\frac{2y-1}{3-y},\qquad y\in\mathbb{R}\setminus\{3\}. \]
\item $f(0)=\dfrac{1}{2}$ : le point $A(0\,;\,\tfrac{1}{2})$ \textbf{appartient} à $(\mathcal{C}_f)$. $f(3)=\dfrac{3\times 3+1}{3+2}=\dfrac{10}{5}=2$ : le point $B(3\,;\,2)$ \textbf{appartient} à $(\mathcal{C}_f)$.
\item Le coût payé est $y=\num{2,6}$ milliers de francs. Le nombre de Go consommés est :
\[ x=f^{-1}(\num{2,6})=\frac{2\times\num{2,6}-1}{3-\num{2,6}}=\frac{\num{4,2}}{\num{0,4}}=\num{10,5}. \]
L'abonné a consommé $\boxed{\num{10,5}\ \text{Go}}$.
\end{enumerate}
\end{corrige}

\begin{attention}[Points de vigilance — Exercice 9]
\begin{itemize}
\item Citer la condition $y\neq 3$ avant de diviser par $3-y$, et vérifier que l'antécédent trouvé appartient bien à l'ensemble de départ.
\item Pour le centre de symétrie $\Omega(a\,;\,b)$, la relation attendue est $f(2a-x)+f(x)=2b$ pour tout $x$ tel que $x$ et $2a-x$ soient dans $D_f$ ; la symétrie de $D_f$ par rapport à $a$ doit être mentionnée.
\end{itemize}
\end{attention}

\begin{corrige}[Exercice 10 — Lecture graphique et confirmation par le calcul]
\emph{Barème indicatif (sur 10 pts) : Partie A : 4 pts (1 pt par question) ; Partie B : 6 pts (2 pts par question).}

\textbf{Partie A : lecture graphique.}
\begin{enumerate}
\item La courbe est tracée pour $x$ allant de $-4$ à $4$ (points extrêmes pleins) : on conjecture $D_f=[-4\,;\,4]$.
\item Graphiquement, l'image de $2$ est $f(2)=-1$ (la courbe passe par le point d'abscisse $2$ et d'ordonnée $-1$). Les antécédents de $0$ sont les abscisses des points d'intersection de $(\mathcal{C})$ avec l'axe des abscisses : on lit environ $x\approx -\num{2,8}$ et $x\approx \num{2,8}$.
\item Tableau de variations conjecturé :
\begin{center}
\begin{tabular}{|c|ccccc|}\hline
$x$ & $-4$ & & $0$ & & $4$ \\ \hline
$f$ & $2$ & $\searrow$ & $-2$ & $\nearrow$ & $2$ \\ \hline
\end{tabular}
\end{center}
$f$ semble décroissante sur $[-4\,;\,0]$ puis croissante sur $[0\,;\,4]$, avec un minimum $-2$ atteint en $0$.
\item La courbe semble symétrique par rapport à l'axe des ordonnées (droite d'équation $x=0$).
\end{enumerate}

\textbf{Partie B : confirmation par le calcul.} On admet $f(x)=\dfrac{x^2}{4}-2$.
\begin{enumerate}
\item[5.] $f(2)=\dfrac{4}{4}-2=1-2=-1$ : $M(2\,;\,-1)\in(\mathcal{C})$. $f(-4)=\dfrac{16}{4}-2=4-2=2$ : $N(-4\,;\,2)\in(\mathcal{C})$.
\item[6.] $D_f=[-4\,;\,4]$ est symétrique par rapport à $0$, et pour tout $x\in D_f$ :
\[ f(-x)=\frac{(-x)^2}{4}-2=\frac{x^2}{4}-2=f(x). \]
$f$ est paire, donc la droite d'équation $x=0$ (axe des ordonnées) est \textbf{axe de symétrie} de $(\mathcal{C})$.
\item[7.] $f(x)=0\iff \dfrac{x^2}{4}=2\iff x^2=8\iff x=-2\sqrt{2}$ ou $x=2\sqrt{2}$. Or $2\sqrt{2}\approx \num{2,83}$, ce qui confirme la lecture graphique de la question 2 ($x\approx \pm\num{2,8}$).
\end{enumerate}
\end{corrige}

\begin{attention}[Points de vigilance — Exercice 10]
\begin{itemize}
\item Une lecture graphique donne une \emph{conjecture} : elle doit être confirmée par le calcul (questions 5 à 7).
\item Pour l'axe de symétrie $x=0$, la justification attendue est la parité : symétrie de $D_f$ puis $f(-x)=f(x)$.
\end{itemize}
\end{attention}

\begin{corrige}[Exercice 11 — Fonctions associées, symétries et périodicité]
\begin{enumerate}
\item Pour tout réel $x$ :
\[ (x-2)^2+1=x^2-4x+4+1=x^2-4x+5=f(x). \]
\item $D_f=\mathbb{R}$ est symétrique par rapport à $2$. Pour tout réel $x$, en utilisant la forme canonique :
\[ f(4-x)=(4-x-2)^2+1=(2-x)^2+1=(x-2)^2+1=f(x). \]
On a donc $f(2\times 2-x)=f(x)$ pour tout $x$ : la droite d'équation $x=2$ est \textbf{axe de symétrie} de $(\mathcal{C}_f)$.
\item $g(x)=2+\dfrac{3}{x-1}$ existe si et seulement si $x\neq 1$ : $D_g=\mathbb{R}\setminus\{1\}$, ensemble symétrique par rapport à $1$ (si $x\neq 1$, alors $2-x\neq 1$). Pour tout $x\in D_g$ :
\begin{align*}
g(2-x)+g(x)&=\left(2+\frac{3}{2-x-1}\right)+\left(2+\frac{3}{x-1}\right)\\
&=4+\frac{3}{1-x}+\frac{3}{x-1}=4-\frac{3}{x-1}+\frac{3}{x-1}=4=2\times 2.
\end{align*}
Donc $\dfrac{g(2-x)+g(x)}{2}=2$ : le point $\Omega(1\,;\,2)$ est \textbf{centre de symétrie} de $(\mathcal{C}_g)$.
\item Pour tout réel $x$ :
\[ h(x+\pi)=\sin\bigl(2(x+\pi)\bigr)=\sin(2x+2\pi)=\sin(2x)=h(x), \]
car la fonction sinus est périodique de période $2\pi$. Donc $h$ est périodique de période $\pi$. Parité : $D_h=\mathbb{R}$ est symétrique par rapport à $0$ et
\[ h(-x)=\sin(-2x)=-\sin(2x)=-h(x), \]
car la fonction sinus est impaire. Donc $h$ est \textbf{impaire} : sa courbe admet l'origine comme centre de symétrie.
\item Transformations à partir de $(\mathcal{C}_f)$ :
\begin{itemize}
\item Courbe de $x\mapsto f(x)-2$ : translation de $(\mathcal{C}_f)$ de vecteur $-2\vec{j}$ (translation verticale vers le bas de 2 unités) ;
\item Courbe de $x\mapsto f(x+1)$ : translation de $(\mathcal{C}_f)$ de vecteur $-\vec{i}$ (translation horizontale vers la gauche de 1 unité) ;
\item Courbe de $x\mapsto |f(x)|$ : on conserve la partie de $(\mathcal{C}_f)$ située au-dessus de l'axe des abscisses et on symétrise par rapport à l'axe des abscisses la partie située en dessous.
\end{itemize}
Or d'après la question 1, $f(x)=(x-2)^2+1\geq 1>0$ pour tout réel $x$ : $f$ est strictement positive, donc $|f(x)|=f(x)$ pour tout $x$. La courbe de $x\mapsto|f(x)|$ est donc \textbf{confondue} avec celle de $f$.
\end{enumerate}
\end{corrige}

\begin{attention}[Points de vigilance — Exercice 11]
\begin{itemize}
\item Pour l'axe de symétrie $x=a$, la relation à établir est $f(2a-x)=f(x)$ ; pour le centre $\Omega(a\,;\,b)$, c'est $f(2a-x)+f(x)=2b$.
\item Attention au sens de la translation : la courbe de $x\mapsto f(x+1)$ s'obtient en décalant $(\mathcal{C}_f)$ vers la \emph{gauche}, pas vers la droite.
\end{itemize}
\end{attention}

\newpage

\coursec{Fiche bilan}

\begin{popfigure}
\begin{tikzpicture}[mindmap, grow cyclic, every node/.style={concept, text=white, font=\small},
  root concept/.append style={concept color=popdark, font=\bfseries},
  level 1/.append style={level distance=3.4cm, sibling angle=60, minimum size=2.4cm},
  level 2/.append style={level distance=2.2cm, minimum size=1.7cm, font=\scriptsize}]
\node[root concept]{Généralités sur les fonctions}
  child[concept color=popblue]{ node {Ensemble de définition}
    child[concept color=popblue!70]{ node {Restriction} }
    child[concept color=popblue!70]{ node {Valeurs interdites} }
  }
  child[concept color=poporange]{ node {Injectivité, surjectivité, bijectivité}
    child[concept color=poporange!70]{ node {Bijection réciproque} }
  }
  child[concept color=popgreen]{ node {Composition $g\circ f$}
    child[concept color=popgreen!70]{ node {Ordre des fonctions} }
  }
  child[concept color=poppurple]{ node {Opérations : $f+g$, $f\times g$, $\frac{f}{g}$}
    child[concept color=poppurple!70]{ node {Polynômes} }
  }
  child[concept color=poppink]{ node {Parité, périodicité}
    child[concept color=poppink!70]{ node {Symétries} }
  }
  child[concept color=popgold]{ node {Courbe $\mathcal{C}_f$}
    child[concept color=popgold!70]{ node {Fonctions associées} }
    child[concept color=popgold!70]{ node {Lecture graphique} }
  };
\end{tikzpicture}
\legende{Carte mentale de la leçon : les six grandes notions à maîtriser.}
\end{popfigure}

\begin{bilan}
\textbf{Définitions clés}
\begin{itemize}
  \item \cle{Ensemble de définition} de $f$ : $D_f=\{x\in\mathbb{R}\mid f(x)\text{ existe}\}$. On exclut : dénominateurs nuls, quantités négatives sous une racine carrée.
  \item \cle{Restriction} de $f$ à $I\subset D_f$ : même règle de calcul, mais définie seulement sur $I$.
  \item $f:E\to F$ est \cle{injective} si $f(a)=f(b)\Rightarrow a=b$ ; \cle{surjective} si tout $y\in F$ admet un antécédent ; \cle{bijective} si elle est injective et surjective (chaque $y$ a un unique antécédent).
  \item \cle{Composée} : $(g\circ f)(x)=g\bigl(f(x)\bigr)$, définie si $x\in D_f$ et $f(x)\in D_g$.
  \item $f$ \cle{paire} : $D_f$ symétrique par rapport à $0$ et $f(-x)=f(x)$ ; \cle{impaire} : $f(-x)=-f(x)$ ; \cle{périodique} de période $T$ : $f(x+T)=f(x)$.
\end{itemize}

\textbf{Formules et résultats à connaître par cœur}
\begin{center}
\begin{tabularx}{\linewidth}{|l|X|}\hline
\rowcolor{popblueL} \textbf{Notion} & \textbf{Résultat} \\ \hline
Bijection réciproque & Résoudre $y=f(x)$ d'inconnue $x$ : $x=f^{-1}(y)$. \\ \hline
Somme, produit & $D_{f+g}=D_{fg}=D_f\cap D_g$. \\ \hline
Quotient & $D_{f/g}=\{x\in D_f\cap D_g\mid g(x)\neq 0\}$. \\ \hline
Composée & $D_{g\circ f}=\{x\in D_f\mid f(x)\in D_g\}$. \\ \hline
Centre de symétrie $\Omega(a;b)$ & $f(a+x)+f(a-x)=2b$ pour tout $x$ tel que $a\pm x\in D_f$. \\ \hline
Axe de symétrie $x=a$ & $f(a+x)=f(a-x)$ pour tout $x$ tel que $a\pm x\in D_f$. \\ \hline
Appartenance & $M(x_0;y_0)\in\mathcal{C}_f \iff x_0\in D_f$ et $f(x_0)=y_0$. \\ \hline
\end{tabularx}
\end{center}

\textbf{Méthodes express}
\begin{enumerate}
  \item Ensemble de définition : repérer dénominateurs et racines, poser les conditions, résoudre, conclure en union d'intervalles.
  \item Injectivité : partir de $f(a)=f(b)$ et prouver $a=b$ (ou exhiber un contre-exemple).
  \item Surjectivité : pour $y$ quelconque, résoudre $f(x)=y$ et vérifier que $x\in E$.
  \item Parité : vérifier d'abord que $D_f$ est centré en $0$, puis comparer $f(-x)$ à $f(x)$.
  \item Fonctions associées : $x\mapsto f(x)+k$ (translation verticale), $x\mapsto f(x-a)$ (translation horizontale), $x\mapsto -f(x)$ et $x\mapsto f(-x)$ (symétries).
\end{enumerate}
\end{bilan}

\begin{aretenir}[Checklist avant le Probatoire]
\begin{itemize}
  \item[$\square$] Je sais déterminer un ensemble de définition (dénominateur non nul, radicande positif).
  \item[$\square$] Je sais écrire la restriction d'une fonction à un intervalle.
  \item[$\square$] Je sais prouver l'injectivité en partant de $f(a)=f(b)$.
  \item[$\square$] Je sais prouver la surjectivité en résolvant $f(x)=y$.
  \item[$\square$] Je sais conclure « bijective » et expliciter la bijection réciproque.
  \item[$\square$] Je sais calculer $(g\circ f)(x)$ sans inverser l'ordre des fonctions.
  \item[$\square$] Je sais donner l'ensemble de définition d'une somme, d'un produit, d'un quotient, d'une composée.
  \item[$\square$] Je sais additionner et multiplier deux polynômes, et donner la condition d'existence d'un quotient.
  \item[$\square$] Je sais montrer qu'une fonction est paire, impaire ou périodique.
  \item[$\square$] Je sais justifier un centre de symétrie $\Omega(a;b)$ et un axe de symétrie $x=a$.
  \item[$\square$] Je sais vérifier qu'un point appartient à une courbe.
  \item[$\square$] Je sais lire graphiquement ensemble de définition, variations, symétries et asymptotes éventuelles.
\end{itemize}
\end{aretenir}

\findecours

\end{document}
